% DER_VII-b: Grothendieck's text with the commentary set in, then the overview, inside a group that applies the chunk's own definitions
\dgrpart{1}{VII-b}{Part 2 of 3, manuscript pages 58--117}{VII-b\enspace\textperiodcentered\enspace Chapter VII, part 2 of 3\enspace\textperiodcentered\enspace manuscript pages 58--117}{Chapter VII \textperiodcentered\ \textup{VII-b}}
\dgrchunkbegin
% DER_VII-b: the definitions of the chunk's preamble, applied inside its group
\providecommand\IfDocumentMetadataT[1]{}
\let\printedpage\relax\newcommand{\printedpage}[1]{\marginpar{\footnotesize p.~#1}}
\let\msp\relax\newcommand{\msp}[1]{\marginpar{\vspace*{10pt}\footnotesize [#1]}}
\let\uTy\relax\newcommand{\uTy}{\underline{\mathrm{Ty}}}
\let\Ch\relax\newcommand{\Ch}{\underline{\mathrm{Ch}}}
\let\Ty\relax\newcommand{\Ty}{\underline{\mathrm{Ty}}}
\let\Nerf\relax\newcommand{\Nerf}{\mathrm{Nerf}}
\let\Fl\relax\newcommand{\Fl}{\mathrm{Fl}}
\let\Hom\relax\newcommand{\Hom}{\mathrm{Hom}}
\let\ch\relax\newcommand{\ch}{\mathrm{ch}}
\let\hot\relax\newcommand{\hot}{\mathrm{hot}}
\let\id\relax\newcommand{\id}{\mathrm{id}}
\let\CH\relax\newcommand{\CH}{\underline{\mathrm{CH}}}
\let\red\relax\newcommand{\red}{\mathrm{r\acute{e}d}}
\let\uHom\relax\newcommand{\uHom}{\underline{\mathrm{Hom}}}
\let\ul\relax\newcommand{\ul}[1]{\underline{#1}}
\let\triv\relax\newcommand{\triv}{\mathrm{triv}}
\let\uCh\relax\newcommand{\uCh}{\underline{\mathrm{Ch}}}
\let\Cat\relax\newcommand{\Cat}{\mathrm{Cat}}
\let\Cathot\relax\newcommand{\Cathot}{\mathrm{Cat\,hot}}
\let\Hot\relax\newcommand{\Hot}{\mathrm{Hot}}
\let\Ob\relax\newcommand{\Ob}{\mathrm{Ob}\,}
\let\sHom\relax\newcommand{\sHom}{\underline{\mathcal{H}\mathit{om}}}
\let\To\relax\newcommand{\To}{T_{\mathrm{o}}}
\let\circled\relax\newcommand{\circled}[1]{\tikz[baseline=(c.base)]\node[draw,circle,inner sep=1pt](c){#1};}
\let\Hotz\relax\newcommand{\Hotz}{\mathrm{Hot}^{\wedge}_{0}}
\let\CNo\relax\newcommand{\CNo}{C_{\mathbf{N}^{o}}}
\let\lab\relax\newcommand{\lab}[2][l]{\text{\footnotesize\begin{tabular}{@{}#1@{}}#2\end{tabular}}}

% ---- from Der-ChVII-456_p041-045.tex
\printedpage{39}
% head: A. Grothendieck : Dérivateurs, Ch. VII, Édition des Universités de Montpellier II et Paris VII
D'autre part, posant $l_i = \mathrm{long}(\pi_{i-1})$ (pour $i \in [1,r]$),
\[
f_i : \Delta_\varepsilon^{l_{i-1}} \longrightarrow \Delta_\varepsilon^{l_i\ (=\mathrm{long}(\pi_{i-1}))}
\]
est défini comme
\[
f_i = \mathfrak{f}_{\varepsilon,\pi_{i-1}(1)} * \mathfrak{f}_{\varepsilon,\pi_{i-1}(2)} * \cdots * \mathfrak{f}_{\varepsilon,\pi_{i-1}(l_{i-1})},
\]
de sorte que l'élément $(*)$ de $\Nerf_r(\Ty)$ est déterminé de façon unique par la donnée des $\varepsilon \in \{\pm 1\}$ et de
\[
\pi_\bullet = (\pi_i)_{i\in[0,r-1]} \in (\mathrm{Hom}_{\mathrm{Ens}}(\mathbf{N}^*,\mathbf{N}))^r
\]
satisfaisant les conditions $(**)$. Soit $F_{\varepsilon,\pi_\bullet}$ cet élément de $\Nerf_r(\Ty)$, et soit $N_r$ l'ensemble des $\pi_\bullet$ qu'on vient d'expliciter. Alors
\[
(F_{\varepsilon,\pi_\bullet})_{(\varepsilon,\pi_\bullet)\in\{\pm1\}\times N_r}
\]
est un système de générateurs libre du monoïde $\Nerf_r(\Ty)$.

Si $l_i = 0$ (i.e. \fbox{$\pi_{i-1} = 0$}), on doit avoir $\mathrm{supp}(\pi_{i+1}) = \emptyset$, i.e. $\pi_{i+1} = 0$, d'où $l_{i+1} = 0$, et de même pour les $l_j$ et $\pi_j$ suivants. Géométriquement l'hypothèse s'exprime par la condition que le but
\msp{page 58}
de $f_i$ est l'élément 1 de Ty (correspondant à l'intervalle ponctuel), et comme $\mathrm{Hom}(1,\tau) = \emptyset$ pour $\tau \neq 1$, cela signifie que les flèches qui suivent $f_i$ doivent être toutes égales à la flèche $\mathrm{id}_1 : 1 \longrightarrow 1$.

\dgrblock{dgrb315}{Degreewise free is a precise assertion}
\begin{dcomment}{exposition}{p57}{Degreewise free is a precise assertion}
Grothendieck has now made the combinatorics of a changing path shape completely explicit. The commentary at manuscript pages~53--57 already explains the free generators in every degree and counts refinement stages. One additional distinction helps us use that calculation correctly: freeness in each degree does not mean that the intrinsic generators themselves form a simplicial set.

To see the distinction, take the generator of degree two represented by
\[
 \Delta_+\longrightarrow\Delta_+^2
       \xrightarrow{\mathrm{id}}\Delta_+^2.
\]
Its first face, which forgets the initial object and first arrow, is
$\mathrm{id}_{\Delta_+^2}=f_{+,1}*f_{+,1}$. This is a product of two nonunit elements in the degree-one monoid. It is therefore not a generator. In the notation of the text, the general formula is
\[
 d_0(f_{\varepsilon,l,\pi})
    =f_{\varepsilon,\pi(1)}*\cdots*f_{\varepsilon,\pi(l)},
\]
where the empty product for $l=0$ is the unit. An isomorphism between free monoids preserves their intrinsic indecomposables. Consequently $N_\bullet\underline{\mathrm{Ty}}$ cannot be the degreewise free monoid on a simplicial set: such a construction would send each generator under every face to a generator. The monoid structure maps are determined by generators, as Grothendieck says; their images can be arbitrary words.

There is also a useful boundary case in his finite-support description. The condition on $\pi_0$ is separate, namely $\operatorname{supp}\pi_0\subseteq\{1\}$. The recursive condition must then run over $1\leq i\leq r-1$, since $\pi_{-1}$ has not been defined. If the length at stage $i$ is zero, the next weighting is $\pi_i=0$, and every later stage remains the empty interval. The edition instead writes $\pi_{i+1}$ at this step. Both index slips are retained in the source and registered in the corrections table. They change neither the free decomposition nor its geometric meaning: a path shape which has disappeared cannot reappear under an endpoint-preserving interval map.
\end{dcomment}


\textbf{Remarque.} Dans la description combinatoire de $\Ty$ et de $\Nerf_\bullet(\Ty)$, l'ensemble des `signatures' $\{\pm1\}$ joue un rôle assez bidon. Il peut être remplacé par n'importe quel ensemble $\Sigma$ (dit `ensemble des signatures'), et on construit alors une catégorie monoïdale (ou un monoïde catégorique) ayant comme ensemble d'objets le monoïde libre engendré par $\Sigma$, interprétés comme des intervalles d'entiers $[0,n]$, munis de `signatures' dans $\Sigma$ pour chaque intervalle de transition, ou ce qui revient au même, des `graphes' orientés $\Sigma$-pondérés de type combinatoire un intervalle $\boldsymbol{\Delta}_i$, ce qui revient à la donnée d'une suite (ordonnée) de $n$ éléments $\varepsilon_1,\ \varepsilon_2,\ \ldots\ ,\ \varepsilon_n \in \Sigma$, i.e. d'un `mot' $\varepsilon_1\varepsilon_2\ldots\varepsilon_n$, qu'on préfère écrire $\Delta_{\varepsilon_1}\Delta_{\varepsilon_2}\cdots\Delta_{\varepsilon_n}$. On peut définir de façon évidente les homomorphismes \textit{croissants} de tels graphes pondérés orientés (sous-entendu : avec l'élément source et l'élément but dans la structure, à respecter par les morphismes), sans avoir besoin d'une interprétation catégorique. La fonctorialité de l'opération de composition des intervalles pondérés, pour cette notion de morphisme, est évidente. Le
\msp{page 59}
choix $\Sigma = \{\pm1\}$ est à ma connaissance le seul utile, à cause de son lien avec la notion de chemins et le calcul de $\pi_0$ d'un objet de $X$ \texttt{[plutôt d'un objet} $X$ \texttt{de} Cat\texttt{]}, etc.
% foot: Version 22 décembre 2001

\dgrblock{dgrb316}{The localization test}
\begin{dcomment}{exposition}{p59:le seul utile, à cause de son lien}{The localization test}
Grothendieck is about to test his category of types against a construction whose universal property is already settled: the category obtained by formally inverting specified arrows. This is a demanding first test because literal concatenation gives too much information. It remembers the length of a zigzag, every inserted identity, and every subdivision of a composite. A localization must identify presentations which differ in these ways, while retaining their common source and target.

Let $P(X)$ denote the category whose arrows are literal finite zigzags in $X$. Length is additive under composition. Thus an invertible arrow of $P(X)$ must have length zero. Moreover the assignment $u\mapsto c(u)$ from arrows of $X$ to one-edge paths does not preserve composition: $c(vu)$ has one edge and $c(v)c(u)$ has two. It does not preserve identities either, since $c(\mathrm{id}_a)$ has length one. These observations concern the literal path construction uniformly in $X$. For particular categories there can of course be useful functors in one direction: if $X$ is a groupoid, every zigzag has an effect in $X$, and evaluation is a functor $P(X)\to X$.

The category $\underline{\mathrm{Ch}}(a,b)$ supplies the missing identifications as morphisms. A type map replaces a string of composable arrows by its composite or inserts identity arrows through pullback. A natural transformation compares two paths of the same type by a ladder of commuting squares. Since the endpoints are fixed, the components at $a$ and $b$ must be identities. Passing to $\pi_0$ identifies objects joined by a finite zigzag of these morphisms. A single directed natural transformation is enough for this purpose; it need not be invertible.

The desired category therefore has
\[
 \operatorname{Hom}(a,b)=\pi_0\underline{\mathrm{Ch}}(a,b).
\]
Its composition comes from concatenation before taking components. This is legitimate because $\pi_0(A\times B)=\pi_0(A)\times\pi_0(B)$ for categories: a zigzag in a product projects to the factors, while zigzags in the factors can be traversed one coordinate at a time. Associativity and units are then inherited, rather than imposed as further equivalence relations.

The relative version allows only a chosen class $\Sigma$ of arrows to be traversed backwards. Its aim is the formula
\[
 \pi_0\underline{\mathrm{Ch}}_\Sigma(a,b)
       \cong \operatorname{Hom}_{X[\Sigma^{-1}]}(a,b).
\]
At this stage the assertion concerns ordinary hom-sets. The homotopy type of the whole path category is a separate, stronger question, and the later pages of the chapter turn precisely to the additional choices and coherence needed for it.
\end{dcomment}


\printedpage{40}
% head: A. Grothendieck : Dérivateurs, Ch. VII, Édition des Universités de Montpellier II et Paris VII
\section*{7\quad Explicitation des catégories de fractions $X\Sigma^{-1}$ en termes des $\Ch(X;a,b)$}

Si $X$ est une catégorie, on désigne par $\Ch(X)$ la catégorie des chemins dans $X$ (cf. §4), et pour $a,b \in \mathrm{Ob}\,X$, par
\[
\Ch(X;a,b) \qquad \text{ou simplement } \Ch(a,b)
\]
celle des chemins de $a$ vers $b$, i.e. de source $a$, but $b$. L'intuition géométrique (qu'on explicitera), c'est qu'un tel chemin $c$ définit un chemin `direct' de $a$ vers $b$ (i.e. un chemin de type $\tau'$ `direct', i.e. toutes les flèches de transition sont directes, i.e. $\tau'$ de la forme $\Delta_+^n$), de même longueur que le chemin initial $c$, obtenu en inversant toutes les flèches rétrogrades dans le chemin $c$. Cette opération est effectivement possible dans $X$ lui-même si et seulement si lesdites flèches sont inversibles. On obtient alors une suite composable de flèches, dont le composé total (qu'on peut appeler `l'effet utile' du chemin)
\msp{page 60}
est une flèche de $a$ vers $b$. Dans le cas général, les chemins de $a$ vers $b$ peuvent être vus comme une sorte de `flèches généralisées' allant de $a$ vers $b$, dont la composition serait la composition évidente, et les identités les `chemins vides' relatifs aux $a \in \mathrm{Ob}\,X$. On trouve d'ailleurs bien ainsi une catégorie, ayant mêmes objets que $X$, mais où les
\[
\mathrm{Hom}_{\mathrm{Ch}}(a,b) = \mathrm{Ch}(a,b) = \mathrm{Ob}\,\Ch(a,b)
\]
sont prohibitivement gros. Cela s'exprime notamment par le fait que les seules flèches inversibles dans $X_{\mathrm{Ch}}$ sont les identités ! On n'a pas de foncteur naturel entre $X_{\mathrm{Ch}}$ et $X$ qui soit l'identité sur les objets. Car à un chemin $a \xrightarrow{c} b$ on ne peut associer de $X$-morphismes $a \longrightarrow b$ (sauf si $c$ \texttt{[est]} de type $\Delta_+$), et à un $X$-morphisme $a \xrightarrow{u} b$ on peut associer un chemin $c(u)$ de type $\Delta_+$, mais si $a \xrightarrow{u} b \xrightarrow{v} c$, $c(vu)$ est encore de type $\Delta_+$ alors que $c(v)c(u)$ est de type $\Delta_+^2$, ils ne sont pas égaux. Par contre, il y a une flèche naturelle entre ces deux éléments de $\Ch(a,c)$, flèche au dessus de l'unique homomorphisme
\[
\Delta_+ \xrightarrow{\mathfrak{f}_{+,2}} \Delta_+^2.
\]
\msp{page 61}
On a alors
\[
\mathfrak{f}_{+,2}^*(\underbrace{a \xrightarrow{u} b \xrightarrow{v} c}_{c(v)c(u)}) = (\underbrace{a \xrightarrow{vu} c}_{c(vu)}),
\]
d'où une flèche canonique
\[
c(vu) \longrightarrow c(v)c(u) \qquad \text{dans } \Ch(a,c).
\]
Cela montre que $c(vu)$ et $c(v)c(u)$ sont tout au moins dans une même composante connexe de $\Ch(a,c)$, ou (comme on dit encore) qu'ils sont `homotopes'. Cela suggère ici de modifier la définition de $\mathrm{Hom}_{\mathrm{Ch}}(a,b)$ en
\[
\pi(a,b) = \pi_0(\Ch(a,b)).
\]
% foot: Version 22 décembre 2001

\printedpage{41}
% head: A. Grothendieck : Dérivateurs, Ch. VII, Édition des Universités de Montpellier II et Paris VII
Le foncteur de composition des chemins
\[
\begin{array}{ccccc}
\Ch(a,b) & \times & \Ch(b,c) & \longrightarrow & \Ch(a,c)\\
(c & , & c') & \longmapsto & c' \circ c
\end{array}
\]
\texttt{[il y a trop de} $c$\texttt{]} (\footnote[17]{Ce foncteur est déduit des foncteurs
\[
\Ch_\tau(a,b) \times \Ch_{\tau'}(b,c) \longrightarrow \Ch_{\tau\tau'}(a,c)
\]
évidents, cf. page 77.}) induit une application sur les $\pi_0$, ce qui donne une loi de composition
\[
\pi(a,b) \times \pi(b,c) \longrightarrow \pi(a,c)
\]
sur les flèches. Il est immédiat qu'elle est associative (la composition des chemins l'étant), et unitaire. Elle admet comme flèche identique en $x$ soit le `lacet vide en $x$', lequel est (directement) homotope dans $\Ch(x,x)$ au chemin `dégénéré' $x \xrightarrow{\mathrm{id}} x$
\msp{page 62}
de type $\Delta_+$, ou aussi au chemin dégénéré de n'importe quel type $\tau$, disons
\[
x \xleftarrow{\mathrm{id}} x \xrightarrow{\mathrm{id}} x \xrightarrow{\mathrm{id}} x \xleftarrow{\mathrm{id}} x \xleftarrow{\mathrm{id}} x,
\]
un tel chemin n'étant autre que $p_\tau^*(x)$, où
\[
p_\tau : \tau \longrightarrow 1
\]
est le morphisme de dégénérescence (l'unique morphisme dans $\Ty$ vers 1, qui est objet final), d'où
\[
p_\tau^* : \Ch_1(x,x) = \{x\} \longrightarrow \Ch_\tau(x,x) \subseteq \Ch(x,x).
\]
La catégorie $\Pi(X)$ ainsi obtenu doit s'appeler le \textit{groupoïde fondamental} de $X$. En effet, % [sic] obtenu -> obtenue, agreement with « La catégorie »

\dgrblock{dgrb317}{A span and its fraction}
\begin{dcomment}{examples}{p61}{A span and its fraction}
The construction becomes concrete as soon as reversing one arrow creates a morphism which the original category did not contain. Consider the category $X$ generated by the span
\[
 a\xleftarrow{s}x\xrightarrow{t}b
\]
with no nonidentity arrows besides $s$ and $t$, and put
$\Sigma=\{\mathrm{id}_a,\mathrm{id}_x,\mathrm{id}_b,s\}$. The zigzag from $a$ to $b$ is a $\Sigma$-path although $\operatorname{Hom}_X(a,b)$ is empty. Its effect after localization is the fraction $ts^{-1}$.

We can compute the entire localization. It has the same three objects and the following hom-sets, with rows denoting sources:
\[
\begin{array}{c|ccc}
 &a&x&b\\ \hline
 a&\{\mathrm{id}_a\}&\{s^{-1}\}&\{ts^{-1}\}\\
 x&\{s\}&\{\mathrm{id}_x\}&\{t\}\\
 b&\varnothing&\varnothing&\{\mathrm{id}_b\}.
\end{array}
\]
Composition is determined by $s^{-1}s=\mathrm{id}_x$,
$ss^{-1}=\mathrm{id}_a$, and $(ts^{-1})s=t$. For any functor $F:X\to Y$ with $F(s)$ invertible, there is exactly one extension to this category: it sends $s^{-1}$ to $F(s)^{-1}$ and $ts^{-1}$ to $F(t)F(s)^{-1}$. This verifies the universal property directly.

The same calculation can be read inside the path categories. A path starting at $b$ can only stay at $b$, because $t$ is not permitted as a backward arrow. A nonconstant path ending at $b$ reaches it by a final occurrence of $t$; before that it moves between $a$ and $x$ using $s$ in either direction. Adjacent traversals of $s$ in opposite directions cancel by the homotopies on manuscript page 63. After removing identity edges, only one of the entries in the displayed table remains. This proves directly that the connected components of the relative path categories give exactly those hom-sets.

The two objects $a$ and $x$ are isomorphic in the result. Choosing $a$ as their representative identifies a skeleton with the walking arrow $a\to b$. Thus localization has changed the hom-set from $a$ to $b$, while it has introduced no path from $b$ to $a$. This is the elementary reason that a relative localization is generally a category rather than a groupoid.
\end{dcomment}


\begin{enumerate}[label=\alph*)]
\item toutes les flèches de $\Pi(X)$ sont inversibles, i.e. $\Pi(X)$ est bien un groupoïde, et
\item le foncteur canonique
\[
X \longrightarrow \Pi(X)
\]
qui est l'identité sur les objets, et qui est défini sur les flèches par $u \longmapsto c(u)$ comme ci-dessus, est \textit{universel} (dans Cat) pour les foncteurs (i.e. les Cat-morphismes) de $X$ vers des \textit{groupoïdes.} (Plus généralement encore, il est universel pour les foncteurs $X \xrightarrow{F} Y$ tels que $F$ transforme toute flèche de $X$ en \texttt{[une]} flèche inversible.)
\end{enumerate}
% foot: Version 22 décembre 2001

\printedpage{42}
% head: A. Grothendieck : Dérivateurs, Ch. VII, Édition des Universités de Montpellier II et Paris VII
\msp{page 63}
Pour prouver que toute flèche dans $\Pi(X)$ est inversible, comme tout chemin non vide est composé de chemins `élémentaires' de type $\Delta_+$ ou $\Delta_-$, il suffit de prouver que ces chemins sont inversibles. En fait, ils se correspondent biunivoquement,
\[
\Ch_{\Delta_+}(a,b) \simeq \Ch_{\Delta_-}(b,a),
\]
en associant au chemin de $a$ vers $b$
\[
c(u) : a \xrightarrow{u} b \qquad \text{de type } \Delta_+
\]
le chemin
\[
c(u)^o : b \xleftarrow{u} a \qquad \text{de type } \Delta_-
\]
de $b$ vers $a$. Je dis que
\[
\begin{array}{rcl}
c(u)^o c(u) & \sim & 1_a\\
c(u) c(u)^o & \sim & 1_b
\end{array}
\]
\[
\left.
\begin{tikzcd}
 & a \arrow[r, "u"] & b & a \arrow[l, "u"'] \\
 {} & a \arrow[u, "\mathrm{id}"'] \arrow[r, "\mathrm{id}"] & a \arrow[u, "u"'] & a \arrow[u, "\mathrm{id}"'] \arrow[l, "\mathrm{id}"'] \\
 {} \arrow[u] & & a &
\end{tikzcd}
\right\}\
\begin{array}{l}
\text{on a un homomorphisme dans } \Ch_{\Delta_+\Delta_-}(a,a)\\
\text{de } 1_a \text{ dans } c(u)^o c(u).
\end{array}
\]
\[
\left.
\begin{tikzcd}
 & b & a \arrow[l, "u"'] \arrow[r, "u"] & b \\
 {} & a \arrow[u, "u"'] & a \arrow[u, "\mathrm{id}"'] \arrow[l, "\mathrm{id}"'] \arrow[r, "\mathrm{id}"] & a \arrow[u, "u"'] \\
 {} \arrow[u] & & a &
\end{tikzcd}
\right\}\
\begin{array}{l}
\text{itou dans } \Ch_{\Delta_-\Delta_+}(b,b),\\
\text{de } 1_b \longrightarrow c(u)c(u)^0 % [sic] c(u)^0 -> c(u)^o, digit 0 printed instead of the letter o used elsewhere
\end{array}
\]
\texttt{[plutôt} $1_b \longleftarrow c(u)c(u)^0$\texttt{, viz.} % [sic] c(u)^0 -> c(u)^o, digit 0 printed instead of the letter o used elsewhere
\[
\begin{tikzcd}
b \arrow[d, "\mathrm{id}"] & a \arrow[l, "u"'] \arrow[r, "u"] \arrow[d, "u"] & b \arrow[d, "\mathrm{id}"] \\
b & b \arrow[l, "\mathrm{id}"'] \arrow[r, "\mathrm{id}"] & b \ \texttt{].}
\end{tikzcd}
\]
% foot: Version 22 décembre 2001

\printedpage{43}
% head: A. Grothendieck : Dérivateurs, Ch. VII, Édition des Universités de Montpellier II et Paris VII
\msp{page 64}
Notons que, plus généralement, pour tout chemin
\[
c : a \longrightarrow b, \qquad \text{de type } \tau,
\]
on définit un \textit{chemin opposé}
\[
c^o : b \longrightarrow a \qquad \text{de type } \tau^o,
\]
où
\[
\tau \longmapsto \tau^o : \mathrm{Ty} \longrightarrow \mathrm{Ty}
\]
est l'antiïnvolution du monoide libre Ty qui échange les deux générateurs $\Delta_-$, $\Delta_+$, et qui est donc donné sur l'élément général par
\[
(\Delta_{\varepsilon_1}\Delta_{\varepsilon_2}\cdots\Delta_{\varepsilon_n})^o = \Delta_{-\varepsilon_n}\Delta_{-\varepsilon_{n-1}}\cdots\Delta_{-\varepsilon_1}.
\]
Par récurrence sur la longueur de $c$ (\footnote[18]{inutile, résulte directement de la formule $(*)$ \texttt{[de la]} page suivante, en décomposant $c$ en composés de chemins `élémentaires' de type $\Delta_+$ ou $\Delta_-$.}), en utilisant le résultat précédent, que $c(u)$, $c(u)^o$ sont inverses l'un de l'autre dans $\Pi(X)$, on trouve
\[
\left.
\begin{array}{rcl}
c^o c & = & \mathrm{id}_a\\
c c^o & = & \mathrm{id}_b
\end{array}
\right\}\text{dans } \Pi(X).
\]

\dgrblock{dgrb318}{Two cancellation homotopies with fixed endpoints}
\begin{dcomment}{exposition}{p63}{Two cancellation homotopies with fixed endpoints}
Grothendieck wants reversal of a single edge to give an inverse after taking components. The proof requires two different directions of natural transformation, because the two loops have different shapes. Keeping those directions visible also explains the editors' correction on this page.

For $u:a\to b$, the loop at $a$ is the diagram
$L_a=(a\xrightarrow{u}b\xleftarrow{u}a)$. Let $K_a$ be the constant $+-$ diagram at $a$. The components
\[
 (\mathrm{id}_a,u,\mathrm{id}_a):K_a\longrightarrow L_a
\]
form a natural transformation. Both endpoint components are identities, so this is a morphism in the fixed-endpoint fibre
$\underline{\mathrm{Ch}}_{+-}(a,a)$. The collapse of the type $+-$ to the empty type gives a cartesian arrow from $K_a$ to the empty path at $a$. We have therefore joined $L_a$ to the empty path by a zigzag in the total fixed-endpoint category.

For the other loop, $L_b=(b\xleftarrow{u}a\xrightarrow{u}b)$, the appropriate constant diagram is $K_b$ at $b$. This time the transformation goes in the other direction:
\[
 (\mathrm{id}_b,u,\mathrm{id}_b):L_b\longrightarrow K_b.
\]
It is a morphism in $\underline{\mathrm{Ch}}_{-+}(b,b)$, and $K_b$ again maps to the empty path through the type collapse. The uncorrected diagram in the edition has a constant diagram at $a$ below $L_b$. Its endpoint components are $u$, so it is not a morphism in this fixed-endpoint fibre. The bracketed correction supplies exactly the downward transformation just described.

Two conventions are doing work here. First, a constant diagram of nonempty type is an object of its fibre, whereas the empty path belongs to the fibre over the unit type. When the source writes $1_a$ or $1_b$ in a fibre, it is temporarily using the constant diagram as a representative of the unit component. Second, connectedness forgets the direction of the arrows used to connect objects. Nothing requires both cancellation homotopies to point the same way.

The relative argument is identical when $u\in\Sigma$: every nonidentity backward edge in these diagrams is then an allowed denominator, and all remaining backward edges are identities. This proves that $\kappa(u)$ is invertible in $\Pi_\Sigma(X)$. There is no condition on a vertical component of a path transformation beyond naturality and the identity endpoint condition.
\end{dcomment}


\textbf{Exemple.} Si
\[
\begin{array}{rcl}
\tau & = & \Delta_+\Delta_-\Delta_+^2\Delta_-\\
\tau^o & = & \Delta_+\Delta_-^2\Delta_+\Delta_-,
\end{array}
\]
et si on a un chemin
\[
c : a_0 \xrightarrow{u_1} a_1 \xleftarrow{u_2} a_2 \xrightarrow{u_3} a_3 \xrightarrow{u_4} a_4 \xleftarrow{u_5} a_5
\]
de type $\tau$, on en déduit le chemin
\[
c^o : a_5 \xrightarrow{u_5} a_4 \xleftarrow{u_4} a_3 \xleftarrow{u_3} a_2 \xrightarrow{u_2} a_1 \xleftarrow{u_1} a_0
\]
de type $\tau^o$. L'opération $c \longmapsto c^o$ satisfait en
\msp{page 65}
effet \texttt{[pour]}
\[
\begin{array}{ccccc}
a & \xrightarrow{\gamma} & b & \xrightarrow{\gamma'} & c\\
c & \xrightarrow{\gamma'^o} & b & \xrightarrow{\gamma^o} & a
\end{array}
\]
la relation
\[
(\gamma'\gamma)^o = \gamma^o\gamma'^o, \tag*{$(*)$}
\]
% foot: Version 22 décembre 2001

\dgrblock{dgrb319}{A circle in a category with two arrows}
\begin{dcomment}{examples}{p64}{A circle in a category with two arrows}
Grothendieck's reversal operation removes backtracking without destroying genuine loops. A small example distinguishes the two phenomena. Let $X$ have two objects $a,b$ and two distinct nonidentity arrows $f,g:a\to b$, and invert every arrow. The zigzag
\[
 a\xrightarrow{f}b\xleftarrow{g}a
\]
represents the loop $h=g^{-1}f$ at $a$.

All morphisms in the resulting groupoid have unique expressions of the following forms, with $n\in\mathbf Z$:
\[
\begin{array}{c|cc}
 &a&b\\ \hline
 a&h^n&gh^n\\
 b&h^ng^{-1}&gh^ng^{-1}.
\end{array}
\]
For existence, use $g$ as a chosen path from $a$ to $b$ and replace $f$ by $gh$. After cancelling adjacent $g^{-1}g$ and $gg^{-1}$, a word has exactly the indicated form. For uniqueness, consider the functor from $X$ to the one-object group $\mathbf Z$ which sends $f$ to $1$ and $g$ to $0$. Since its target is a groupoid, it extends to the groupoid completion, and it sends $h^n$ to $n$. Thus different integers cannot give equal arrows. Composition in the table adds the integer exponents, with the displayed $g$ and $g^{-1}$ providing the changes of endpoint.

It follows that $\operatorname{Aut}_{\Pi(X)}(a)\cong\mathbf Z$ and
$\pi_0\underline{\mathrm{Ch}}(a,a)\cong\mathbf Z$. The nerve supplies a geometric check: the only nondegenerate positive-dimensional simplices of $NX$ are the two edges $f,g$. There are no composable pairs of nonidentity arrows. Its realization is therefore the graph consisting of two vertices joined by two edges, which is a circle.

If instead $\Sigma$ consists only of identities, every backward edge is an identity and the effect of a path is already defined in $X$. The relative localization is then $X$ itself. In particular the two arrows $f,g$ remain distinct and there is no nonidentity loop at $a$. The difference lies entirely in which arrows may be traversed backwards, not in a choice of representative for the same unrestricted path space.
\end{dcomment}


% ---- from Der-ChVII-456_p046-050.tex
\printedpage{44}
% head: A. Grothendieck : Dérivateurs, Ch. VII, Édition des Universités de Montpellier II et Paris VII
% foot: Version 22 décembre 2001

ce qui est en accord avec l'intuition que l'opération $\gamma \longmapsto \gamma^o$ joue le rôle d'un passage à l'inverse (`à homotopie près'), et on a la formule
\[
(\gamma'\gamma)^{-1} = \gamma^{-1}\gamma'^{-1}
\]
pour deux flèches inversibles $\gamma$, $\gamma'$ composables (dans l'ordre $\gamma'$, $\gamma$).

Pour la propriété universelle de
\[
X \longrightarrow \Pi(X)
\]
(qui permettra d'écrire
\[
\Pi(X) \simeq X\Sigma^{-1}, \qquad \text{avec } \Sigma = \Fl(X) \ \ ),
\]
autant l'établir tout de suite dans le cadre général où on se donne une partie
\[
\Sigma \subseteq \Fl(X),
\]
et on définit, pour $\tau \in \Ty$,
\[
\Ch_{\Sigma,\tau}(X) \subseteq \Ch_{\tau}(X)
\]
comme la sous-catégorie pleine de $\Ch_{\tau}(X)$ formée des chemins dans $X$ dont \textit{toutes les flèches rétrogrades}

\msp{page 66}

sont dans $X$ \texttt{[plutôt dans $\Sigma$]}. On notera que pour qu'un chemin $c$ et son opposé $c^o$ soient tous les deux des $\Sigma$-\textit{chemins,} il faut et il suffit que \textit{toutes} les flèches du chemin $c$ soient dans $\Sigma$ (on dira alors que $c$ est un $\Sigma$-\textit{chemin stricte}), puisque les flèches rétrogrades de $c^o$ sont les flèches directes de $c$.

Pour que les sous-catégories pleines $\Ch_{\tau,\Sigma}(X)$ soient stables par les foncteurs images inverses
\[
f^* : \Ch_{\tau'}(X) \longrightarrow \Ch_{\tau}(X)
\]
associés aux flèches
\[
f : \tau \longrightarrow \tau'
\]
dans $\uTy$, ou simplement pour qu'y ait stabilité pour les deux foncteurs $f^*$ associés à un $f : \Delta_- \longrightarrow 1$ et $\Delta_- \longrightarrow \Delta_-^2$, il faut et il suffit que $\Sigma$ soit stable par composition et contienne toutes les flèches identiques % [sic] qu'y ait -> qu'il y ait, missing word
\[
\left\{
\begin{aligned}
& x \in \mathrm{Ob}\,X \ \Longrightarrow \ \id_x \in \Sigma \\
& u, v \in \Sigma, \ \ u,\, v \text{ composables} \Longrightarrow uv \in \Sigma.
\end{aligned}
\right.
\]
En effet, les flèches dans un chemin $f^*(c') \in \Ch_{\tau}(X)$ sont soit des flèches identiques (qui correspondent aux flèches de transition de $I_\tau$ dégénérées pour $f$), soit des flèches parmi celles de $\tau$, ou des composés d'une suite de flèches consécutives composables de $\tau'$ (qui correspondent aux flèches de transition dans $I_{\tau'}$ incluses dans la partie d'une flèche de transition de $I_\tau$), avec conservation des signatures.

\printedpage{45}
% head: A. Grothendieck : Dérivateurs, Ch. VII, Édition des Universités de Montpellier II et Paris VII
% foot: Version 22 décembre 2001

\msp{page 67}

\dgrblock{dgrb320}{Why denominators must form a subcategory}
\begin{dcomment}{exposition}{p66}{Why denominators must form a subcategory}
Grothendieck imposes exactly the closure needed to let every interval map act on relative paths. A backward source edge may cover several consecutive backward target edges, or it may collapse to a vertex. Pullback must therefore allow composition of denominators and insertion of identity denominators. These are the two tests singled out in the text.

The map $\Delta_-\to1$ pulls the empty path at $x$ back to the one-edge path $x\xleftarrow{\mathrm{id}_x}x$. Stability forces $\mathrm{id}_x\in\Sigma$ for every $x$. Next, for composable denominators $v:z\to y$ and $u:y\to x$, the backward path
\[
 x\xleftarrow{u}y\xleftarrow{v}z
\]
pulls back along $\Delta_-\to\Delta_-^2$ to
$x\xleftarrow{uv}z$. Thus stability forces $uv\in\Sigma$.
Conversely, assume these two closure properties. The image of a source edge under an endpoint-preserving increasing type map is either a vertex or a consecutive monochromatic block. A backward source edge consequently becomes either an identity or a composite of backward target arrows. Both belong to $\Sigma$. This proves stability under every type map, including maps with both collapses and subdivisions.

The condition says that $\Sigma$ is the arrow set of a wide subcategory of $X$. It does not require an Ore condition, saturation, or a two-out-of-three property. If we begin with an arbitrary collection $S$ of arrows to invert, we may first take its closure under identities and finite composition. A functor inverts this closure exactly when it inverts $S$, so the intended universal localization is unchanged. Saturation is a later operation: it adds every arrow whose image in the localization happens to be invertible.

The adjective ``full'' in $\underline{\mathrm{Ch}}_{\Sigma,\tau}(X)$ deserves equal attention. The restriction is on backward \emph{horizontal} arrows of the path. Given two such paths, every natural transformation between them remains available, with identity components at prescribed endpoints. Its internal vertical arrows need not lie in $\Sigma$. This fullness is consistent with the ordinary localization theorem, as the calculation at manuscript page 68 will show. It must be remembered when making a higher-homotopical comparison with constructions which restrict vertical arrows as well.

Finally, reversal of an individual path remains allowed precisely when its forward arrows also belong to $\Sigma$. Closure of the entire path system under reversal therefore forces $\Sigma=\operatorname{Fl}(X)$. This statement concerns literal permitted paths; it does not say that a localization can be a groupoid only in that case. Arrows outside $\Sigma$ may already be invertible in $X$, or become invertible after localization.
\end{dcomment}


On suppose ces conditions satisfaites, de sorte que les $\Ch_{\tau,\Sigma}(X)$ forment les fibres d'une sous-catégorie fibrée pleine de $\Ch(X)$ sur $\uTy$, qu'on dénotera par
\[
\Ch_{\Sigma}(X) \hookrightarrow \Ch(X).
\]
Il est évident que cette sous-catégorie est stable par l'opération de composition des chemins (mais non pas, en général, on l'a dit, par passage de $c$ à $c^o$ (il faudrait pour cela que $\Sigma$ contienne \textit{toutes} les flèches de $X$, puisque pour toute $u : a \longrightarrow b$, qui est donc dans $\Ch_{\Delta_+,\Sigma}(X)$, $u^o$ donné par $b \xleftarrow{\ u\ } a$ dans $\Ch_{\Delta_-}(X)$ devrait être un $\Sigma$-chemin, d'où $u \in \Sigma$)).

On va définir une catégorie canonique
\[
\Pi_\Sigma(X),
\]
ayant même objets que $X$, et un foncteur canonique
\[
X \longrightarrow \Pi_\Sigma(X)
\]
qui est l'identité sur les objets, et qui transforme les $u \in \Sigma$ en flèches inversibles. Pour cela on calque la description du groupoïde fondamental (mais $\Pi_\Sigma(X)$ ne sera pas en général un groupoïde), en posant
\[
\Hom_\Sigma(x,y) = \pi_0(\Ch_\Sigma(x,y)),
\]

\msp{page 68}

et en définissant la composition des `$\Sigma$-flèches' en termes de la composition des chemins. L'argument qui prouvait, dans le cas de $\Pi(X)$, que \textit{toute} flèche de $X$, vue comme chemin $c$ de type $\Delta_+$, était inversible dans $\Pi(X)$, ayant comme inverse $c^o$ de type $\Delta_-$, est encore valable, pourvu que $c^o$ soit encore bien un $\Sigma$-chemin, ce qui exige que $u \in \Sigma$.

On va l'interpréter en termes d'un foncteur canonique
\[
\kappa : X \longrightarrow \Pi_\Sigma(X) \qquad \text{(\footnotemark[19])}
\tag*{$(*)$}
\]
\footnotetext[19]{Sur les flèches, $\kappa$ est défini, pour $u : a \longrightarrow b$, par $\kappa(u) =$ classe d'homotopie de $\ch(u)$ dans $\Ch_\Sigma(a,b)$, où $\ch(u)$ est le chemin de type $\Delta_+$ de $a$ vers $b$ défini par $u$.}
qui se décrit encore comme dans le cas du groupoïde fondamental, en associant à toute flèche $u$ de $X$ la classe d'homotopie du chemin $c(u)$ qu'elle définit. On voit donc que ce foncteur canonique transforme les $u \in \Sigma$ en flèches inversibles.

Quand les flèches de $\Sigma$ sont déjà inversibles dans $X$, je dis que le foncteur $(*)$ est un isomorphisme, i.e. que pour tout $a, b \in \mathrm{Ob}\,X$,
\[
\alpha : \Hom_X(a,b) \longrightarrow \pi_0\Ch_\Sigma(a,b), \qquad u \longmapsto \underbrace{\ch_\Sigma(u)}_{= \ch(u)}
\]
est bijectif. Pour le voir, on va définir
\[
\beta : \pi_0(\Ch_\Sigma(a,b)) \longrightarrow \Hom_X(a,b),
\]
qui sera inverse de $\beta$. On définit % [sic] inverse de $\beta$ -> inverse de $\alpha$, $\beta$ is meant to be the inverse of $\alpha$

\printedpage{46}
% head: A. Grothendieck : Dérivateurs, Ch. VII, Édition des Universités de Montpellier II et Paris VII
% foot: Version 22 décembre 2001

\msp{page 69}

\dgrblock{dgrb321}{The effect of a zigzag descends to components}
\begin{dcomment}{exposition}{p68}{The effect of a zigzag descends to components}
Grothendieck now uses evaluation in a category where the denominators are already invertible. Before checking that evaluation is inverse to $\alpha$, we should verify that it is defined on a connected component of the path category. This is a short verification, but it explains why the internal vertical arrows of a path transformation need not be invertible.

The useful statement is slightly more general. Let $F:X\to Y$ send every arrow of $\Sigma$ to an isomorphism. For a $\Sigma$-path
$c=(x_0,\ldots,x_n)$, let
$e_i:F(x_{i-1})\to F(x_i)$ be $F(u_i)$ if its $i$th edge is forward and $F(u_i)^{-1}$ if that edge is backward. Define
\[
 E_F(c)=e_n\cdots e_1,
 \qquad E_F(\text{empty path at }a)=\mathrm{id}_{F(a)}.
\]
This formula is defined solely by the inversion hypothesis.

First compare paths $c,c'$ of the same type by a natural transformation with components $\eta_i:F(x_i)\to F(x'_i)$. For a forward edge the commuting square reads
$\eta_i e_i=e'_i\eta_{i-1}$. For a backward edge the original square is
\[
 \eta_{i-1}F(u_i)=F(u'_i)\eta_i.
\]
Multiplication by the inverses of $F(u_i)$ and $F(u'_i)$ gives exactly the same formula
\[
 \eta_i e_i=e'_i\eta_{i-1}.
\]
Successively composing these equations yields
\[
 \eta_n E_F(c)=E_F(c')\eta_0.
\]
At fixed endpoints, $\eta_0$ and $\eta_n$ are identities. Hence the two effects agree. No inverse of any $\eta_i$ has been taken.

We must also compare different types. If $q:\tau\to\tau'$ and $c'$ is a path of type $\tau'$, then each edge of $q^*c'$ is a composite over the corresponding consecutive block of edges of $c'$, or an identity for an empty block. A forward block has the same effect as its composite. A backward block does too, because inversion reverses the order of composition. The blocks partition the target edges in their traversal order, so
\[
 E_F(q^*c')=E_F(c').
\]
A general morphism in the total path category factors into a morphism within a fibre followed by a cartesian morphism of this form. We have therefore proved invariance under every morphism and hence under every connecting zigzag. Evaluation defines a map on $\pi_0\underline{\mathrm{Ch}}_\Sigma(a,b)$.

Taking $F=\mathrm{id}_X$ gives the missing descent check for $\beta$ in the text. The identity $\beta\alpha=\mathrm{id}$ is immediate. For the other identity, a backward traversal of an invertible $u\in\Sigma$ represents the inverse of $\kappa(u)$; this is also $\kappa(u^{-1})$, since $\kappa$ is a functor. Replacing each backward edge in this way gives the all-forward path $c^\natural$, and repeated composition identifies its class with the elementary path of its composite. These equalities take place in $\Pi_\Sigma(X)$. The omitted subscripts in equations (1) and (2) on manuscript page 69 should not lead us to infer equality in a relative localization merely from equality in its groupoid completion.
\end{dcomment}


$\beta(c)$, pour une $\Sigma$-flèche de type $\tau$, le longueur $n$, en inversant toutes les flèches rétrogrades (ce qui est possibles, puisqu'elles sont dans $\Sigma$), de façon à obtenir un chemin $c^\natural$ de type $\Delta_+^n$, et en prenant le composé (dans $X$) des flèches de ce chemin $c^\natural$. Il est évident alors que $\beta\alpha = \id$, il faut prouver que l'on a aussi $\alpha\beta = \id$, i.e. que dans $\Pi_\Sigma(X)$, les flèches $u = \int c^\natural$ et $c$ sont égales, i.e. $u$ et $c$ sont homotopes dans $\Ch_\Sigma(a,b)$. Considérons le chemin $c^\natural$, de type $\Delta_+^n$ : % [sic] le longueur -> la longueur, gender; possibles -> possible, number
\[
a = a_0 \xrightarrow{u_1} a_1 \xrightarrow{u_2} a_2 \longrightarrow \ \cdots \ \longrightarrow a_{n-1} \xrightarrow{u_n} a_n = b,
\]
et soit $u_i$ ($i \in [1,n]$) une flèche de ce chemin, correspondant à une flèche de transition $\underline{u}_i$ \textit{rétrograde} de $\tau$, donc par construction
\[
u_i = v_i^{-1} \qquad \text{dans } X\text{, avec } v_i \in \Sigma
\]
la flèche d'indice $i$ dans le chemin initial $c$ de type $\tau$. Si $\underline{u}_i$ par contre est directe, on aura
\[
u_i = v_i.
\]
(NB On désigne par $v_i$ la $i$-flèche dans le chemin $c$.) Or on a
\[
c = c_n * c_{n-1} * \cdots * c_1 \qquad \text{dans } \Pi_\Sigma(X)
\tag*{$(*)$}
\]
et le facteur $c_i$ est soit la classe d'homotopie de $\ch(v_i) = \ch(u_i)$ (cas où la flèche de transition $\underline{u}_i$ de $\tau$ est directe), soit celle du chemin \textit{opposé} au chemin $\ch(v_i) : a_i \xrightarrow{v_i} a_{i-1}$ de type $\Delta_+$, donc le chemin $\ch(v_i)^o : a_{i-1} \xleftarrow{v_i} a_i$ de type $\Delta_-$. Je dis que
\[
\boxed{\ \ch(v_i)^o = \ch(u_i) \qquad \text{\textit{dans} } \Pi(X),\ }
\tag*{(1)}
\]
ce qui prouve que dans $\Pi_\Sigma(X)$, $c$ est égal au composé des $\ch(u_i)$ ($i \in [1,n]$), donc à $c^\natural$. Il faut prouver de plus que
\[
c^\natural = \ch(u) \qquad (u = u_nu_{n-1}\cdots u_1) \qquad \text{dans } \Pi(X).
\tag*{(2)}
\]

\msp{page 70}

\dgrblock{dgrb322}{Relations before higher homotopies}
\begin{dcomment}{context}{p69}{Relations before higher homotopies}
Grothendieck's path construction is testing the same universal problem as the classical category of fractions, while retaining the witnesses for its elementary relations. In \citet[Chapter I, \S1.1 and Lemma 1.2]{GabrielZisman1967}, localization is constructed from paths with formal reverse arrows by imposing composition, identity, and inverse relations. Lemma 1.2 also gives the universal property on functor categories and natural transformations. Here composition and identity relations come from type changes; the inverse relations come from the two fixed-endpoint ladders. The evaluation calculation above shows that arbitrary path ladders respect these relations. The universal property on the following pages will prove that taking components introduces no additional identifications.

This is also the appropriate scale on which to read the connection with the derivator programme. In \citet[\S69, pp.~1--4 of the separately edited text]{Grothendieck1983}, the diagramwise construction
\[
 I\longmapsto
 \operatorname{Fun}(I,M)[W_I^{-1}]
\]
retains a whole family of localized diagram categories, where $W_I$ consists of pointwise weak equivalences. The present calculation makes morphisms in such ordinary localizations accessible by diagrams of arrows and commuting squares. It does not by itself supply the homotopy Kan extensions or the higher comparisons required for a derivator.

The chronological comparison with Dwyer and Kan is particularly instructive, since their simplicial localization papers appeared in 1980, before these manuscript pages. Their hammock construction restricts every vertical arrow to the chosen subcategory of weak equivalences, as well as every backward horizontal arrow; its component category is the ordinary localization, and it is weakly equivalent to their standard simplicial localization \citep[Definition 2.1 and Propositions 2.2, 3.1]{DwyerKan1980Calculating}. Grothendieck's full relative path fibres allow arbitrary internal vertical arrows. Agreement of component categories is therefore a common starting point, not an automatic comparison theorem for the entire nerves. The stronger assertions later in the chapter must be justified with the actual path category and its actual indexing maps.
\end{dcomment}


Donc il y a deux vérifications à faire (où je suppose à nouveau $\Sigma$ quelconque) :

\begin{enumerate}[label={}]
\item[$1^o$)] Soit $u \in \Sigma$, $u : a \longrightarrow b$, supposons $\Sigma$ inversible dans $X$. Alors
\[
\ch(u)^o = \ch(u^{-1}) \qquad \text{dans } \Pi_\Sigma(X),
\]
i.e. les deux chemins de $b$ vers
\[
\begin{array}{cccl}
b & \xleftarrow{\ u\ } & a & \text{de type } \Delta_- \\
b & \xrightarrow{u^{-1}} & a & \text{de type } \Delta_+
\end{array}
\]
\printedpage{47}
% head: A. Grothendieck : Dérivateurs, Ch. VII, Édition des Universités de Montpellier II et Paris VII
% foot: Version 22 décembre 2001
sont homotopes dans $\Ch_\Sigma(b,a)$. Or on sait déjà, comme $\kappa : X \longrightarrow \Pi_\Sigma(X)$ est un foncteur, que
\[
\hot\,\ch(u^{-1}) = (\hot\,\ch(u))^{-1} \qquad \text{dans } \Pi_\Sigma(X)
\]
où $\hot(\bullet)$ désigne la classe d'homotopie d'une chaine (ici classe dans $\Ch_\Sigma(b,a)$ et dans $\Ch_\Sigma(a,b)$ respectivement). Donc l'énoncé à prouver c'est que $\hot(\ch(u)^o)$ est égal à l'\textit{inverse} de $\hot(\ch(u))$ dans $\Pi_\Sigma(X)$. C'est donc ce qu'on a prouvé pour \textit{tout} $u \in \Fl(X)$, dans le cas de $\Pi(X)$. L'argument doit marcher, du moment qu'il a un sens, i.e. que $\ch(u)^o$ est bel et bien un $\Sigma$-chemin, i.e. $u \in \Sigma$.

\item[$2^o$)] La formule (2) signifie que pour tout chemin $c$ de type $\Delta_+^n$ (donc tautologiquement un \msp{page 71} $\Sigma$-chemin), allant de $a$ vers $b$, sa classe dans $\Hom_\Sigma(a,b)$ n'est autre que celle de $\ch(u)$, où $u$ est le composé $u_nu_{n-1}\cdots u_1$. Mais on a
\[
c = c_1 * c_2 * \cdots * c_n,
\]
d'où dans $\Pi_\Sigma(X)$
\[
\hot(c) = \hot(c_n)\hot(c_{n-1})\cdots\hot(c_1),
\]
or $\hot(c_i) = \hot\,\ch(u_i) = \kappa(u_i)$, donc
\[
\begin{aligned}
\hot(c) \ &= \ \kappa(u_n)\kappa(u_{n-1})\cdots\kappa(u_1) \\
&= \ \kappa(\underbrace{u_nu_{n-1}\cdots u_1}_{u}),
\end{aligned}
\]
puisque $\kappa$ est un foncteur, donc
\[
\hot(c) = \hot\,\ch(u),
\]

\hfill c.q.f.d.
\end{enumerate}

Notons maintenant que si $(X', \Sigma')$ est un deuxième couple d'une catégorie $X'$ et d'un ensemble \textit{`admissible'} de flèches de $X'$, et si
\[
f : X \longrightarrow X'
\]
est un foncteur tel que
\[
f(\Sigma) \subseteq \Sigma',
\]
alors il existe un \textit{unique} foncteur
\[
\bar{f} = \Pi_\Sigma(f) : \Pi_\Sigma(X) \longrightarrow \Pi_{\Sigma'}(X')
\]
qui `prolonge' $f$, i.e. tel que le carré
\printedpage{48}
% head: A. Grothendieck : Dérivateurs, Ch. VII, Édition des Universités de Montpellier II et Paris VII
% foot: Version 22 décembre 2001
\msp{page 72}
\[
\begin{tikzcd}
X \arrow[r, "\kappa_X"] \arrow[d, "{\scriptstyle f}"] & \Pi_\Sigma(X) \arrow[d, "{\scriptstyle \bar{f} = \Pi_\Sigma(f)}"] \\
X' \arrow[r, "\kappa_{X'}"] & \Pi_{\Sigma'}(X')
\end{tikzcd}
\]
soit commutatif. L'unicité résulte du fait que sur les objets $\bar{f}$ doit être donné par $f$, et que sur les flèche la commutativité du carré implique que $\bar{f}$ est connu sur les flèches de la forme $\kappa_X(u)$, ainsi donc que sur les inverses de telles flèches, quand ces inverses sont inversibles, donc en particulier sur les $\kappa_X(u)^{-1}$ pour $u \in \Sigma$. Or toute flèche de $\Pi_\Sigma(X)$, provenant d'un chemin, est composé d'une suite de flèches de cette forme. Pour l'existence de $\bar{f}$, on note que $f$ induit des foncteurs % [sic] les flèche -> les flèches, plural
\[
\Ch_\tau(X; a, b) \longrightarrow \Ch_\tau(X'; f(a), f(b))
\]
qui commutent aux opérations $\varphi^*$ sur les chemins associées aux foncteurs
\[
\varphi : \tau \longrightarrow \tau',
\]
d'où un homomorphisme de catégories fibrées sur $\uTy$
\[
f_{a,b} : \Ch(X; a, b) \longrightarrow \Ch(X'; f(a), f(b)).
\]

\msp{page 73}

D'où une application sur les $\pi_0$, i.e.
\[
\bar{f}_{a,b} : \Hom_\Sigma(a,b) \longrightarrow \Hom_\Sigma(f(a), f(b)),
\]
donnée par
\[
\bar{f}_{a,b}(\hot(c)) = \hot(f_{a,b}(c)).
\]
On a évidemment, pour deux chemins $\gamma : a \longrightarrow b$ et $\gamma' : b \longrightarrow c$,
\[
f_{a,c}(\gamma\gamma') = f_{b,c}(\gamma')f_{a,b}(\gamma),
\]
ce qui prouve que $\bar{f}$ est compatible avec la composition. Il est trivial qu'il est compatible avec les identités (car $f_{a,a}$ transforme le chemin vide en $a$ en le chemin vide en $f(a)$), donc on a un foncteur $\bar{f}$, et il rend bel et bien le carré commutatif.

On peut maintenant prouver la propriété universelle de $\Pi_\Sigma(X)$ et de
\[
X \xrightarrow{\ \kappa\ } \Pi_\Sigma(X),
\]
savoir que tout foncteur $f : X \longrightarrow X'$ tel que $u \in \Sigma \Longrightarrow f(u)$ inversible, se factorise \textit{de façon unique} par $\Pi_\Sigma(X)$. L'unicité se voit comme tantôt (\footnote[20]{C'est inutile, cf. ci-dessous.}). Pour l'existence, on note que $f$ applique $\Sigma$ dans l'ensemble $\Sigma'$

% ---- from Der-ChVII-456_p051-055.tex
\setcounter{footnote}{20}

\printedpage{49}
% head: A. Grothendieck : Dérivateurs, Ch. VII, Édition des Universités de Montpellier II et Paris VII
\msp{page 74}
des flèches inversibles de $X'$ (lequel ensemble satisfait bien aux deux conditions d'admissibilité). Donc par ce qui précède on a un diagramme commutatif
\[
\begin{tikzcd}
X \ar[r,"\kappa"] \ar[d,"f"'] & \Pi_{\Sigma}(X) \ar[d,"\bar f = \Pi_{\Sigma}(f)"] \ar[dl,dotted,"\tilde f"'] \\
X' \ar[r,"\kappa'","\sim"'] & \Pi_{\Sigma'}(X'),
\end{tikzcd}
\]
où $\bar f$ est uniquement déterminé. Comme $\kappa'$ est un isomorphisme (comme on a vu tantôt, puisque $\Sigma'$ \texttt{[est]} formé des flèches inversibles), l'existence et l'unicité de $\bar f$ équivaut à la conclusion annoncée,

\hfill c.q.f.d.

\textbf{Encore un sorite} sur les classes d'homotopie d'éléments d'une catégorie $X$. De façon générale, on voit que $X$ et $\Pi_{\Sigma}(X) = X\Sigma^{-1}$ ont mêmes composantes connexes, car si
\msp{page 75}
$a$ et $b$ dans $X$ sont reliés par une flèche de $X\Sigma^{-1}$, i.e. par un $\Sigma$-chemin, alors ils sont dans une même composante connexe de $X$. Mais l'inverse est également vrai, comme on voit p.~ex. en prenant $\Sigma = \mathrm{Fl}(X)$, et en disant que dans le \textit{groupoïde} $\Pi(X)$, où deux sommets sont dans une même composante \texttt{[connexe]} si et seulement si $\exists\ x \longrightarrow y$ (relation qui est non seulement transitive et reflexive, mais aussi symétrique dans un groupoïde). Un chemin
\[
a \xrightarrow{\quad c \quad} b
\]
sera appelé aussi une \textit{homotopie} de $a$ vers $b$. On dit que $a$ et $b$ sont \textit{homotopes} s'il existe une homotopie $c$ de $a$ vers $b$, auquel cas il existe aussi une homotopie `inverse' $c^{o}$ de $b$ vers $a$. Cela signifie aussi que $a$ et $b$ sont dans la même composante connexe de $X$, i.e. définissent le même élément de $\pi_0(X)$.

\dgrblock{dgrb323}{The universal property includes natural transformations}
\begin{dcomment}{exposition}{p74}{The universal property includes natural transformations}
Grothendieck has now brought the path construction back to the familiar category of fractions. The work with arbitrary lengths and their changes of parameter pays for a universal property requiring only that the denominators contain identities and be closed under composition. No condition allowing fractions to be shortened to spans has entered the argument. This is the general localization of \citet[Chapter~I, \S1.1 and Lemma~1.2]{GabrielZisman1967}; a calculus of left or right fractions is additional structure that can simplify its representatives.

Write $Q:X\longrightarrow L=\Pi_\Sigma(X)$. A functor $F:X\longrightarrow Y$ carrying $\Sigma$ into isomorphisms must send a zigzag
\[
a=x_0\xrightarrow{u_1}x_1\xleftarrow{s_2}x_2
 \xrightarrow{u_3}\cdots\xleftarrow{s_n}x_n=b
\]
to the corresponding composite of the $F(u_i)$ and $F(s_j)^{-1}$, in their order of traversal. The preceding discussion of the evaluation map proves that this prescription is constant on components of the path category. It therefore defines $\overline F:L\longrightarrow Y$. Every arrow of $L$ is represented by a zigzag, so these prescribed values also prove uniqueness. The factorization through $\Pi_{\mathrm{Iso}(Y)}(Y)\simeq Y$ used in the text packages precisely this argument.

There is a useful categorical refinement, already part of the classical universal property. Precomposition induces an isomorphism
\[
 Q^*:\operatorname{Fun}(L,Y)\xrightarrow{\ \sim\ }
       \operatorname{Fun}_{\Sigma}(X,Y),
\]
where the right side is the full subcategory of $\operatorname{Fun}(X,Y)$ on the functors inverting $\Sigma$. The statement concerns all natural transformations, not merely natural isomorphisms. To see the extra assertion directly, let $\eta:F\Rightarrow G$ be a natural transformation between such functors. Since $Q$ is the identity on objects, the only possible extension has components $\overline\eta_x=\eta_x$. Naturality for $s:x\to y$ gives
\[
 \eta_yF(s)=G(s)\eta_x,
 \qquad
 \eta_xF(s)^{-1}=G(s)^{-1}\eta_y.
\]
The second equality supplies naturality for the formal inverse of $s$. Naturality then passes to every composite representing an arrow of $L$. These components define the required transformation, and uniqueness follows again from the identity on objects. This also explains why relative functors induce functors on localizations without further choices: both their action on arrows and their action on transformations are determined before any presentation of a fraction is chosen.

For the present chapter, this refinement matters because a path whose vertices are themselves functors is a diagram of natural transformations. The later passage to diagram categories will repeatedly use the same distinction between determining objects and controlling the morphisms between them. The localization established here already supplies both at the ordinary categorical level.
\end{dcomment}


L'homotopie sera dite \textit{simple} si $c$ est de type $\Delta_{+}$ ou $\Delta_{-}$, i.e. si elle est donnée par une flèche $u : a \longrightarrow b$ ou $v : b \longrightarrow a$. Si on cherche une homotopie entre $a$ et $b$, on essaiera de la trouver la plus `courte' possible. Si on a une homotopie $c$ de longueur $l$ de $a$ vers $b$, alors en remplaçant $\tau$ par $\tau_{\red}$, donc $I_{\tau}$ par $I_{\tau\,\red}$, le sous-intervalle réduit
\msp{page 76}
correspondant, on trouve une homotopie de $a$ à $b$ de type $\tau_{\red}$, dont la longueur est égale à la \textit{`longueur utile'} du type $\tau$ de $c$ (ce qui justifie la terminologie). Notons que si
\[
\tau = \Delta_{\varepsilon}^{l_1}\Delta_{-\varepsilon}^{l_2}\Delta_{\varepsilon}^{l_3}\cdots\Delta_{(-1)^{r}\varepsilon}^{l_r},
\]
alors
\[
\tau_{\red} = \Delta_{\varepsilon}\Delta_{-\varepsilon}\Delta_{\varepsilon}\cdots\Delta_{(-1)^{r}\varepsilon},
\]
% foot: Version 22 décembre 2001
\printedpage{50}
% head: A. Grothendieck : Dérivateurs, Ch. VII, Édition des Universités de Montpellier II et Paris VII
\texttt{[plutôt deux fois} $(-1)^{r-1}$\texttt{]} où $r$ est la longueur utile de $\tau$ (la longueur pas utile étant $n = \sum\limits_{i\in[1,r]} l_i$).

Quand on part d'une homotopie entre \textit{foncteurs}
\[
X \longrightarrow Y,
\]
il est entendu qu'il s'agit d'homotopies dans la catégorie $\uHom(X,Y)$. (Ou, suivant le contexte, dans une sous-catégorie précisée de $\uHom(X,Y)$.)

\msp{page 77}
À la page 61, j'ai invoqué l'existence d'un \textit{foncteur} de \textit{`composition des chemins'}
\[
\Ch_{\Sigma}(X;a,b)\times\Ch_{\Sigma}(X;b,c) \ \longrightarrow\ \Ch_{\Sigma}(X;a,c),
\]
qui sur la fibre du premier membre au dessus de $(\tau,\tau') \in \operatorname{Ob}(\Ty\times\Ty)$, serait donné par
\[
\Ch_{\Sigma,\tau}(X;a,b)\times\Ch_{\Sigma,\tau'}(X;b,c) \ \longrightarrow\ \Ch_{\Sigma,\tau\tau'}(X;a,c), \tag*{$(*)$}
\]
la composition des chemins de types précisés $\tau$, $\tau'$. Qu'on ait là un \textit{foncteur,} pas seulement une application sur les ensembles d'objets, est assez évident, et peut s'exprimer par exemple à l'aide du diagramme d'amalga\texttt{[ma]}tion
\[
\begin{tikzcd}[column sep=small]
 & \varepsilon = \{0,1\} \ar[dl,hook'] \ar[dr,hook] & \\
I_{\tau} \ar[dr,hook] & & I_{\tau'} \ar[dl,hook'] \\
 & I_{\tau\tau'} &
\end{tikzcd}
\]
en le transformant par $\uHom(\cdot,X)$ (au lieu de le transformer par $\mathrm{Hom}_{\mathrm{Cat}}(\cdot,X)$ seulement), du moins dans le cas particulier où $\Sigma = \mathrm{Fl}(X)$ (qui était celui de la page 61). Mais le passage aux flèches est assez évident - pour deux composés $c'_1\circ c_1$ et $c'_2\circ c_2$, de
\[
\begin{tikzcd}[column sep=small]
a \ar[rr,"c_1"] \ar[d,equal] & {} \ar[d,"w"'] & b \ar[rr,"c'_1"] \ar[d,equal] & {} \ar[d,"w'"] & c \ar[d,equal] \\
a \ar[rr,"c_2"] & {} & b \ar[rr,"c'_2"] & {} & c
\end{tikzcd}
\]
\texttt{[on]} voit que les homomorphismes $w''$ du premier dans le second, dans $\Ch_{\Sigma}(X;a,c)$, qui sont l'identité sur le sommet d'indice $n+1$ des deux chemins (où $n = \operatorname{long}\tau$), correspondant à $1_{\tau}$ et à $0_{\tau'}$, `sont' exactement les couples
% foot: Version 22 décembre 2001
\printedpage{51}
% head: A. Grothendieck : Dérivateurs, Ch. VII, Édition des Universités de Montpellier II et Paris VII
\msp{page 78}
formés d'un homomorphisme $w : c_1 \longrightarrow c_2$ et $w' : c'_1 \longrightarrow c'_2$ dans $\Ch_{\Sigma,\tau}(a,b)$ et $\Ch_{\Sigma,\tau'}(b,c)$ respectivement, et si $w''$, $w$, $w'$ se correspondent ainsi, on posera
\[
w'' = w * w'.
\]
Il y a une transitivité évidente à vérifier pour voir que ça définit un \textit{foncteur} $(*)$. On peut exprimer la situation à l'aide du diagramme
\[
\begin{tikzcd}[column sep=small]
 & \Ty\times\Ty \ar[rrr,"\pi\ (\text{composition})"] \ar[dl,"p_1"'] \ar[dr,"p_2"] & & & \Ty \\
\Ty & & \Ty, & &
\end{tikzcd}
\]
où $p_1$, $p_2$ sont les deux projections. On considère sur $\Ty\times\Ty$ les trois catégories fibrés % [sic] fibrés -> fibrées, agreement with "catégories"
\[
p_1^*(\CH_{\Sigma,\tau}(a,b)),\ \ p_2^*(\CH_{\Sigma,\tau'}(b,c)),\ \ \pi^*(\Ch_{\Sigma,\tau\tau'}(a,c))
\]
\texttt{[plutôt} $\pi^*(\CH_{\Sigma,\tau\tau'}(a,c)$\texttt{]}, dont les fibres en $(\tau,\tau')$ sont respectivement % [sic] \pi^*(\CH_{\Sigma,\tau\tau'}(a,c) -> \pi^*(\CH_{\Sigma,\tau\tau'}(a,c)), one closing parenthesis missing
\[
\Ch_{\Sigma,\tau}(a,b),\ \ \Ch_{\Sigma,\tau'}(b,c),\ \ \Ch_{\Sigma,\tau\tau'}(a,c).
\]
(On écrit $\CH$ au lieu de $\Ch$, pour indiquer que l'on considère les catégories envisagées comme des catégories (fibrées) \textit{au dessus de} $(\Ty)$ (\footnote{Mais il s'agit pour $p_1^*$, $p_2^*$ et $\pi^*$ de produits cartésiens ordinaires dans Cat, il n'y a pas à avoir peur et à écrire des $\CH$ \ldots}).) Au dessus de $\Ty\times\Ty$, on a un homomorphisme (cartésien) de catégories fibrées
\[
p_1^*(\CH_{\tau}(a,b))\times_{\Ty\times\Ty} p_2^*(\CH_{\tau}(b,c)) \ \longrightarrow\ \pi^*(\CH_{\tau}(a,c)) \tag*{$(**)$}
\]
\texttt{[plutôt} $p_1^*(\CH_{\tau}(a,b))\times_{\Ty\times\Ty} p_2^*(\CH_{\tau'}(b,c)) \ \longrightarrow\ \pi^*(\CH_{\tau\tau'}(a,c))$\texttt{]},
\msp{page 79}
car c'est OK sur les fibres ($(*)$ page 77), et si
\[
(\varphi,\varphi') : (\tau_1,\tau'_1) \longrightarrow (\tau,\tau') \qquad (\varphi : \tau_1 \longrightarrow \tau,\ \varphi' : \tau'_1 \longrightarrow \tau')
\]
est une flèche dans $\Ty\times\Ty$, on en déduit un diagramme commutatif
\[
\begin{tikzcd}
\Ch_{\Sigma,\tau}(a,b)\times\Ch_{\Sigma,\tau'}(b,c) \ar[r,"\text{comp.}"] \ar[d,"\varphi^*_{a,b}\times(\varphi'_{b,c})^*"] & \Ch_{\Sigma,\tau}(a,c) \ar[d,"(\varphi*\varphi')^*_{a,c}"] \\
\Ch_{\Sigma,\tau_1}(a,b)\times\Ch_{\Sigma,\tau'_1}(b,c) \ar[r,"\text{comp.}"] & \Ch_{\Sigma,\tau_1\tau'_1}(a,c)
\end{tikzcd}
\]
% foot: Version 22 décembre 2001
\printedpage{52}
% head: A. Grothendieck : Dérivateurs, Ch. VII, Édition des Universités de Montpellier II et Paris VII
\texttt{[plutôt} $\Ch_{\Sigma,\tau\tau'}(a,c)$\texttt{]}, i.e. la composition des chemins d'extrémités fixés est compatible avec les images inverses de chemins par flèches entre les types de chemins.

\dgrblock{dgrb324}{Connectedness does not impose an orientation}
\begin{dcomment}{examples}{p76}{Connectedness does not impose an orientation}
Grothendieck next isolates the least information carried by his paths: whether two objects belong to one component. For any admissible $\Sigma$, the map $Q:X\to X[\Sigma^{-1}]$ induces a bijection on connected components. Indeed, every arrow of $X$ supplies an edge between its two objects in the localization, so a zigzag in $X$ remains a zigzag there. Conversely, an arrow of the localization is represented by a $\Sigma$-zigzag in $X$; forgetting which of its arrows were designated denominators gives an ordinary zigzag in $X$. Applying this observation to the finitely many arrows of a zigzag in the localization proves the converse on components.

The orientation must be kept separate from this argument. Let $X=[1]$, with its single nonidentity arrow $0\to1$, and let $\Sigma$ contain only identities. Then $X[\Sigma^{-1}]=X$ and $0,1$ belong to the same component, although there is no arrow $1\to0$ in the localization. Thus the phrase about the converse on manuscript page~75 is read as a statement about connected components. It would be false if read as the assertion that connectedness supplies a directed $\Sigma$-path in every prescribed direction. When $\Sigma=\mathrm{Fl}(X)$, reversal is available, and the groupoid formulation does permit that stronger assertion.

Reduction of a type has an equally concrete meaning. A string
\[
 a_0\longrightarrow a_1\longrightarrow a_2
 \longleftarrow a_3\longleftarrow a_4\longrightarrow a_5
\]
can be replaced by
\[
 a_0\longrightarrow a_2\longleftarrow a_4\longrightarrow a_5
\]
by composing the arrows in each maximal run. The signs of the three remaining edges are $+,-,+$. For $r$ runs beginning with $\varepsilon$, their signs are $(-1)^{j-1}\varepsilon$, $1\leq j\leq r$; this is the reason for the edition's correction of the last exponent on pages~76--77. If the original backward arrows belong to $\Sigma$, closure under composition ensures that the reduced backward arrows still do. The same construction in $\operatorname{Fun}(A,B)$ composes successive natural transformations of the same direction. Here a homotopy between functors means a finite zigzag of natural transformations; its endpoints need not be naturally isomorphic.
\end{dcomment}


Mais l'homomorphisme cartésien $(**)$ de catégories fibrées sur $\Ty\times\Ty$ n'est pas tout à fait la même chose qu'un anodin foncteur
\[
\Ch_{\Sigma}(a,b)\times\Ch_{\Sigma}(b,c) \ \longrightarrow\ \Ch_{\Sigma}(a,c). \tag*{$(*)$}
\]
La formule $(*)$ définit en tout cas une application sur les ensembles d'objets, puisque les ensembles d'objets de chacune des trois catégories envisagées se décompose en somme d'ensembles $\Ch_{\Sigma,\tau}(\bullet,\bullet)$, où $\tau$ parcourt $\mathrm{Ty} = \operatorname{Ob}\Ty$. Mais cela se prolonge-t-il de façon naturelle aux morphismes?

\msp{page 80}
Mais voilà que je me rappelle que ce fameux produit cartésien $p_1^*(\text{fourbi})\times p_2^*(\text{machin})$ n'est autre que le produit cartésien tout bête
\[
\Ch_{\Sigma}(a,b)\times\Ch_{\Sigma}(b,c),
\]
envoyé sur $\Ty\times\Ty$ par le produit cartésien des deux foncteurs de projection
\[
\begin{aligned}
p_{a,b} : \Ch_{\Sigma}(a,b) &\ \longrightarrow\ \ \Ty\\
p_{b,c} : \Ch_{\Sigma}(b,c) &\ \longrightarrow\ \ \Ty,
\end{aligned}
\]
d'où
\[
\begin{tikzcd}
\Ch_{\Sigma}(a,b)\times\Ch_{\Sigma}(b,c) \ar[d,"p_{a,b}\times p_{b,c}"] \\
\Ty\times\Ty
\end{tikzcd}
\]
\msp{page 81}
qui est bel et bien fibrant (comme produit cartésien de deux foncteurs fibrants) et a les fibres qu'il faut! On trouve donc bien, par le sorite général sur les foncteurs fibrants (ici sur $\Ty\times\Ty$) un foncteur de catégories fibrées sur $\Ty\times\Ty$
\[
\Ch_{\Sigma}(a,b)\times\Ch_{\Sigma}(b,c) \ \longrightarrow\ \qquad \underbrace{\pi^*(\CH_{\Sigma}(a,c))}_{=\,(\Ty\times\Ty)\times_{\Ty}\Ch_{\Sigma}(a,c)\ \ \text{(\footnotemark)}} \qquad .
\]
\footnotetext{produit fibré par $\pi : \Ty\times\Ty \longrightarrow \Ty$.}
Comme on a d'autre part un foncteur du dernier membre (produit fibré) vers le facteur $\Ch_{\Sigma}(a,c)$, on gagne. (On l'aurait retrouvé par le raisonnement direct que je m'apprêtais de faire, et que j'ai biffé - mais je suis plus heureux comme ça!)

% foot: Version 22 décembre 2001
\printedpage{53}
% head: A. Grothendieck : Dérivateurs, Ch. VII, Édition des Universités de Montpellier II et Paris VII
Pour bien faire, il faudrait prouver l'\textit{associativité} des foncteurs de composition, pas seulement au niveau des objets, par des arguments comme ceci, mais ça revient simplement à une relation du type
\[
(w''*w')*w = w''*(w'*w)
\]
pour un système de trois flèches dans
\[
\Ch_{\Sigma}(a,b),\ \Ch_{\Sigma}(b,c)\ \ \text{et}\ \ \Ch_{\Sigma}(c,d)
\]
respectivement. Ça ne peut être que trivialement
\msp{page 82}
vrai, mais fastidieux à expliciter.

\dgrblock{dgrb325}{The total composition and its associativity}
\begin{dcomment}{exposition}{p81}{The total composition and its associativity}
Grothendieck returns to composition because the localization argument uses a functor on the total path categories, including morphisms that change the type. The commentary at manuscript pages~48--50 describes the fixed-type operation. The additional verification here is that this operation respects the arrows of the Grothendieck construction. Writing these arrows explicitly also settles the associativity left to the reader on pages~81--82.

For fixed endpoints, put $P_{a,b}(\tau)=\Ch_{\Sigma,\tau}(X;a,b)$. An arrow in its total category,
\[
 (\tau,c)\longrightarrow(\sigma,d),
\]
is a pair $(\varphi,\eta)$ with $\varphi:\tau\to\sigma$ and an endpoint-fixed natural transformation $\eta:c\Rightarrow d\varphi$. If a second arrow is $(\psi,\theta):(\sigma,d)\to(\rho,e)$, its composite with the first is
\[
 (\psi\varphi,(\theta\varphi)\eta).
\]
This is the concrete form of the fibred-category formalism used in the text.

To concatenate two such arrows, concatenate the type maps and glue their transformations at the common endpoint:
\[
 ((\varphi,\eta),(\varphi',\eta'))
       \longmapsto(\varphi*\varphi',\eta*\eta').
\]
Both components at the joining object are the identity. Hence the two definitions there agree, and the naturality square for every edge comes from exactly one of the two given naturality squares. The resulting objects remain $\Sigma$-paths because concatenation does not change an edge or its direction. A vertical component of a transformation is unrestricted: $P_{a,b}(\tau)$ is a full subcategory on its allowed objects, with the endpoint condition imposed on transformations.

Compatibility with composition is checked at vertices. On the first interval the two possible composites have component $(\theta\varphi)_i\eta_i$; on the second they have the corresponding primed component; at the joining vertex both have the identity. They are therefore equal. The base maps agree by functoriality of concatenation in $\underline{\mathrm{Ty}}$. Thus the formula defines the desired functor on total categories. Equivalently,
\[
 p_1^*\Ch_\Sigma(a,b)\times_{\underline{\mathrm{Ty}}\times\underline{\mathrm{Ty}}}
 p_2^*\Ch_\Sigma(b,c)
 \cong \Ch_\Sigma(a,b)\times\Ch_\Sigma(b,c),
\]
and the cartesian functor to the pullback along concatenation can be followed by the projection to $\Ch_\Sigma(a,c)$. This is the recognition made on page~80.

For associativity, number the vertices of three paths of lengths $n,m,l$ by successive blocks in $[0,n+m+l]$. The first block keeps its indices, the second is shifted by $n$, and the third by $n+m$. Either parenthesization produces those same indices, objects, and arrows. The same statement holds for the components of a glued transformation and for the underlying type map. Consequently both parenthesizations agree as functors, and the length-zero path is a strict unit. With these standard representatives, no later strictification is required.

One diagrammatic slip deserves explicit repair. The amalgamation on page~77 is
\[
 I_\tau\amalg_{\{1_\tau=0_{\tau'}\}}I_{\tau'},
\]
a pushout over a singleton category. The displayed $\varepsilon=\{0,1\}$ would glue two pairs of endpoints. Already for two forward one-edge intervals, that different pushout has two parallel arrows between two objects, rather than the two successive arrows of $I_{\Delta_+^2}$. The joining vertex has index $n$ in the numbering $0,\ldots,n+m$; the printed $n+1$ uses the wrong index. Neither slip affects the correct componentwise construction.
\end{dcomment}


Les sueurs froides sont terminées!

En résumé, le formalisme des $\Ch_{\Sigma}(X)$ et des $\Ch_{\Sigma}(X;a,b)$ développé jusqu'à présent donne satisfaction complète sur les points suivants :

\begin{enumerate}
\item[a)] Bon formalisme de composition des chemins (qui était dégueulasse avant).

\item[b)] Traitement élégant de la construction des catégories de fractions $\mathcal{M}\Sigma^{-1}$. Mais je m'attends que ce sera moins joli qu'avant, quand il s'agira de remplacer un $(\mathcal{M},\Sigma)$ par des $(\ \underbrace{\mathcal{M}(X)}_{=\,\uHom(X,\mathcal{M})}\ ,\Sigma_X)$ ($X$ une catégorie d'indices) \texttt{[plutôt} $\uHom(X^{o},\mathcal{M})$\texttt{]}. C'est pour y pallier qu'il me faudra développer des énoncés de comparaison, entre la catégorie $\Ch_{\Sigma}$ construite ici, et d'autres dont le type dans Hot se prête aux calculs de façon plus directe, comme dans l'état premier de la question. De plus, on sera sûrement obligé d'en venir à des Ind-catégories, pour pouvoir suivre la variance de $\mathcal{M}(X)\Sigma_X^{-1}$ pour $X$ variable.

\item[c)] Bon formalisme des `homotopies' et des compositions des homotopies (ce qui est une autre façon de parler des chemins et de leur composition).
\end{enumerate}
% foot: Version 22 décembre 2001

% ---- from Der-ChVII-456_p056-060.tex
\printedpage{54}
% head: A. Grothendieck : Dérivateurs, Ch. VII, Édition des Universités de Montpellier II et Paris VII
% foot: Version 22 décembre 2001

\msp{page 83}

\dgrblock{dgrb326}{Why the category of diagrams returns}
\begin{dcomment}{context}{p82}{Why the category of diagrams returns}
The reservation at the end of this section brings the calculation back to the purpose of a derivator. Localizing $\mathcal M$ answers the question at one point. To vary coefficients over an indexing category $A$, one must also localize
\[
 \mathcal M(A)=\operatorname{Fun}(A^{\mathrm{op}},\mathcal M),
 \qquad
 \Sigma_A=\{\eta\mid\eta_a\in\Sigma\text{ for every }a\in A\}.
\]
The opposite in the edition's correction records the presheaf convention. It is a variance convention, not an alteration of the localization theorem: the latter applies to either functor category once its convention is fixed.

There is a small but decisive issue of uniformity. An object of a path category for $\mathcal M(A)$ is a finite zigzag of natural transformations with one common word of signs over all $a\in A$. An arbitrary family of zigzags in $\mathcal M$ may use a different word, or require an unbounded length, as $a$ varies. Thus a description that works perfectly for individual objects still has work to do before it describes localization naturally in the parameter $A$. The ind-objects considered later organize precisely the freedom to enlarge a finite path type while retaining its functorial meaning.

This interpretation agrees with Grothendieck's subsequent account in his letter to Thomason of 2 April 1991. In \S2, printed page~9, he explicitly defines the associated prederivator by $A\mapsto\mathcal M(A)[\Sigma_A^{-1}]$, with presheaf diagrams and objectwise denominators; on printed page~8 he emphasizes the transformations between restriction functors as part of the structure \citep{Grothendieck1991Letter}. The present concern about types is a concrete instance of that requirement. It asks for controlled variation of the entire diagram theory, beyond the universal property of a single localization.
\end{dcomment}


\section*{8\quad Le lemme d'homotopie pour les morphismes d'intervalles}

C'est le suivant :

\textbf{Lemme 1.} \textit{Soient $\tau,\tau'\in\mathrm{Ty}$, et $f,g:\tau\rightrightarrows\tau'$ deux morphismes dans $\uTy$. Alors $f$, $g$ sont homotopes dans la catégorie $\uHom\,\underline{\mathrm{int}}(I_\tau,I_{\tau'})$ des morphismes croissants d'intervalles orientés de $I_\tau$ dans $I_{\tau'}$.}

Cette catégorie pourrait licitement se désigner par $\uHom_{\uTy}(\tau,\tau')$, ou $\uHom(\tau,\tau')$. C'est une sous-catégorie de la catégorie $\uHom(I_\tau,I_{\tau'})$ de \textit{tous} les foncteurs de $I_\tau$ dans $I_{\tau'}$, et une sous-catégorie pleine de la catégorie $\uHom^\varepsilon(I_\tau,I_{\tau'})$ des foncteurs $f$ qui sont compatibles avec la donnée des sources et des buts, i.e. qui rendent commutatif le diagramme
\[
\begin{tikzcd}[column sep=small]
 & \varepsilon=\{0,1\} \arrow[dl, hook', "{(\sigma_\tau,\beta_\tau)}"'] \arrow[dr, hook, "{(\sigma_{\tau'},\beta_{\tau'})}"] & \\
I_\tau \arrow[rr, "f"] & & I_{\tau'}.
\end{tikzcd}
\]

Si $f$, $g$ sont dans cette catégorie (pas forcément morphismes croissants d'ailleurs), une flèche de $f$ dans $g$ est donc la donnée, pour tout $i\in I_\tau$, d'une flèche $f(i)\xrightarrow{w_i}g(i)$ dans $I_{\tau'}$, avec les deux conditions :

\begin{enumerate}[label={}, leftmargin=*]
\item[a)] Si $i\in\{0_\tau,1_\tau\}$, $w_i=\mathrm{id}$, et
\item[b)] le carré commutatif habituel pour une flèche $i\longrightarrow j$ dans $I_\tau$. (Il suffit d'avoir cette compatibilité pour les flèches de transition $\ul{u}_i:i-1\longrightarrow i$ dans $I_\tau$ ($i\in[1,\operatorname{long}\tau]$).)
\end{enumerate}

Une autre façon de décrire cette catégorie $\uHom^\varepsilon$, c'est de dire que c'est la catégorie
\msp{page 84}
\[
\underline{\mathrm{Ch}}_\tau(I_{\tau'};0_{\tau'},1_{\tau'})\qquad(\simeq\uHom^\varepsilon(I_\tau,I_{\tau'})),
\]
mais $\uHom_{\uTy}(\tau,\tau')$ est la sous-catégorie pleine de cette catégorie formée des \textit{chemins croissants} dans $I_{\tau'}$. Le lemme d'homotopie signifie aussi que cette catégorie $\uHom_{\uTy}(\tau,\tau')$ est toujours \textit{connexe.} Elle n'est d'ailleurs pas forcément 0-connexe, car elle peut être vide. Notons ici, pour plus tard, la

\textbf{Proposition 1 :} \textit{Soient $\tau,\tau'\in\mathrm{Ty}$, soient $r$, $r'$ leurs longueurs utiles. Pour que $\mathrm{Hom}_{\uTy}(\tau,\tau')$ soit non vide, il faut et il suffit que l'une ou l'autre des deux conditions suivantes soit satisfaite :}

\begin{enumerate}[label={}, leftmargin=*]
\item[\textit{a)}] $r\geq r'+1$

\printedpage{55}
% head: A. Grothendieck : Dérivateurs, Ch. VII, Édition des Universités de Montpellier II et Paris VII
% foot: Version 22 décembre 2001

\item[\textit{b)}] \textit{$r=r'$, et (si $r\neq0$) $\tau$ et $\tau'$ ont même direction initiale (de sorte que $\tau_{\red}=\tau'_{\red}$).}
\end{enumerate}

(Démonstration $\pm$ immédiate.)

Notons que si $\tau$, $\tau'$ n'ont pas même direction initiale, alors tout morphisme $\tau\longrightarrow\tau'$ est forcément dégénéré sur le premier intervalle critique de $I_\tau$. On aura donc $\tau=\Delta_\varepsilon^{l_0}\cdot\tau_0$ (où $l_0$ est la longueur du premier intervalle critique, et où $\tau_0$ a \texttt{[la]} direction initiale opposée de celle de $\tau$), et un isomorphisme
\[
\uHom(\tau,\tau')\simeq\uHom(\tau_0,\tau'),
\]
où cette fois $\tau_0$, $\tau'$ ont même direction initiale. On supposera donc désormais que \textit{la direction initiale de $\tau$, $\tau'$ est la même.} La condition $\uHom(\tau,\tau')\neq\emptyset$ équivaut maintenant à $\underline{r\geq r'}$, ce que nous pouvons supposer. Je vais alors décrire un élément canonique
\msp{page 85}
$F_0$ de $\uHom(\tau,\tau')$, ainsi. Soient
\[
I_1,I_2,\dots,I_r
\]
les intervalles critiques de $I_\tau$, et
\[
I'_1,I'_2,\dots,I'_{r'}\qquad(r'\leq r)
\]
ceux de $I_{\tau'}$. Le foncteur $F_0$ proviendra de morphismes d'intervalles orientés (comme
\[
\left\{
\begin{array}{ccc}
I_1 & \xrightarrow{f_1} & I'_1\\[4pt]
I_2 & \xrightarrow{f_2} & I'_2\\[4pt]
 & \dots & \\[4pt]
I_{r'} & \xrightarrow{f_{r'}} & I'_{r'}
\end{array}
\right.\qquad)
\]
et sera donc forcément trivial sur les intervalles critiques $I_\alpha$ pour $r'<\alpha\leq r$ (il y en a si et seulement si $r>r'$). Comme $I_\tau$, $I_{\tau'}$ ont même direction initiale, il s'ensuit que pour $\alpha\in[1,r']$, $I_\alpha$ et $I'_\alpha$ ont même signature (à savoir $\varepsilon_\alpha=(-1)^{\alpha+1}\varepsilon_1$, où $\varepsilon_1$ est la signature commune de $I_1$ et $I'_1$). Donc
\[
I_\alpha=\Delta^{l_\alpha}_{\varepsilon_\alpha},\qquad I'_\alpha=\Delta^{l'_\alpha}_{\varepsilon_\alpha},
\]
et on définit $f_\alpha$ comme
\[
f_\alpha=\mathfrak{f}_{\varepsilon_\alpha,l'_\alpha}*\triv_\alpha,
\]
où
\[
\mathfrak{f}_{\varepsilon_\alpha,l'_\alpha}:\Delta_{\varepsilon_\alpha}\longrightarrow\Delta^{l'_\alpha}_{\varepsilon_\alpha}
\]
et où
\[
\triv_\alpha:\Delta^{l_{\alpha-1}}_{\varepsilon_\alpha}\longrightarrow1. % [sic] l_{\alpha-1} -> l_\alpha-1, the exponent is printed as l with subscript \alpha-1, while I_\alpha = \Delta^{l_\alpha} = \Delta * \Delta^{l_\alpha-1}
\]
En d'autres termes, pour $\varepsilon=+1$ (p. ex.) on doit définir un morphisme des intervalles type $\underbrace{\mathbf{\Delta}^+_{l_\alpha}}_{=[0,l_\alpha]}\longrightarrow\underbrace{\mathbf{\Delta}^+_{l'_\alpha}}_{=[0,l'_\alpha]}$, et on pose (avec les numérotations internes à $\mathbf{\Delta}_{l_\alpha}$, $\mathbf{\Delta}_{l'_\alpha}$)
\[
f_\alpha(0)=0,\qquad f_\alpha(i)=l'_\alpha\ \text{ pour }i\in[1,l_\alpha].
\]

\dgrblock{dgrb327}{The two orders in the interval lemma}
\begin{dcomment}{exposition}{p84}{The two orders in the interval lemma}
Grothendieck now wants the choice of a type map to disappear after homotopy. The statement is about a particularly constrained mapping category, and its constraints explain both the nonemptiness criterion and the proof. Number the source and target vertices by $[0,n]$ and $[0,m]$, and write their edge signs as $\varepsilon_i$ and $\delta_j$. An object is a nondecreasing function $F:[0,n]\to[0,m]$ satisfying $F(0)=0$, $F(n)=m$, and
\[
 F(i-1)<F(i)\quad\Longrightarrow\quad
 \delta_{F(i-1)+1}=\cdots=\delta_{F(i)}=\varepsilon_i.
\]
The condition says that one source edge can traverse a run of target edges only when their directions agree. It cannot cross a change of direction in the target. An arrow $F\to G$ is instead defined by the categorical order of the target: for every $i$, there must be an arrow $F(i)\to G(i)$ in $I_{\tau'}$. Such arrows, when present, are unique, and their naturality is automatic because the target is a poset. In particular the ordering of maps is generally different from pointwise numerical comparison.

Let $r,r'$ count the maximal constant-sign runs. Each target run must be traversed using at least one source edge of the same sign. Distinct target runs require distinct source runs in the same order, because their signs alternate and a source run cannot pass a target turning vertex. This proves $r\geq r'$ when the initial signs agree. If the initial signs disagree, the whole first source run is forced to be constant at the initial target vertex, and the necessary inequality is $r-1\geq r'$. Conversely, under either inequality, map the first edge of each of the earliest matching source runs across the corresponding whole target run, collapse its remaining edges, and collapse every unused source run at the final target vertex. This constructs an admissible map and proves Proposition~1. If the target has no edges, there is exactly one map from every source; if the source has no edges and the target has an edge, there is none.

These observations also clarify the word ``connected'' in the text. Its assertion is that any two objects are connected by a finite zigzag of transformations; it allows the mapping category to be empty. For a nonempty mapping category, the lemma asserts that its nerve is connected. Grothendieck then asks whether a stronger contraction exists. The following example separates that question from the proposed terminal-object argument, and the development after it gives a uniform contraction.
\end{dcomment}


\printedpage{56}
% head: A. Grothendieck : Dérivateurs, Ch. VII, Édition des Universités de Montpellier II et Paris VII
% foot: Version 22 décembre 2001

\msp{page 86}
L'idée donc, c'est d'être dans les dispositions d'un voyageur qui, tenu à un plan de voyage par étapes décrit par $\tau$, voudrait arriver aussi vite que possible à $1_{\tau'}$, en partant de $0_{\tau'}$. Il va dès la première étape $\Delta_\varepsilon\subseteq I_1$, se porter jusqu'au bout de $I'_1$, regrettant ne pas pouvoir aller plus loin (à cause du plan de voyage). Il doit rester inactif jusqu'à ce qu'il soit arrivé au moment marqué par l'extrémité finale de $I_1$ (si celle-ci est au delà de celle de sa première flèche de transition, i.e. si $l_1>1$), puis il fait son deuxième grand pas jusqu'à l'extrémité terminale de $I'_2$ etc. Il peut espérer, à chaque moment de son voyage, être en avance (ou du moins pas en retard) sur les autres voyageurs, qui auraient suivi de tactiques de progression différentes. Si cet espoir était justifié, cela signifierait que $F_0$ serait un plus grand élément de $\uHom(\tau,\tau')$ (\textbf{NB} toutes ces catégories $\uHom$ entre intervalles sont ‘ordonnées’, i.e. il y a au plus une flèche entre deux objets), i.e. un objet final. Cela impliquerait que $\uHom(\tau,\tau')$ serait même contractile - chose peut-être vraie, mais sûrement pas de façon si simpliste je crains. Je me suis d'ailleurs canulé : $F_0$ est sans doute bien un ‘plus grand’ élément, mais pour l'ordre défini par la numérotation des sommets de $I_{\tau'}$ (donc $F\leq G$ si et seulement si $F(i)\leq G(i)$ $\forall i\in I_\tau$), et nullement pour l'ordre défini par la structure de catégorie de $\uHom(\tau,\tau')$.

\msp{page 87}
L'idée pour prouver la connexité de $\uHom(\tau,\tau')$, c'est de ramener tout $F\in\mathrm{Hom}(\tau,\tau')$, par une suite d'homotopies simples, à la forme normale. Rappelons que $F$ est donné par une suite croissante de flèches de transition
\[
\ul{u}_{i_1},\ \ul{u}_{i_2},\ \dots\ ,\ \ul{u}_{i_s},\qquad i_1<i_2<\cdots<i_s,
\]
et une suite d'entiers $\geq1$
\[
\lambda_1,\ \lambda_2,\ \dots\ ,\ \lambda_s
\]
avec la condition que
\[
\Delta^{\lambda_1}_{\varepsilon_{i_1}}\Delta^{\lambda_2}_{\varepsilon_{i_2}}\cdots\Delta^{\lambda_s}_{\varepsilon_{i_s}}=\tau'. \tag*{$(*)$}
\]
Les $\ul{u}_{i_\alpha}$ sont les flèches de transition de $I_\tau$ \textit{critiques} pour $F$, i.e. non dégénérées pour $F$, et $\lambda_\alpha$ est la \textit{partie} de $\ul{u}_{i_\alpha}$ pour $F$. Le foncteur normalisé (‘pressé’) $F_0$ est distingué par les conditions suivantes a) et b) \texttt{[et c)]}. Tout d'abord :

\dgrblock{dgrb328}{A quickest map that is neither initial nor terminal}
\begin{dcomment}{examples}{p86}{A quickest map that is neither initial nor terminal}
The normal form reaches each target turning vertex as early as the source permits. It maximizes numerical progress, but the categorical direction reverses whenever the target changes direction. A four-object mapping category makes the distinction visible without any large diagram.

Take
\[
 \tau=\Delta_+^2\Delta_-^2,
 \qquad \tau'=\Delta_+\Delta_-.
\]
The target is $0\to1\leftarrow2$. Every increasing endpoint-preserving functor has vertex list
\[
 F_{a,b}=(0,a,1,b,2),\qquad a\in\{0,1\},\quad b\in\{1,2\}.
\]
Indeed the middle source vertex must map to the target turning vertex $1$: it cannot lie below $1$, since the remaining negative edges could not reach $2$, and it cannot lie above $1$, since the initial positive edges could not reach it. The choices of $a,b$ are then unrestricted in the displayed ranges. The exact arrow relation is
\[
 F_{a,b}\longrightarrow F_{a',b'}
 \quad\Longleftrightarrow\quad a\leq a'\ \text{ and }\ b\geq b',
\]
where the inequalities on the right are numerical. The reversal of the second inequality is forced by $2\to1$ in the target.

The quickest functor is $F_0=F_{1,2}$. It is numerically greater than or equal to every other functor, but the terminal object of this mapping category is $F_{1,1}$ and its initial object is $F_{0,2}$. In particular $F_{1,1}$ has no arrow to $F_0$, and $F_0$ has no arrow to $F_{0,2}$. The objection on manuscript page~86 is therefore exact, even in this small case.

Nevertheless a contraction to $F_0$ is immediate once we normalize one target run at a time. Define $R(F_{a,b})=F_{1,b}$. There is a natural transformation $\mathrm{id}\Rightarrow R$, and another $\mathrm{const}_{F_0}\Rightarrow R$. Hence
\[
 \mathrm{id}\ \Longrightarrow\ R\ \Longleftarrow\ \mathrm{const}_{F_0}.
\]
This is the pattern of the general argument: each step advances across only one constant-sign target run, so all its components have one common categorical direction. The entire numerical maximum need not define a single transformation.
\end{dcomment}


\dgrblock{dgrb329}{A uniform contraction of the interval mapping category}
\begin{dcomment}{development}{p86}{A uniform contraction of the interval mapping category}
Grothendieck's proposed normal form can be made into a contraction by advancing across one target run at a time. This answers the question on page~86 for the increasing endpoint-preserving mapping categories used in this section. We did not find this particular argument in the literature checked for the present commentary: the source's \S8 proves connectedness, while the institutional abstract of \citet{FernandesKoppitzMusunthia2019} concerns order-preserving self-transformations of a fence, without these simultaneous endpoint and linear-order constraints. We checked that abstract and metadata; its manuscript could not be retrieved. The result below is an elementary supplement, without a claim of priority.

\textbf{Proposition (run-by-run contraction).}\ours\
\textit{Let $P(\tau,\tau')$ be the category of nondecreasing endpoint-preserving functors $I_\tau\to I_{\tau'}$, with endpoint-fixed natural transformations. If $P(\tau,\tau')$ is nonempty, its identity functor is connected to the constant functor at the quickest map by a finite zigzag of natural transformations. In particular its nerve is contractible. One can use at most $r'$ transformations, where $r'$ is the number of target runs.}

\textit{Proof.} We induct on $r'$. When the target has no edges, every functor is the same constant functor and $P(\tau,\tau')$ is the terminal category. Assume that the target is nontrivial. The source is then nontrivial too. If its initial sign differs from the target's, every map is constant on the entire first source run, as proved above. Every transformation is the identity at those vertices. Restriction past that run is consequently an isomorphism of mapping categories. We may use this isomorphism to assume that the first source and target runs have the same sign $\varepsilon$.

Let the first source run have vertices $0,\ldots,a$, and the first target run have vertices $0,\ldots,b$, where $a,b\geq1$. The numerical clamping map
\[
 h:[0,m]\longrightarrow[0,m],\qquad h(j)=\max\{j,b\},
\]
is a functor of the target interval category. On its first run it collapses every arrow to an identity at $b$; on every later edge it is the identity. These checks on generating edges prove functoriality, including at the turning vertex $b$.

For each $F\in P(\tau,\tau')$, define
\[
 (RF)(0)=0,\qquad (RF)(i)=h(F(i))\quad(1\leq i\leq n).
\]
The first source run cannot pass the first target turning vertex. Thus $F(i)\leq b$ for $0\leq i\leq a$, and $(RF)(i)=b$ for $1\leq i\leq a$. The new first edge is an allowed traversal from $0$ to $b$ of sign $\varepsilon$; all later edges are obtained by applying the functor $h$ to their old images. This proves that $RF$ is a functor. Its vertex map remains nondecreasing, it fixes both endpoints, and hence it is again an object of $P(\tau,\tau')$.

For a transformation $F\Rightarrow G$, apply $h$ to its components at vertices $i\geq1$ and use the identity at vertex $0$. The resulting components define a transformation $RF\Rightarrow RG$. Indeed both functors are already valid interval functors, their target is thin, and each component exists by functoriality of $h$. Identities and composites are preserved componentwise. We have therefore constructed an endofunctor $R$ of $P(\tau,\tau')$.

The comparison with the identity has one uniform direction. If $\varepsilon=+$, the unique arrows $j\to h(j)$ give $\mathrm{id}_P\Rightarrow R$; if $\varepsilon=-$, the arrows $h(j)\to j$ give $R\Rightarrow\mathrm{id}_P$. At vertex $0$ the component is the identity, and at the final vertex it is also the identity. Naturality in the source interval and in $F$ follows from thinness of the target and of the mapping category. Only values in the first target run have changed, which is why this comparison exists despite the failure of a global numerical-maximum argument.

Let $Q\subseteq P$ be the full subcategory of maps with values $0,b,\ldots,b$ on the first source run. The endofunctor $R$ takes values in $Q$ and restricts to the identity there. Delete the first source run, retain its last vertex as the new initial vertex, and do the same for the target. Restriction identifies $Q$ with the endpoint-preserving mapping category of these two remaining intervals. To verify this identification on morphisms, observe that every component over the deleted prefix is forced to be an identity; restriction and the prescribed extension are inverse functors. The remaining target has $r'-1$ runs. Since $Q$ contains $RF$ for every $F$, it is nonempty, so induction applies.

Compose the inductive zigzag on $Q$ with $R:P\to Q$ and the inclusion $Q\hookrightarrow P$. Adjoining the single comparison between $\mathrm{id}_P$ and $R$ gives a zigzag to the constant functor at the map that crosses the first target run on the first available source edge and then follows the inductively quickest map. This is exactly Grothendieck's normal form, after the forced initial deletion when necessary. The number of comparisons increases by one, hence is at most $r'$. A natural transformation induces a simplicial homotopy after taking nerves, so this categorical contraction contracts the nerve as asserted.\hfill$\square$

The proposition supplies more than a path joining two chosen maps: its choices of intermediate functors are simultaneous and natural in the object of the mapping category. It concerns the increasing subcategory. No conclusion about all endpoint-preserving functors, without numerical monotonicity, is being used here.
\end{dcomment}


\begin{enumerate}[label={}, leftmargin=*]
\item[a)] La suite des signes $\varepsilon_{i_1},\varepsilon_{i_2},\dots,\varepsilon_{i_s}$ est ‘alternée’, i.e. $\varepsilon_{i_\alpha}\neq\varepsilon_{i_{\alpha+1}}$ ($\alpha\in[1,s-1]$). (Cela traduit le fait que le voyageur pressé ne fait pas en plusieurs pas ce qu'il pourrait faire en un seul.)

Comme le premier de ces signes est forcément $\varepsilon'_1=\varepsilon_1$, cela signifie que la suite des signes est la même que celle des $\varepsilon_1$, $\varepsilon_2$, \dots, $\varepsilon_{r'}$ (à la longueur près). Mais mieux, que l'expression $(*)$ de $\tau'$ est son unique expression réduite, donc que $s=r'$ (longueur utile de $\tau'$), et que les $\lambda_i$ sont les longueurs $l'_i$ des intervalles critiques de $I_{\tau'}$.

\printedpage{57}
% head: A. Grothendieck : Dérivateurs, Ch. VII, Édition des Universités de Montpellier II et Paris VII
% foot: Version 22 décembre 2001

\msp{page 88}
Soit donc $\alpha\in[1,s=r']$, et considérons la flèche de transition $F$-critique $\ul{u}_{i_\alpha}$, soit $j_\alpha$ l'indice marquant l'intervalle critique de $I_\tau$ contenant la partie de $\ul{u}_{i_\alpha}$ dans $I_\tau$,
\[
\ul{u}_{i_\alpha}\in\mathrm{Fl}(I_{j_\alpha}),
\]
alors $F$ induit des morphismes (croissants) d'intervalles orientés
\[
F_\alpha=F|I_{j_\alpha}:I_{j_\alpha}\longrightarrow I'_\alpha
\]
et est trivial sur les autres intervalles critiques (ceux pas de la forme $I_{j_\alpha}$). La donnée des $F_\alpha$ détermine donc $F$ sur les $I_{j_\alpha}$, donc sur \textit{tous} les $I_j$, donc partout (par l'opération $*$ de composition par recollement). De plus la condition a) assure que chaque $F_\alpha$ est de nature assez particulière : une seule flèche de transition critique $\ul{u}_{i_\alpha}$. De plus, pour qu'on ait $F=F_0$, il faudra :

\item[b)] $\forall\alpha\in[1,r']$, $\ul{u}_{i_\alpha}$ est la \textit{première} flèche de transition de l'intervalle critique $I_{j_\alpha}$ qui la contient.

Et enfin :

\item[c)] La suite des intervalles critiques de $I_\tau$
\[
I_{j_1},\ I_{j_2},\ \dots\ ,\ I_{j_{r'}}\qquad(\text{de signes alternés }\varepsilon_1,\ -\varepsilon_1,\ \varepsilon_1,\ \ldots)
\]
associés aux flèches $F$-critiques, n'est autre que celle
\[
I_1,\ I_2,\ \dots\ ,\ I_{r'}
\]
des $r'$ premiers intervalles critiques de $I_\tau$, i.e.
\[
(j_1,j_2,\dots,j_{r'})=(1,2,\dots,r').
\]
\end{enumerate}

\msp{page 89}
Ceci vu, l'idée naturelle pour ‘déformer’ $F$ en $F_0$, c'est de trouver une suite d'homotopies qui transforment $F$ de façon à lui faire acquérir successivement les caractéristiques a), b), c) du voyageur pressé. Au besoin, pour aller plus vite au but, on utilisera l'hypothèse de récurrence sur $n=\operatorname{long}\tau$, que le théorème de connexité est déjà prouvé quand $I$ est de longueur $\leq n-1$. (NB Ce théorème est trivial pour $n\leq1$, car alors les catégories $\uHom(\tau,\tau')$ sont toutes ou bien vides, ou bien ponctuelles.)

\msp{page 90}
1) On veut assurer a). Pour ceci, on utilise l'hypothèse de récurrence, pour remplacer un groupe de \texttt{[mot illisible]} flèches critiques consécutives de même signe, par une seule flèche critique. Cette méthode ne foire que si les flèches critiques sont toutes de même signe $\varepsilon$ et incluent $\ul{u}_1$ et $\ul{u}_n$. Mais l'hypothèse de récurrence appliquée à $I=[1,n-1]\subseteq I_\tau$ nous ramène au cas où il n'y a que deux flèches critiques $\ul{u}_1$ et $\ul{u}_n$. Si le premier intervalle critique n'est pas réduit à $\Delta_\varepsilon$, i.e. soit de longueur 2, alors l'hypothèse de récurrence appliquée à $I=[1,n]\subseteq I_\tau$ nous ramène au cas où les deux flèches critiques sont $\ul{u}_1$ et $\ul{u}_2$, et on est
\printedpage{58}
% head: A. Grothendieck : Dérivateurs, Ch. VII, Édition des Universités de Montpellier II et Paris VII
% foot: Version 22 décembre 2001
\msp{page 91}
alors ramené au cas où $\boxed{I=\Delta^2_\varepsilon,\ \tau'=\Delta^2_\varepsilon.}$ Sinon, au cas où $\boxed{\tau=\Delta_\varepsilon\Delta^{l_2}_{-\varepsilon}\Delta_\varepsilon,\ \tau'=\Delta^2_\varepsilon,}$ et où les deux flèches critiques sont encore la première et la dernière, et $\tau'=\Delta^2_\varepsilon$. Ce sont là les deux cas irréductibles, qu'il faudra traiter à part.

Ces cas admis, on a donc moyen pour définir $F$ de façon à satisfaire a). Ensuite on veut :

2) Assurer que la suite des intervalles critiques $I_{j_1}$, $I_{j_2}$, \dots, $I_{j_{r'}}$ qui contient les flèches critiques $\ul{u}_{i_1}$, \dots, $\ul{u}_{i_{r'}}$, soit la suite $I_1$, \dots, $I_{r'}$. C'est encore immédiat par l'hypothèse de récurrence, le seul cas irréductible étant celui où $r'=1$ (une seule flèche critique), et où il s'agit de faire descendre cette flèche critique jusqu'à $\ul{u}_1$. Si $I_1$ est de longueur $>1$, alors on fait descendre en $\ul{u}_2$ par hypothèse de récurrence, et le cas irréductible est celui où $\boxed{\tau=\Delta^2_\varepsilon,\ \tau'=\Delta_\varepsilon,}$ une flèche critique $\ul{u}_2$, la faire descendre en $\ul{u}_1$. Si par contre $I_1$ \texttt{[est]} de longueur 1, le cas irréductible est $\boxed{\tau=\Delta_\varepsilon\Delta^{l_2}_{-\varepsilon}\Delta_\varepsilon,\ \tau'=\Delta_\varepsilon,}$ une flèche critique $\ul{u}_{l_2+2}$, la faire descendre en $\ul{u}_1$.

Sous réserve de traiter les deux nouveaux cas irréductibles, on peut déformer $F$ en un foncteur qui satisfait à a) et à c). Reste à normaliser chacun des $F_\alpha:I_\alpha\longrightarrow I'_\alpha$ ($\alpha\in[1,r']$). Donc on est ramené à $\Delta^n_\varepsilon\longrightarrow\Delta_\varepsilon$ : une seule flèche critique $\ul{u}_\alpha$, $\alpha\in[1,n]$, à faire descendre en $\ul{u}_1$. Si $n\geq2$, par hypothèse de récurrence elle descend en $\ul{u}_2$, cas supposé traité en $2^o$). Donc
\msp{page 92}
on gagne, sous réserve de traiter les 4 cas irréductibles.

A) Cas $\tau=\Delta^2_\varepsilon$, $\tau'=\Delta^2_\varepsilon$ ou $\Delta_\varepsilon$.

La catégorie $\uHom(\Delta^2_\varepsilon,\Delta_\varepsilon)$ a deux éléments, les morphismes $\Delta^2_\varepsilon\longrightarrow\Delta_\varepsilon$ correspondent aux flèches de transition de $\Delta^2_\varepsilon$ ($=\mathbf{\Delta}^+_2$ ou $\mathbf{\Delta}^-_2$). Dans le cas $\varepsilon=+1$, ce sont les deux applications de $\{0,1,2\}$ dans $\{0,1\}$, qui appliquent 0 en 0, 2 en 1, et 1 soit en 1 ($F_0$), soit en 1 \texttt{[plutôt en 0]} ($F_1$).
\[
\begin{tikzcd}
0 \arrow[r] & 1 \arrow[r] & 2
\end{tikzcd}
\]
\[
\begin{tikzcd}
0 \arrow[r] & 1
\end{tikzcd}
\]

On a évidemment une flèche $F_1\longrightarrow F_0$ (le voyageur pressé arrive plus vite \ldots). Similaire pour $\varepsilon=-1$,
\[
\uHom(\Delta^2_\varepsilon,\Delta_\varepsilon)\simeq\mathbf{\Delta}^1,\qquad\text{catégorie 0-connexe.}
\]

La catégorie $\uHom(\Delta^2_\varepsilon,\Delta^2_\varepsilon)$ a trois éléments, s'identifiant à l'ensemble des parties non vide de l'ensemble $\mathrm{Tr}(\tau)$ à deux éléments des flèches de transition $\ul{u}_1$, $\ul{u}_2$ de $\tau$ : il y a $F_{\ul{u}_1}=F_0$, $F_{\ul{u}_2}$, $F_{\ul{u}_1,\ul{u}_2}$. Si $\varepsilon=+1$, il s'agit des applications croissantes (sans plus) de $\{0,1,2\}$ dans $\{0,1,2\}$, appliquant 0 en 0, 2 en 2, et 1 en ce qu'elle voudra. % [sic] parties non vide -> parties non vides, agreement

% ---- from Der-ChVII-456_p061-065.tex
\printedpage{59}
% head: A. Grothendieck : Dérivateurs, Ch. VII, Édition des Universités de Montpellier II et Paris VII
% foot: Version 22 décembre 2001

\begin{tabular}{lll}
Donc & $1\longmapsto 2$ & correspond à \quad $F_0$\\[1ex]
 & $1\longmapsto 1$ & correspond à \quad $F_1$\\[1ex]
 & $1\longmapsto 0$ & correspond à \quad $F_2$
\end{tabular}

et on a bien sur $F_2\longrightarrow F_1\longrightarrow F_0$, % [sic] bien sur -> bien sûr, accent missing
\[
\uHom(\Delta^2_\varepsilon,\Delta^2_\varepsilon)\simeq\boldsymbol{\Delta}^2,\qquad\quad\text{catégorie 0-connexe.}
\]

Dans ces deux cas, deux éléments de $\mathrm{Hom}(\tau,\tau')$ sont reliés par une homotopie \textit{simple.}

\msp{page 93}

B) Cas $\tau=\Delta_\varepsilon\Delta^l_{-\varepsilon}\Delta_\varepsilon$, $\tau'=\Delta^2_\varepsilon$ ou $\Delta_\varepsilon$.

Notons que l'on a, pour $\tau,\tau'\in\uTy$ quelconque, un isomorphisme de catégories
\[
\uHom(\tau,\tau')\longrightarrow\uHom(\tau'^o,\tau^o),
\]
exprimant simplement que
\[
\tau\longmapsto\tau^o
\]
est un \texttt{[2]}-foncteur de $\uTy$ dans $\uTy$, et même un automorphisme de cette \texttt{[2]}-catégorie. Donc dire que $\uHom(\tau,\tau')$ est connexe revient à dire que $\uHom(\tau^o,\tau'^o)$ l'est. Donc on peut supposer ci-dessus que $\varepsilon=+1$ (et pareil dans A!).
\[
\begin{aligned}
\tau :\quad & \longrightarrow\ \underbrace{\longleftarrow\ \longleftarrow\ \cdots\ \longleftarrow}_{l}\ \longrightarrow\\
\tau' :\quad & \longrightarrow\ \longrightarrow\quad\text{ou}\quad\longrightarrow\ .
\end{aligned}
\]

Mais tous les morphismes $\tau\longrightarrow\tau'$ sont triviaux sur l'intervalle critique $I_2$ (rétrograde), ce qui implique que
\[
\uHom(\tau,\tau')\xleftarrow{\ \sim\ }\uHom(\Delta^2_\varepsilon,\tau'),
\]
via le foncteur de contraction canonique
\[
I_\tau\longrightarrow I_{\Delta^2_\varepsilon}=\boldsymbol{\Delta}^2.
\]

Or dans A on a vu que $\uHom(\Delta^2_\varepsilon,\tau')$ est connexe (pour $\tau'$ de la forme $\Delta_\varepsilon$ ou $\Delta^2_\varepsilon$).

On a gagné!

\newpage
\printedpage{60}
% head: A. Grothendieck : Dérivateurs, Ch. VII, Édition des Universités de Montpellier II et Paris VII
% foot: Version 22 décembre 2001

\msp{page 94}

\dgrblock{dgrb330}{The four small cases and the completed proof}
\begin{dcomment}{exposition}{p93}{The four small cases and the completed proof}
Grothendieck's induction reduces the normalization to two monotone cases and two cases with one intervening run of the opposite sign. The local calculations are sound; the reversal formula printed at the start of case~B needs its arguments restored to their original order. Reversal of the numbered interval gives
\[
 P(\tau,\tau')\cong P(\tau^o,\tau'^o),
 \qquad F\longmapsto r_{\tau'}F r_\tau^{-1},
\]
where $r_\tau:I_\tau\to I_{\tau^o}$ reverses the vertex numbering. It is a covariant isomorphism of the underlying categories, since the reversed word also changes every sign. The displayed $P(\tau'^o,\tau^o)$ in the edition would reverse the direction of a functor and is not this operation. The next sentence in the source already has the correct order.

For $\tau=\Delta_+^2$ and target $\Delta_+^k$ with $k=1$ or $2$, a map is determined by $F(1)=j$, $0\leq j\leq k$. There is a unique transformation from $j$ to $j'$ exactly when $j\leq j'$. Thus the two mapping categories are the chains $[1]$ and $[2]$. With negative signs the order reverses, producing isomorphic chains. In particular any two objects are joined by one simple homotopy. These are the calculations on pages~92--93, including the edition's correction of the repeated value $1$ in the first example.

For $\tau=\Delta_+\Delta_-^l\Delta_+$ and a positive target, every edge of the negative middle run must collapse: a nondecreasing map cannot carry a negative source edge to a nonconstant positive target arrow. It therefore factors uniquely through the contraction $I_\tau\to I_{\Delta_+^2}$. The factorization also holds for transformations, because all components over the collapsed middle run coincide. This proves the displayed isomorphism of mapping categories and reduces the remaining cases to the preceding chains. Two small corrections clarify the longer induction on page~90. A first source run that is not a singleton has length at least two; the printed ``length two'' omits that qualification. More substantially, to concentrate all activity except the final edge, the inductive subinterval must be $[0,n-1]$, not the printed $[1,n-1]$. The latter fixes both interior boundary values, and cannot eliminate a nonzero traversal between them: the identity on a positive three-edge interval already demonstrates the problem. On the prefix $[0,n-1]$, induction instead leaves its first edge active and its remaining edges constant, while the final edge is unchanged. The following use of $[1,n]$ then moves the tail activity to the second edge when the first run has length at least two. With these corrections the stated reductions have their intended meaning. The run-by-run contraction above gives a full alternative proof of their connectedness conclusion.
\end{dcomment}


\dgrblock{dgrb331}{The path model and simplicial localization}
\begin{dcomment}{later}{p93}{The path model and simplicial localization}
The interval lemma makes changes of parameter homotopically harmless in a precise local sense. It is natural to compare the resulting path constructions with Dwyer--Kan localization, but the categories being compared must first be specified. In the present text, $\Ch_{\Sigma,\tau}(X;a,b)$ is full on the diagrams whose backward edges lie in $\Sigma$; its interior vertical transformations are arbitrary arrows of $X$. In the hammock construction, the vertical arrows are required to lie in the chosen weak-equivalence subcategory as well. These prescriptions agree when $\Sigma=\mathrm{Fl}(X)$, and can differ for a proper subcategory of denominators.

\citet[Definition~2.1, Proposition~2.2 and Proposition~3.1]{DwyerKan1980Calculating} construct the hammock mapping simplicial sets, compare them with simplicial localization, and identify their components with ordinary localization. Their definition also permits strict composition by concatenation followed by reduction. This gives a useful historical comparison: their work appeared in 1980, before these notes of 1990. It supplies an established construction retaining homotopical information about fractions, but its comparison theorem does not identify the nerve of the particular full path category in this manuscript without an additional argument.

Here the elementary result is already valuable on its own. It says that a category of allowed parameter maps carries a contraction whenever it has an object. The next section transports natural transformations between such maps through precomposition on path diagrams. That step explains why Grothendieck needed the homotopies to remain inside the increasing endpoint-preserving category, rather than merely to exist between maps of the underlying realizations.
\end{dcomment}


\section*{9\quad Théorème d'homotopie pour les catégories de chemins}

Soient $X$ une catégorie, $\Sigma\subseteq\mathrm{Fl}(X)$ \texttt{[un]} ensemble admissible de flèches dans $X$, $a,b\in\mathrm{Ob}\,X$. Considérons pour $\tau\in\mathrm{Ty}$ variable, les catégories
\[
\uCh_{\Sigma,\tau}(X),\qquad\quad\uCh_{\Sigma,\tau}(X;a,b),
\]
qui forment les fibres des catégories Cat-fibrées
\[
\begin{tikzcd}[column sep=small]
\uCh_\Sigma(X) \arrow[dr] & & \uCh_\Sigma(X;a,b) \arrow[ll, "\text{incl.}"'] \arrow[dl]\\
 & \uTy &
\end{tikzcd}
\]
sur la catégorie $\uTy$. Ainsi, pour $\tau,\tau'\in\mathrm{Ty}$, et pour toute
\[
\varphi:\tau\longrightarrow\tau',
\]
on a un carré commutatif d'homomorphismes
\[
\begin{tikzcd}
\uCh_{\Sigma,\tau'}(X) \arrow[d, "\varphi^*_X"'] & \uCh_{\Sigma,\tau'}(X;a,b) \arrow[l, "\text{incl.}"'] \arrow[d, "\varphi^*_{a,b}"]\\
\uCh_{\Sigma,\tau}(X) & \uCh_{\Sigma,\tau}(X;a,b), \arrow[l, "\text{incl.}"']
\end{tikzcd}
\]
avec transitivité pour $\varphi$ variable (composition $\psi\circ\varphi$ de flèches dans $\uTy$).

En fait, pour $\varphi$ variable dans $\uHom(\tau,\tau')$, les foncteurs images inverses $\varphi^*_X$, $\varphi^*_{a,b}$ dépendent \textit{fonctoriellement} de $\varphi$; i.e. on a des \textit{foncteurs}
\[
\begin{tikzcd}[row sep=small]
 & \uHom(\uCh_{\Sigma,\tau'}(X),\uCh_{\Sigma,\tau}(X))\\
\uHom(\tau,\tau') \arrow[ur] \arrow[dr] & \\
 & \uHom(\uCh_{\Sigma,\tau'}(X;a,b),\uCh_{\Sigma,\tau}(X;a,b))
\end{tikzcd}
\tag*{$(*)$}
\]

\newpage
\printedpage{61}
% head: A. Grothendieck : Dérivateurs, Ch. VII, Édition des Universités de Montpellier II et Paris VII
% foot: Version 22 décembre 2001

\msp{page 95}

\dgrblock{dgrb332}{Why a homotopy of types acts on every path}
\begin{dcomment}{exposition}{p94}{Why a homotopy of types acts on every path}
Grothendieck now uses the interval calculation to remove a choice from the theory: the operation induced on paths should depend on the source and target types, rather than on the particular subdivision or collapse chosen between them. The relevant assertion concerns natural transformations, not merely the induced maps of connected components. This stronger form is what will allow a diagram of categories to descend to a diagram in $\Cathot$.

Write $P_\tau=\uCh_{\Sigma,\tau}(X)$ and let $\mathcal T(\tau,\tau')$ be the category of endpoint-preserving type maps used in the preceding section. If $\eta:\varphi\Rightarrow\psi$ is a morphism of this category and $F:I_{\tau'}\to X$ is a permitted path, the transformation
\[
 F\eta:F\varphi\Longrightarrow F\psi
\]
is a morphism in $P_\tau$. Indeed, both its objects are permitted paths because $\varphi$ and $\psi$ carry backward arrows to backward arrows or identities. The fibre $P_\tau$ is the \emph{full} subcategory of the functor category on these objects: the components of $F\eta$ need not belong to $\Sigma$. Naturality in $F$ is precisely the interchange law for natural transformations. Thus whiskering gives a functor
\[
 \mathcal T(\tau,\tau')\longrightarrow\uHom(P_{\tau'},P_\tau),
 \qquad \varphi\longmapsto\varphi^*.
\]
This is the elementary two-dimensional mechanism behind the theorem.

The endpoint condition does a second piece of work. For a transformation of type maps its two endpoint components are identities, so $(F\eta)_0=\id_{F(0)}$ and $(F\eta)_1=\id_{F(1)}$. The displayed functor therefore lands in the category of functors and vertical transformations over $X\times X$. Restriction to the strict fibre at $(a,b)$ consequently preserves the entire homotopy. An arbitrary homotopy of functors over an unspecified base would not have this property. The words ``au dessus de $X$'' on manuscript page~97 should here be understood as ``over $X\times X$'', as the following diagram makes explicit.

Finally, connectedness of $\mathcal T(\tau,\tau')$, proved by the interval homotopy lemma, supplies a finite zigzag from $\varphi$ to $\psi$. Applying the whiskering functor supplies the required finite zigzag of natural transformations. Taking nerves turns each transformation into a simplicial homotopy. No assertion that the individual pullback functors are equivalences is made: the theorem compares parallel operations, and this is exactly what the next construction needs.
\end{dcomment}


Cela provient de l'inclusion
\[
\uHom(\tau,\tau')\hookrightarrow\uHom^{\varepsilon}(I_\tau,I_{\tau'})
\]
et de la description des catégories $\uCh_\tau$, $\uCh_{\tau'}$ comme catégories de foncteurs, telle
\[
\uCh_{\Sigma,\tau}(X)\simeq\uHom_{\Sigma_\tau,\Sigma}(I_\tau,X),
\]
où $\Sigma_\tau$ désigne l'ensemble des flèches dans $I_\tau$ qui sont soit inversibles, soit \textit{rétrogrades.} (Cet ensemble de flèches est admissible, car stable par composition.) On désigne par $\uHom_{\Sigma_\tau,\Sigma}$ la sous-catégorie pleine de $\uHom$ formée des foncteurs qui appliquent $\Sigma_\tau$ dans $\Sigma$. Or de façon générale, si
\[
(I,\Sigma_I),\ \ (I',\Sigma_{I'}),\ \ (X,\Sigma)
\]
sont trois catégories, munies d'ensembles admissibles de flèches, considérant les foncteurs $\varphi:I\longrightarrow I'$ avec $\varphi(\Sigma_I)\subseteq\Sigma_{I'}$, alors le foncteur familier
\[
\uHom(I,I')\longrightarrow\uHom(\uHom(I',X),\uHom(I,X))
\]
(transposé du foncteur de composition
\[
\begin{matrix}
\uHom(I,I') & \times & \uHom(I',X) & \longrightarrow & \uHom(I,X) & \\
(\varphi & , & F') & \longmapsto & F'\circ\varphi & ),
\end{matrix}
\]
se généralise en un foncteur
\[
\uHom_{\Sigma_I,\Sigma_{I'}}(I,I')\longrightarrow\uHom(\uHom_{\Sigma_{I'},\Sigma}(I',X),\uHom_{\Sigma_I,\Sigma}(I,X)),
\tag*{$(*)$}
\]
transposé du foncteur de composition
\[
\uHom_{\Sigma_I,\Sigma_{I'}}(I,I')\times\uHom_{\Sigma_{I'},\Sigma}(I',X)\longrightarrow\uHom_{\Sigma_I,\Sigma}(I,X)
\]
induit par $(*)$.

\msp{page 96}

Comme toute
\[
\varphi:\tau\longrightarrow\tau'
\]
induit
\[
\varphi:I_\tau\longrightarrow I_{\tau'},
\]
transformant flèches rétrogrades en flèches identiques ou rétrogrades, i.e.
\[
\varphi(\Sigma_\tau)\subseteq\Sigma_{\tau'},
\]
\printedpage{62}
% head: A. Grothendieck : Dérivateurs, Ch. VII, Édition des Universités de Montpellier II et Paris VII
% foot: Version 22 décembre 2001
on trouve le premier des deux foncteurs $(*)$ annoncés page 94, qui est associé par transposition au foncteur $\gamma$
\[
\begin{tikzcd}[column sep=large]
\uHom(\tau,\tau')\times\ \underbrace{\uCh_{\Sigma,\tau'}(X)}_{\simeq\,\uHom_{\Sigma_{\tau'},\Sigma}(I_{\tau'},X)} \arrow[r, "\gamma"] \arrow[d, "\text{can.}"'] & \uCh_{\Sigma,\tau}(X) \arrow[d, equal]\\
\uHom_{\Sigma_\tau,\Sigma_{\tau'}}(I_\tau,I_{\tau'})\times\uHom_{\Sigma_{\tau'},\Sigma}(I_{\tau'},X) \arrow[r, dotted, "\gamma_0"] & \uHom_{\Sigma_\tau,\Sigma}(I_\tau,X)
\end{tikzcd}
\tag*{$(*)$}
\]
déduit du foncteur composition $\gamma_0$ dans ce diagramme : $\gamma=\gamma_0\circ\mathrm{can}.$ D'autre part, ce foncteur de composition $\gamma$ induit visiblement un foncteur sur les sous-catégories de chemins à source et but fixés
\[
\uHom(\tau,\tau')\times\uCh_{\Sigma,\tau'}(X;a,b)\longrightarrow\uCh_{\Sigma,\tau}(X;a,b),
\]
d'où la deuxième des deux flèches annoncées dans $(*)$ p. 94.

Cela montre donc que si
\[
\varphi_1,\varphi_2:\tau\rightrightarrows\tau'
\]
sont deux morphismes de $\tau$ dans $\tau'$ homotopes dans $\uHom(\tau,\tau')$, alors les foncteurs induits entre catégories de $\Sigma$-chemins ou entre catégories de $\Sigma$-chemins à extrémités précisés : % [sic] précisés -> précisées, agreement with extrémités

\msp{page 97}
\[
\begin{aligned}
\varphi^*_{1\,X},\varphi^*_{2\,X} : && \uCh_{\Sigma,\tau'}(X) &\quad\rightrightarrows\quad\uCh_{\Sigma,\tau}(X),\\
(\varphi_1)^*_{a,b},(\varphi_2)^*_{a,b} : && \uCh_{\Sigma,\tau'}(X;a,b) &\quad\rightrightarrows\quad\uCh_{\Sigma,\tau}(X;a,b)
\end{aligned}
\]
sont également homotopes.

Or le lemme d'homotopie nous assure que l'hypothèse $\varphi_1\sim\varphi_2$ est \textit{toujours} satisfaite. Donc :

\textbf{Théorème d'homotopie sur les opérations universelles sur les chemins.} \textit{Soient $X$, $\Sigma$, $a$, $b$ comme dessus, et $\tau,\tau'\in\mathrm{Ty}$ telles que $\mathrm{Hom}(\tau,\tau')\neq\emptyset$. Alors les foncteurs $\varphi^*_X$, $\varphi^*_{X;a,b}$ induits par un $\varphi\in\mathrm{Hom}(\tau,\tau')$}
\[
\left\{
\begin{aligned}
\varphi^*_X : && \uCh_{\Sigma,\tau'}(X) &\quad\longrightarrow\quad\uCh_{\Sigma,\tau}(X)\\
\varphi^*_{X;a,b} : && \uCh_{\Sigma,\tau'}(X;a,b) &\quad\longrightarrow\quad\uCh_{\Sigma,\tau}(X;a,b)
\end{aligned}
\right.
\]
\textit{ne dépendent pas de $\varphi$,} à homotopie près.

On peut préciser cet énoncé, de telle façon que l'énoncé pour les $\varphi^*_{X;a,b}$ apparaisse, comme il devrait, comme corollaire pour celui sur $\varphi^*_X$. Pour ceci, on va regarder les catégories de chemins dans $X$ comme catégories \textit{au dessus de $X$,} par les foncteurs `source et but' :
\[
\begin{tikzcd}[column sep=small]
\uCh_{\Sigma,\tau}(X) \arrow[dr, "(\sigma_\tau{,}\beta_\tau)"'] & & \uCh_{\Sigma,\tau'}(X) \arrow[dl, "(\sigma_{\tau'}{,}\beta_{\tau'})"]\\
 & X\times X &
\end{tikzcd}
\]

\newpage
\printedpage{63}
% head: A. Grothendieck : Dérivateurs, Ch. VII, Édition des Universités de Montpellier II et Paris VII
% foot: Version 22 décembre 2001

\msp{page 98}

Alors le foncteur $(*)$ (p. 94) se factorise par la sous-catégorie $\uHom_{\mathrm{Cat}/X\times X}(\cdot,\cdot)$ de $\uHom_{\mathrm{Cat}}(\cdot,\cdot)$ formée des foncteurs qui commutent aux foncteurs de projection dans $X\times X$, et les homomorphismes entre tels foncteurs qui s'effectuent par des flèches `verticales' de la catégorie but (i.e. par des flèches au dessus de flèches identiques de la catégorie base $X\times X$) :
\[
\uHom(\tau,\tau')\longrightarrow\uHom_{\mathrm{Cat}/X\times X}(\uCh_{\Sigma,\tau'}(X),\uCh_{\Sigma,\tau}(X)).
\]

Donc dans l'énoncé du théorème d'homotopie, on peut préciser que l'homotopie se fait dans cette sous-catégorie, i.e. s'effectue par une suite d'homotopies simples qui sont des morphismes de foncteurs entre catégories au dessus de $X\times X$. Cela implique alors l'énoncé d'homotopie sur les $\varphi^*_{X;a,b}$.

% ---- from Der-ChVII-456_p066-070.tex
\leavevmode\printedpage{64}
% head: A. Grothendieck : Dérivateurs, Ch. VII, Édition des Universités de Montpellier II et Paris VII
\msp{page 99}

\section*{10\quad Question de comparaison pour les catégories \textmd{Cat}-fibrées satisfaisant au ‘théorème d'homotopie' (base filtrante)}

Nous allons généraliser maintenant au cas d'une catégorie (Cat)-fibrée sur une autre
\[
\begin{tikzcd}
C \arrow[d, "\pi"'] \\
T,
\end{tikzcd}
\]
en pensant surtout au cas $T = \uTy$, $C = \uCh_{\Sigma}(X;a,b)$. L'hypothèse essentielle est la suivante :

\noindent\begin{tabular}{@{}p{0.3\textwidth}p{0.64\textwidth}@{}}
(Hypothèse d'homotopie) & Si $t,t' \in \Ob T$ sont tels que $\Hom(t,t') \neq \emptyset$, alors les foncteurs $\varphi^{*} : C_{t'} \longrightarrow C_{t}$ associés aux $\varphi \in \Hom(t,t')$ sont homotopes.
\end{tabular}

Ils définissent donc une seule et même flèche dans Cat\,hot. Je désigne par $T_{\mathrm{ord}}$ ou $\To$ l'ensemble $\Ob(T)$, avec comme relation d'ordre celle déduite de la structure de catégorie :
\[
t \leq t' \qquad\Longleftrightarrow\qquad \Hom_{T}(t,t') \neq \emptyset,
\]
de sorte qu'on a un foncteur
\[
T \longrightarrow \To
\]
qui est \textit{universel} pour les foncteurs de $T$ dans des catégories ‘ordonnées' (i.e. où toute flèche est monomorphique \texttt{[non, cf. }$\bullet\rightrightarrows\bullet$\texttt{; plutôt, où toujours }$\#\Hom(a,b) \leq 1$ \texttt{(admettant des isomorphismes non identiques)]}) ou pour les foncteurs $F$ sur $T$ tels que $F$ soit constant sur tout $\Hom(t,t')$. L'hypothèse d'homotopie peut s'exprimer en disant que le foncteur composé
\[
\begin{tikzcd}
T^{o} \arrow[r] & \Cat \arrow[r] & \Cathot \\
t \arrow[r, mapsto] & C_{t} &
\end{tikzcd}
\]
\msp{page 100}
se factorise (évidemment de façon unique) en un foncteur
\[
H_{C} : \To^{o} \longrightarrow \Cathot,
\]
i.e. qu'on a un diagramme commutatif
\printedpage{65}
% head: A. Grothendieck : Dérivateurs, Ch. VII, Édition des Universités de Montpellier II et Paris VII
\[
\begin{tikzcd}
T^{o} \arrow[r, "t\,\longmapsto\, C_{t}"] \arrow[d] & \Cat \arrow[d] \\
\To^{o} \arrow[r, "H_{C}"] & \Cathot.
\end{tikzcd}
\]

\dgrblock{dgrb333}{The preorder remembers existence, not multiplicity}
\begin{dcomment}{exposition}{p99}{The preorder remembers existence, not multiplicity}
Grothendieck's purpose is to retain the common value of all operations between two fibres while discarding their multiplicity. Let $q:T\to P$ be the identity on objects, with
\[
 P(t,t')=\begin{cases}\{*\},&T(t,t')\ne\varnothing,\\
 \varnothing,&T(t,t')=\varnothing.\end{cases}
\]
The category $P$ is a preorder; distinct objects may be mutually comparable. Its universal property is among thin categories, whose hom-sets have at most one element. Having every arrow monic is a different condition: the category with two parallel arrows already separates the two notions, as the editors' correction observes.

Let $F:T^{\mathrm{op}}\to\Cat$ describe the fibration. In this discussion
\[
 \Cathot(A,B)=\pi_0\uHom(A,B)
\]
means functors modulo finite zigzags of natural transformations. It is not the localization $\Hot$ of categories at weak equivalences of nerves. Natural transformations become equalities in $\Cathot$, and hence also in $\Hot$, but not every weak equivalence of nerves has a functorial homotopy inverse. The distinction matters throughout the comparison problem.

The homotopy hypothesis says exactly that the composite of $F$ with $\Cat\to\Cathot$ factors through $P^{\mathrm{op}}$. If $P$ is nonempty and any two objects have a common lower bound, $P^{\mathrm{op}}$ is filtered and the resulting diagram defines an ind-object $H_C$. Here nonemptiness is part of the usual convention for ``filtrant''. The displayed pairwise condition in the manuscript must be read with that convention; otherwise the empty base would satisfy it vacuously and would have no cofinal sequence of the proposed form.

One must keep this modest hypothesis distinct from $T$ being cofiltered. A cofiltered category must also admit a common precomposition equalizing any pair of parallel arrows. Nothing of that kind is assumed here. The usual definition and its role in finite-limit arguments are recorded in \citet[Categories, Definition~4.19.1 and Lemma~4.19.2]{StacksProject}. For $\Ty$, concatenation supplies common lower bounds because either factor can be collapsed to the empty type. It does not identify the different ways to perform a collapse. The homotopy theorem, rather than strict equalization, is what allows those different ways to be forgotten in $H_C$.
\end{dcomment}


Nous allons supposer de plus :

\noindent\begin{tabular}{@{}p{0.36\textwidth}p{0.58\textwidth}@{}}
(Hypothèse de filtration) & L'ensemble ordonné $\To$ associé à $T$ est \textit{filtrant} décroissant (filtrant à gauche), i.e. si $t,t' \in \Ob T$, il existe $t'' \in \Ob T$ et des flèches $t'' \longrightarrow t$, $t'' \longrightarrow t'$.
\end{tabular}

Cette hypothèse est bien satisfaite dans $\uTy$, car pour $\tau,\tau' \in \Ob \uTy = \mathrm{Ty}$, il suffit de prendre $\tau\tau'$, qui s'envoie dans $\tau$ par $\mathrm{id}_{\tau} * p_{\tau'}$, dans $\tau'$ par $p_{\tau} * \mathrm{id}_{\tau'}$, où $p_{\tau} : \tau \longrightarrow 1$ est le morphisme de $\tau$ dans l'objet final.

Ainsi, $H_{C}$ définit un ind-objet de Cat\,hot, que nous allons désigner par $\underline{\underline{C}}$. Si on suppose de plus :

\noindent\begin{tabular}{@{}p{0.36\textwidth}p{0.58\textwidth}@{}}
(Hypothèse de dénombrabilité, toujours satisfaite si $\mathrm{card}\,\To \leq \aleph_{0}$) & $\To$ admet un sous-ensemble \textit{dénombrable} cofinal à gauche.
\end{tabular}

alors cet ind-objet de Cat\,hot peut être défini, de beaucoup de façons, en termes
\msp{page 101}
d'un système inductif dans Cat (pas seulement dans Cat\,hot)
\[
(C_{n})_{n\in\mathbf{N}},
\]
d'où on récupère l'ind-objet $\underline{\underline{C}}$ dans Cat\,hot en passant au système inductif correspondant dans Cat\,hot, puis à l'ind-objet associé.

Pour construire le système inductif $C_{n}$, il suffit de prendre une chaîne infinie
\[
t_{0} \longleftarrow t_{1} \longleftarrow t_{2} \longleftarrow \;\cdots\; \longleftarrow t_{n} \longleftarrow t_{n+1} \longleftarrow \;\cdots
\]
de morphismes dans $T$, telle que l'ensemble des $t_{i}$ soit cofinal à gauche dans $\To$. On va supposer de plus $C$ \textit{strictement fibrée} sur $T$ (i.e. provenant d'un véritable foncteur $T^{o} \longrightarrow \Cat$, i.e. on peut choisir les foncteurs images inverses
\[
\varphi^{*} : C_{t'} \longrightarrow C_{t}
\]
de façon transitive stricte, pas seulement transitive à isomorphisme près) - cas auquel on peut toujours se ramener d'ailleurs par l'astuce de \textsc{Giraud}. Cette hypothèse est d'ailleurs vérifiée dans le cas $\uCh_{\Sigma}(X;a,b)$ sur $\uTy$. Ainsi, les
\[
(C_{t_{i}})_{i\in\mathbf{N}}
\]
forment un système inductif de catégories indexées par $\mathbf{N}$, le morphisme de transition
\[
C_{n} \longrightarrow C_{n'} \qquad \text{i.e. } C_{t_{n}} \longrightarrow C_{t_{n'}}
\]
\printedpage{66}
% head: A. Grothendieck : Dérivateurs, Ch. VII, Édition des Universités de Montpellier II et Paris VII
\msp{page 102}
étant le foncteur image inverse $\varphi^{*}_{n,n'}$, où
\[
\varphi_{n,n'} : t_{n'} \longrightarrow t_{n}
\]
est le composé des $t_{n'} \longrightarrow t_{n'-1} \longrightarrow \;\cdots\; \longrightarrow t_{n}$. (Si on n'avait pas transitivité stricte, on définirait $\varphi^{*}_{n,n'}$ comme le composé des foncteurs image inverse par des flèches de transition successives $t_{n'} \longrightarrow t_{n'-1} \longrightarrow \;\cdots\; \longrightarrow t_{n+1} \longrightarrow t_{n}$, ce ne serait $\varphi^{*}_{n,n'}$ qu'à isomorphisme canonique près.)

\dgrblock{dgrb334}{Making a sequence strict}
\begin{dcomment}{context}{p101}{Making a sequence strict}
Grothendieck wants an actual category as a sequential colimit, so he temporarily makes the reindexing functors strictly composable. This is a familiar issue from the theory of fibred categories: a cleavage supplies canonical comparison isomorphisms, whereas a splitting makes those comparisons identities. The distinction is already developed in \citet[Expos\'e~VI, \S\S8--9]{SGA1}. The replacement invoked here is an equivalence of fibred categories; it is not a claim that a chosen cleavage on the original total category can always be made split without changing that category.

There is a concrete construction behind the reference to Giraud. For $t\in T$, let $\widetilde F(t)$ be the category of cartesian functors $T/t\to C$ over $T$, with vertical natural transformations. Precomposition along $T/s\to T/t$, induced by $s\to t$, is strictly associative. Evaluation at the terminal object $\id_t$ gives an equivalence $\widetilde F(t)\simeq C_t$: a cartesian functor is reconstructed from its value there by cartesian lifts, and any morphism between those values extends uniquely to a vertical transformation. Thus $\widetilde F$ is a strict replacement of the original pseudofunctor. The right splitting and its equivalence property are treated explicitly in \citet[Theorem~5.4 and Remark~5.5]{EmmeneggerMesitiRosoliniStreicher2024}; this gives a precise modern reference for the ``astuce de Giraud''.

Countability has a different role. Enumerate a cofinal family $s_0,s_1,\ldots$ in $P$. Choose $t_0$ below $s_0$, then choose $t_{n+1}$ below both $t_n$ and $s_{n+1}$, together with an arrow $t_{n+1}\to t_n$. The resulting chain is cofinal in the preorder. Pullback produces
\[
 C_{t_0}\longrightarrow C_{t_1}\longrightarrow C_{t_2}\longrightarrow\cdots.
\]
Even without replacing the entire fibration, this particular sequence can be made strict simply by defining its long transition maps as composites of its adjacent ones. There are no competing paths in the freely generated chain to compare. The arrow $C_{t_n}\to C_{t_{n-1}}$ printed on manuscript page~102 runs in the wrong direction; the intended transition is $C_{t_{n-1}}\to C_{t_n}$.
\end{dcomment}


Soit
\[
C_{\infty} = \varinjlim_{n} C_{n} = \varinjlim_{n} C_{t_{n}}
\]
la catégorie $\varinjlim$ . (Elle sera surtout sympathique si on a pu choisir le système $(t_{n})_{n\in\mathbf{N}}$ de telle façon que les foncteurs de transition
\[
C_{t_{n}} \longrightarrow C_{t_{n-1}}
\]
soient pleinement fidèles, et peut-être même \texttt{[des]} monomorphismes, de telle façon que les $C_{n}$ s'identifient à des sous-catégories pleines de leur limite inductive $C_{\infty}$.) Je voudrais comprendre les relations précises entre les trois genres d'objets suivants :

\begin{itemize}[leftmargin=2.5em]
\item[\circled{a}] $C$ elle-même (comme objet dans Cat, pas de $\Cat/T$), et surtout les objets de Cat\,hot et de Hot qui lui correspondent.

\item[\circled{b}] \msp{page 103}$C_{\infty}$, et les objets de Cat\,hot et de Hot qui lui correspondent. Dans quelle mesure ces deniers % [sic] deniers -> derniers, missing letter
sont-ils indépendant % [sic] indépendant -> indépendants, agreement with the plural subject
du choix du système projectif cofinal à gauche $(t_{n})_{n\in\mathbf{N}}$ dans $T$? Quel est le rôle de $C_{\infty}$? Est-elle canoniquement isomorphe à $C$ dans Cat\,hot, ou dans Hot? (Auquel cas la question de son ‘rôle' me semblerait résolue!)

\item[\circled{c}] Le ind-objet % [sic] Le ind-objet -> L'ind-objet, missing elision
$H_{C}$ dans Cat\,hot, et l'ind-objet de Hot qu'il définit.

Je note que l'ind-objet $H_{C}$, et aussi les catégories $C_{\infty}$ construites comme dessus, ne changent pas quand on remplace $T$ par une sous-catégorie $T'$ cofinale à gauche, ou aussi par une catégorie de la forme $T/t_{0}$ ($t_{0} \in \Ob T$), qui donne lieu à l'ensemble ordonnée % [sic] ordonnée -> ordonné, agreement with "ensemble"
associé
\[
(T/t)_{\mathrm{o}} = \To/t = \{t \in \To \mid t \leq t_{0}\}
\]
\texttt{[plutôt} $(T/t_{0})_{\mathrm{o}} = \To/t_{0} = \{t \in \To \mid t \leq t_{0}\}$\texttt{]}. Si la relation entre les objets dans a), et ceux dans b) et dans c), est très étroite, on en concluerait % [sic] concluerait -> conclurait, spelling of the conditional
que le ‘type' de $C$ dans Cat\,hot, ou dans Hot suivant le cas, ne change pas en remplaçant $C$ par
\printedpage{67}
% head: A. Grothendieck : Dérivateurs, Ch. VII, Édition des Universités de Montpellier II et Paris VII
\msp{page 104}
$C \times_{T} T'$ ou par $C/t_{0} \simeq C \times_{T} T/t_{0}$. Bien mieux, elle ne changerait pas en remplaçant $T$ par $(\mathbf{N})^{o}$, la catégorie (ordonnée) opposée à $\mathbf{N}$ (ayant donc un unique objet final), et la catégorie fibrée $C$ sur $T$ par la catégorie $C_{\underline{t}}$ induite par $C$ sur $(\mathbf{N})^{o}$, via le foncteur
\[
\underline{t} : (\mathbf{N})^{o} \longrightarrow T
\]
défini par le système projectif $\underline{t} = (t_{n})_{n\in\mathbf{N}}$ dans $T$. Cela signifierait donc que pour comprendre le type homotopique (au sens strict, i.e. dans Cat\,hot, ou large, i.e. dans Hot) de $C$, on peut le ‘nettoyer' de façon draconienne, la base $T$ étant vue comme une sorte de gangue un peu vaseuse, ou du moins pléthorique.

\dgrblock{dgrb335}{A slice does not have the same preorder}
\begin{dcomment}{development}{p103}{A slice does not have the same preorder}
Grothendieck proposes restricting the construction to types below a chosen type. This is sound for a lower ideal of the preorder, but a slice category retains the actual arrow to the chosen object. The equality of preorders printed here therefore needs a correction that affects the proposed argument.

\textbf{A slice obstruction.}\ours\quad There is a finite category $T$ whose associated preorder is cofiltered, while the associated preorder of $T/t_0$ is not cofiltered.

\emph{Proof.} Take two objects $s,t_0$, their identities, and exactly two further arrows $f,g:s\rightrightarrows t_0$. The preorder $P$ has $s<t_0$, so $s$ is a lower bound for every pair. The slice $T/t_0$ has three objects
\[
 (s,f),\qquad (s,g),\qquad(t_0,\id_{t_0}),
\]
with one arrow from each of the first two objects to the last. An arrow between the first two would have to be $\id_s$ and would force $f=g$. There is no fourth object that could map to both. Its preorder is consequently not cofiltered, although the lower ideal $P/t_0=P$ is. This proves the assertion.\hfill$\square$

The example also prevents interpreting the displayed equality as an equivalence after removing duplicate isomorphic objects: the two incomparable objects remain distinct in the skeletal preorder. It does not dispute the connectedness of the slice, which has a terminal object, nor the separate construction of $H_C$ over $P^{\mathrm{op}}$.

There are two precise repairs. If only the ind-object is wanted, use the full subcategory $T_{\le t_0}\subseteq T$ on objects admitting some arrow to $t_0$. Its preorder is exactly the lower ideal of $P$. Its inclusion is cofinal on associated preorders: given $t$, choose a common lower bound of $t$ and $t_0$. It is cofiltered as a preorder for the same reason. Restriction therefore leaves $H_C$ unchanged. The actual slice $T/t_0$ can instead be used when $T$ is cofiltered in the categorical sense. Given $u:a\to t_0$ and $v:b\to t_0$, first choose $r\to a,b$, then precompose once more to equalize the resulting two arrows $r\rightrightarrows t_0$. This supplies a common predecessor in the slice. Parallel arrows in the slice can likewise be equalized in $T$, and the resulting equality still lies over $t_0$. Thus the slice is then cofiltered.

We checked the standard filtering results in \citet[Categories, \S4.19]{StacksProject} and the homotopy-colimit comparison discussed below; we did not find this particular correction to manuscript page~103 stated there. The example and the two repairs are recorded as a local development of the argument, without a claim of priority.
\end{dcomment}


\item[\circled{d}] De toutes façons, si $T'$ est une autre catégorie filtrante à gauche, et
\[
\alpha : T' \longrightarrow T
\]
un foncteur cofinal à gauche, de sorte que $C$ sur $T$, et $C' = C \times_{T} T'$ sur $T'$, définissent essentiellement le ‘même' objet de Ind\,Cat\,hot, et que la catégorie $C'_{\infty}$ pour $C'$ sur $T'$ fait partie des catégories $C_{\infty}$ pour $C$ sur $T$ - on peut se demander
\msp{page 105}
si le foncteur canonique
\[
C' \;\longrightarrow\; C
\]
est également toujours un isomorphisme dans Cat\,hot, ou sinon, du moins dans Hot.
\end{itemize}

Un intérêt particulier s'attache à $\pi_{0}(C)$, et de plus à sa variance pour $C$ sur $T$ variable de la façon suivante. Si $S$ est dans Cat, on regarde les catégories
\[
\uHom_{\Cat}(S,C_{t}) \qquad (t \in T),
\]
on trouve un foncteur
\[
t \longmapsto \uHom_{\Cat}(S,C_{t}), \qquad T^{o} \longrightarrow (\Cat),
\]
d'où une catégorie fibrée sur $T$, que je vais désigner par
\[
\sHom_{T}(S_{T},C)
\]
(Sauf erreur, elle s'interprète comme la catégorie fibrée sur $T$ correspondante au champ sur $\mathrm{Top}(T) = T^{\wedge}$ du morphisme de $T$-champs de $S_{T}$ dans $C$, où $S_{T} = S \times T$ (considérée comme catégorie fibrée sur $T$ à fibre essentiellement constante $S$) dans $C$. Ici j'identifie catégorie fibrée sur $T$, et champs sur le topos $T^{\wedge}$.) Il est immédiat que la propriété d'homotopie (p. 99) est stable par passage de $C$ à $\sHom_{T}(S_{T},C)$. On s'intéresse à la variation de
\[
\pi_{0}(\sHom_{T}(S_{T},C))
\]
\printedpage{68}
% head: A. Grothendieck : Dérivateurs, Ch. VII, Édition des Universités de Montpellier II et Paris VII
\msp{page 106}
comme contrafoncteur en $S$,
\[
H : S \longmapsto \pi_{0}\sHom_{T}(S_{T},C) \qquad (\Cat)^{o} \longrightarrow (\mathrm{Ens}),
\]
et à une éventuelle expression de ce contrafoncteur $H$ en termes de l'ind-objet $H_{C}$ de Cat\,hot, par la formule hypothétique
\[
H(S) \simeq \Hom_{\mathrm{Ind}(\Cathot)}(\gamma(S),H_{C}),
\]
où
\[
\gamma : \Cat \longrightarrow \Cathot
\]
est le foncteur canonique. Si on avait un tel isomorphisme canonique fonctoriel en $S$, cela donnerait une description (à isomorphisme unique près) de l'ind-objet $H_{C}$ en termes du foncteur précédent $H$ sur Cat.

Il me vient un doute ici, en vue de l'exemple $T = \uTy$, $C = \uCh_{\Sigma}(X;a,b)$. Les trois types d'objets a) b) c) ont le même $\pi_{0}$, et il est sûr que ‘c'est le bon'. D'autre part, pour une catégorie paramètre variable $S$, il est sûr aussi que le bon foncteur en $S$ est
\[
S \longmapsto \pi_{0}(\uHom_{T}(S_{T},C)) \simeq \Hom_{\mathrm{Ind}\,\Cathot}(S,H_{C}),
\]
isomorphe aussi, \textit{si $S$ est finie,} à
\[
\Hom_{\Cathot}(S,C_{\infty}),
\]
ce qui montre que, du moins pour des arguments $S \in (\Cat)$ qui sont des catégories \textit{finies}, les objets
\msp{page 107}
$C_{\infty}$ et $H_{C}$ sont ‘bons', il nous permettent % [sic] il nous permettent -> ils nous permettent, agreement with the plural subject
de décrire un foncteur important qui domine la situation envisagée. Moyennant une hypothèse très forte, mais satisfaite dans les cas importants, on peut dans ces deux formules remplacer Cat\,hot par Hot.

\dgrblock{dgrb336}{The functor represented by the ind-object}
\begin{dcomment}{exposition}{p106}{The functor represented by the ind-object}
Grothendieck's parameter category $S$ restores information that the connected components of a single fibre cannot carry. The proposed formula for $H(S)$ is valid under the hypotheses already imposed, and its proof explains both the ind-object and the restriction to finite $S$ in the next formula.

For any strict diagram $A:T^{\mathrm{op}}\to\Cat$, the connected components of its fibred total category are
\[
 \pi_0\left(\int_T A\right)
 \cong\mathop{\mathrm{colim}}_{T^{\mathrm{op}}}\pi_0 A(t).
\]
Indeed, the total category is generated by vertical arrows and cartesian transport arrows. Vertical arrows impose exactly the component relation inside each $A(t)$; a cartesian arrow imposes exactly the relation identifying a class with its restriction along an arrow of $T$. Those are the generators and relations defining the colimit of sets, with no further relation. Apply this observation to $A(t)=\uHom(S,C_t)$. A zigzag between two reindexing functors induces a zigzag after postcomposition with functors from $S$, so the diagram of component sets factors through $P^{\mathrm{op}}$. We obtain, naturally in $S$,
\[
 \begin{split}
 H(S)&=\pi_0\sHom_T(S_T,C)\\
 &\cong\mathop{\mathrm{colim}}_{P^{\mathrm{op}}}
       \Cathot(S,C_t)
 =\Hom_{\operatorname{Ind}(\Cathot)}(S,H_C).
 \end{split}
\]
The last equality is the defining formula for maps from a constant object into an ind-object. The usual fully faithful embedding of ind-objects into set-valued presheaves then gives the asserted uniqueness of $H_C$. Setting $S=e$ proves the equality of the three component sets in the text.

To compare with $C_\infty$, suppose $S$ has finitely many objects and arrows. A functor $S\to\mathop{\mathrm{colim}}_n C_n$ involves finitely many objects, arrows, and composition equations. Represent its objects and arrows at a common stage, then move to a later stage where all the finitely many required equations hold. It therefore descends to some $C_n$. A natural transformation between two such functors also consists of finitely many arrows and naturality equations and descends in the same way. Finally, if two descended functors lie in the same component, a witnessing homotopy is a \emph{finite} zigzag of such transformations; the whole zigzag descends to one later stage. This proves the bijection
\[
 \mathop{\mathrm{colim}}_n\Cathot(S,C_n)
 \xrightarrow{\ \sim\ }\Cathot(S,C_\infty).
\]
It is an instance of the finite-presentation argument underlying filtered colimits, not a fibrancy assertion.

The finiteness cannot simply be omitted. Let $C_n$ be the discrete category on $\{0,\ldots,n\}$ and let $S$ be the discrete category on $\mathbf N$. Then $C_\infty$ is discrete on $\mathbf N$. The left side consists of functions $\mathbf N\to\mathbf N$ with bounded image; the right side consists of all functions. The identity function lies only in the latter. The ind-object retains the information that a whole $S$-family must occur at one stage, while the ordinary colimit allows the stage to vary with the individual element of $S$.
\end{dcomment}


Par contre, sauf pour le cas où $S$ est réduit à un point et où $\Hom_{\Cathot}$ ou $\Hom_{\Hot}$ de $S$ dans $C$ n'est autre que $\pi_{0}(C)$ ($\simeq \pi_{0}(C_{\infty}) \simeq \pi_{0}(H_{C})$), donc effectivement ‘bon', je ne sais pas si
\[
\underbrace{\Hom_{\Cathot}(S,C)}_{=\,\pi_{0}\uHom(S,C)} \qquad \text{ou} \qquad \Hom_{\Hot}(S,C)
\]
est lié de façon raisonnable aux ensembles précédents (déterminés par $H_{C}$, resp. par $C_{\infty}$), p. ex. s'ils lui sont isomorphes de façon naturelle, et s'ils ont une signification géométrico-topologique raisonnable. Je commence à suspecter que \textit{non} - donc qu'on aurait ici un cas où une relation de type très général concernant un $\pi_{0}$, ici donc
\[
\pi_{0}(C) \simeq \pi_{0}(C_{\infty}) \simeq \pi_{0}(H_{C}),
\]
ne se prolonge pas en un isomorphisme sur les $\pi_{i}$ supérieurs, et en une propriété de Hot-équivalence.

% ---- from Der-ChVII-456_p071-075.tex
\printedpage{69}
% head: A. Grothendieck : Dérivateurs, Ch. VII, Édition des Universités de Montpellier II et Paris VII

\msp{page 108}
Je ne trouve d'ailleurs aucune flèche directe entre $C$ et l'un des objets $C_{\infty}$, $H_{C}$; tout ce que je vois dans ce sens, c'est
\[
\begin{tikzcd}[column sep=large, row sep=huge]
\CNo \arrow[r, "{\substack{\text{foncteur de}\\ \text{localisation}}}"] \arrow[d] & C_{\infty} \\
C, &
\end{tikzcd}
\]
où la flèche verticale est la flèche de projection de $\CNo = C \times_{T} \mathbf{N}^{o}$ (avec $\mathbf{N}^{o} \xrightarrow{\ t\ } T$). Ceci n'établit une relation topologique étroite entre $C$ (en tant qu'objet de Cat\,hot ou de Hot simplement) et $C_{\infty}$, que si les deux flèches précédentes sont tout au moins des Hot-équivalences. Et le cas-type (problématique) où la question de la relation entre $C$ et $C_{\infty}$ se pose, est donc celui où $T = \mathbf{N}^{o}$, i.e. celui d'un système inductif quelconque de (petites) catégories, indexé par $\mathbf{N}$, et interprété comme défini par une catégorie fibrée $C = \CNo$ sur $\mathbf{N}^{o}$. Dans ce cas particulier, on a du moins une flèche canonique
\[
C \longrightarrow C_{\infty}
\]
sur laquelle on peut se poser la question si c'est
\msp{page 109}
une Hot-équivalence. Sinon, cela prouvera que dans le cas général, le Hot-type de $C$ (et a fortiori sa classe d'isomorphie dans Cat\,hot) n'a pas de sens géométrique intéressant. C'est à dire que $C$ est intéressante uniquement \textit{en tant que catégorie fibrée sur $T$}, permettant notamment de construire l'ind-objet $H_{T}$ de Cat\,hot ou de Hot, mais pas en tant que `catégorie nue'.

\dgrblock{dgrb337}{A finite category need not be a finite homotopy probe}
\begin{dcomment}{examples}{p107}{A finite category need not be a finite homotopy probe}
Grothendieck explicitly reserves the passage from $\Cathot$ to $\Hot$ for an additional hypothesis. The reservation is necessary. The finite-presentation proof above counts objects, arrows, and equations in a category. Compactness in the homotopy category of spaces concerns the whole nerve, and a finite category can have nondegenerate simplices in arbitrarily high degrees.

Let $S=\mathbf B(\mathbf Z/2)$, which has one object and two arrows, and put $K=NS$, $K_n=\operatorname{sk}_n K$. Let $D_n=\Delta/K_n$ be the category of simplices of $K_n$. Every simplex of $K$ belongs to one skeleton, so
\[
 D_\infty=\mathop{\mathrm{colim}}_nD_n=\Delta/K.
\]
The simplex-category resolution gives $|ND_n|\simeq|K_n|$ and $|ND_\infty|\simeq|K|$, naturally for the inclusions. The precise last-vertex statement for an arbitrary simplicial set is \citet[Corollary~11.9.11.5]{Kerodon2026Simplex}. Thus the filtered colimit has the homotopy type of $B(\mathbf Z/2)$, while each stage has an $n$-dimensional realization.

The identity of $B(\mathbf Z/2)$ cannot factor in $\Hot$ through any $|K_n|$. Such a factorization would make its identity map on cohomology factor through a group vanishing in degrees greater than $n$. But
\[
 H^*(B(\mathbf Z/2);\mathbf F_2)\cong\mathbf F_2[u],\qquad |u|=1.
\]
For completeness, the normalized group cochains with trivial $\mathbf F_2$ coefficients have one basis element in each degree: a normalized cochain is determined by its value on the string consisting entirely of the nonidentity group element. Its differential is zero, since the two end terms cancel and every interior product gives the identity, on which a normalized cochain vanishes. The degree-one generator and its cup powers take value one on these strings. This proves the displayed calculation and the obstruction.

Consequently the canonical map
\[
 \mathop{\mathrm{colim}}_n\Hom_{\Hot}(S,D_n)
 \longrightarrow\Hom_{\Hot}(S,D_\infty)
\]
is not surjective, although $S$ is finite as a category. There is no contradiction with the preceding $\Cathot$ formula. Finite CW homotopy types are compact probes for filtered homotopy colimits, but the nerve of this finite group is not one. This is the point at which the two meanings of ``finite'' must be separated when reading the notation $\Hotz$ later on these pages.
\end{dcomment}


\dgrblock{dgrb338}{The sequential total category is a homotopy colimit}
\begin{dcomment}{later}{p108}{The sequential total category is a homotopy colimit}
Grothendieck has reduced one part of the comparison to a single sequence. That part has a uniform answer: the canonical functor from its fibred total category to its ordinary filtered colimit is a weak equivalence of nerves. No full-faithfulness assumption on the transition functors is needed.

Write $E=\int_{\mathbf N^{\mathrm{op}}} C_n$ and $D=\mathop{\mathrm{colim}}_n C_n$. With the variance of this manuscript, an arrow $(m,y)\to(n,x)$ over $m\to n$, where $m\ge n$, is an arrow $y\to F_{n,m}x$ in $C_m$. The canonical functor $q:E\to D$ sends it to the corresponding arrow in $D$. Cartesian transport arrows become identities. The localization description in the text is understood in the usual categorical sense: after inverting those arrows, $E$ becomes \emph{equivalent} to $D$. Literal equality of its object set with that of $D$ is neither needed nor generally true.

The homotopy-colimit theorem of \citet[Theorem~1.2]{Thomason1979}, in its dual form for a fibration, identifies $|NE|$ with $\operatorname{hocolim}_n|NC_n|$. The convention for a contravariant diagram and the comparison map are displayed explicitly in \citet[\S3, equation~(9)]{Cegarra2020}. Two elementary filtered-colimit facts complete the argument. First,
\[
 N D\cong\mathop{\mathrm{colim}}_n NC_n:
\]
an $r$-simplex is a finite string of composable arrows, and finitely many representatives and equalities can all be realized at one stage. Second, for a filtered diagram of simplicial sets $K_n$, the augmentation
\[
 \operatorname{hocolim}_n K_n\longrightarrow
 \mathop{\mathrm{colim}}_n K_n
\]
is a weak equivalence. One way to see the second fact is to use the simplicial replacement. In each fixed simplicial degree, the fibre over a simplex of the colimit is the category of its representatives at the different stages. Any two representatives eventually agree, and finitely many such agreements hold simultaneously at a later stage. This category is nonempty and filtered, hence its nerve is contractible: every finite diagram admits a cone, which contracts each finite simplicial sphere. Thus the augmentation is a weak equivalence in each row of the bisimplicial replacement, and diagonal realization preserves these rowwise weak equivalences. The composite is the map induced by $q$.

This identifies the genuine issue in the span on manuscript page~108. The horizontal arrow is always a weak equivalence. Whether the vertical arrow is one depends on how the chosen sequence sits inside $T$, including the homotopy carried by its parallel arrows. Passing to the preorder has deliberately forgotten that information.
\end{dcomment}


\dgrblock{dgrb339}{The total category and the telescope}
\begin{dcomment}{exposition}{p108:sur laquelle on peut se poser la question}{The total category and the telescope}
Grothendieck wants a path category that carries composition and at the same time has the intended homotopy type. The two objects in this opening comparison solve different parts of that problem. The total category remembers how the types of paths vary; the filtered colimit forgets the stages at which paths were represented. We write $BC=|NC|$, and call a functor a weak equivalence when its nerve is a weak homotopy equivalence. This is the meaning of an isomorphism in $\Hot$. It is stronger than a bijection on connected components and different from an isomorphism in $\mathrm{Cat\,hot}$, where morphisms are connected components of ordinary functor categories.

Here is a proof of the assertion announced on the following page. Let $F_0\to F_1\to\cdots$ be a strict sequence of small categories, let $E\to\mathbf N^o$ be its fibred Grothendieck construction, and let $L=\mathop{\rm colim}_nF_n$. The canonical functor $q:E\to L$ is a weak equivalence. For clarity, a morphism $(n,x)\to(m,y)$ in $E$, with $n\geq m$, is a morphism $x\to F_{m,n}y$ in $F_n$.

To prove the assertion, fix $z\in L$ and consider the comma category $q/z$. Any finite collection of its objects, arrows, and commuting triangles can be represented at one sufficiently large stage $r$: filtered colimits of categories are computed from finite diagrams and their eventual relations. We may also represent $z$ at stage $r$. Transporting the finitely many objects to that stage gives a diagram in $F_r/z_r$. Its projection to $q/z$ has a natural transformation to the original diagram, supplied by the cartesian arrows in $E$. It has a natural transformation to the constant diagram $(r,z_r,1)$, supplied by the arrows to $z_r$. Thus every finite simplicial diagram in $N(q/z)$ contracts after passing to a stage. The same argument with a chosen vertex proves connectedness and vanishing of all homotopy groups; the comma category is nonempty because $z$ has a representative. Quillen's Theorem A now gives the assertion \citep[\S1, Theorem A]{Quillen1973}. This argument does not require the transition functors to be faithful or injective on objects.

The second leg of the displayed roof is a separate matter. Take a nontrivial group $G$, let $T=BG$ be its one-object category, and take $C=T$ over itself. The preorder obtained by forgetting the distinction between parallel arrows is a point. A sequence constant on the identity arrow is therefore cofinal in that preorder, whereas the comparison is $*\to BG$. Its fundamental group is $1\to G$, so it is not a weak equivalence. This proves exactly the limitation noticed on manuscript page~109: cofinality in the preorder of objects does not account for the automorphisms of the indexing category. It does not predict failure for the particular category of interval types, which has additional structure.
\end{dcomment}


Mais je viens de me convaincre que $\CNo \longrightarrow C_{\infty}$ est bien une Hot-équivalence. Donc la flèche $\CNo \longrightarrow C$ définit une flèche \textit{dans} Hot
\[
C_{\infty} \longrightarrow C
\]
(en composant avec l'inverse dans Hot de la précédente). Est-ce un isomorphisme? Donc pour cette question, on est ramené à la question pour $\CNo \longrightarrow C$. Mais si on prend $C = T$, cela pose déjà la même question pour $\mathbf{N}^{o} \longrightarrow T$, donc la question si $T$ est asphérique, sous la seule hypothèse que $T_{\mathrm{o}}$ soit filtrante à gauche et dénombrable à l'infini (à gauche). Mais bien sûr c'est faux (p. ex. si $T$ n'a qu'un seul objet, il n'est pas pour autant asphérique !).

Cela me convainc donc qu'à défaut
% foot: Version 22 décembre 2001
\printedpage{70}
% head: A. Grothendieck : Dérivateurs, Ch. VII, Édition des Universités de Montpellier II et Paris VII
\msp{page 110}
d'hypothèse supplémentaire draconienne sur $T$, la catégorie $C$ `nue' n'offre pas d'intérêt particulier, à part son $\pi_{0}$ (qui est donné aussi par $C_{\infty}$, et par $H_{C}$). Il n'est pas dit, certes, que \underline{Ty} ne satisfasse pas à ces `conditions draconiennes'. Mais comme $C$ `nue' ne s'insère pas dans un formalisme sympathique (p.ex. on n'a pas d'interprétation géométrique de $\mathrm{Hom}_{\mathrm{Cat\,hot}}(S,C)$, ou de $\mathrm{Hom}_{\mathrm{Hot}}(S,C)$), je crois que par la suite je peux en faire abstraction. Exit a) !

\dgrblock{dgrb340}{What a one-object base can retain}
\begin{dcomment}{examples}{p109}{What a one-object base can retain}
Grothendieck's one-object example can be made completely explicit, and it shows exactly which information is absent from the ind-object. Let $G$ be a nontrivial discrete group and let $T=\mathbf BG$ be the category with one object and endomorphism group $G$. Take $C=T$ with its identity projection. This is a split fibration with terminal fibre and identity reindexing functors. Its homotopy hypothesis holds strictly; its associated preorder is the one-point category, which is cofiltered and has a countable cofinal family.

Choose the constant chain at that object, with identity transition arrows. Every $C_n$ and $C_\infty$ is the terminal category, and $H_C$ is the constant terminal ind-object. The total category remains $\mathbf BG$. Consequently the comparison in $\Hot$ is
\[
 *\longrightarrow BG,
\]
where $BG=|N\mathbf BG|$. Its source has trivial fundamental group and its target has fundamental group $G$. It is a bijection on components and is not a weak equivalence. More generally the constant fibration with fibre $D$ has total category $T\times D$, so it retains the product $BG\times|ND|$ while its sequential colimit is $D$.

There is also a direct answer to the mapping-set question of manuscript page~107. Set $G=\mathbf Z/2$ and $S=\mathbf BG$. Since the fibres of $C\to T$ are terminal, $H(S)$ is a singleton. The identity and the trivial functor $S\to C$ are not homotopic in $\Cathot$: natural transformations between functors of one-object groupoids express conjugacy of the corresponding homomorphisms, and these two homomorphisms are not conjugate. They remain distinct in $\Hot$, because their maps on $\pi_1$ differ. Thus neither of the proposed mapping sets into the bare total category recovers $H(S)$ in this generality.

The example concerns the abstract hypotheses of \S10. It does not decide the special path-category comparison over $\Ty$. That case has additional geometric structure, used later in the chapter. Nor does it make $C$ topologically meaningless: $C$ remembers the base monodromy, which is precisely the datum suppressed by the thin reflection. The warning is against identifying two different invariants.
\end{dcomment}


\begin{flushright}
{\footnotesize 5.11.90}
\end{flushright}

J'ai quand même passé la journée entière hier à essayer de comprendre la relation de $C$ et de son type dans Hot, à celui de $C_{\infty}$. En effet, comme on a finalement une flèche dans Hot
\[
C_{\infty} \longrightarrow C,
\]
déduite du diagramme dans (Hot)
\[
\begin{tikzcd}[column sep=huge, row sep=huge]
\CNo \arrow[r, "\sim"] \arrow[d] & C_{\infty} \\
C, &
\end{tikzcd}
\]
de prouver que $C_{\infty} \longrightarrow C$ est un isomorphisme dans Hot
\msp{page 111}
semble la façon la plus simple, et peut-être unique, de prouver que l'image de $C_{\infty}$ dans Hot ne dépend pas, à isomorphisme unique près, des choix faits pour construire $C_{\infty}$. Sans compter que dans le cas motivant, $C = \Ch_{\Sigma}(X)$ ou $\Ch_{\Sigma}(X;a,b)$, $T = \Ty$, c'est la catégorie des chemins $C$ qui donne lieu à une loi de composition des chemins impeccable, fonctorielle, associative et tout, alors qu'il n'en est rien pour $C_{\infty}$. On ne trouve pas une loi de composition
\[
C_{\infty} \times C_{\infty} \longrightarrow C_{\infty}
\]
dans Cat, ni même dans Cat\,hot ou dans Hot - ou plutôt, si on arrive pourtant à construire une flèche $C_{\infty} \times C_{\infty} \longrightarrow C_{\infty}$ dans Cat, c'est encore laborieusement à coups de choix $\pm$ arbitraires, cran par cran, et sans savoir si ce qu'on obtient a une propriété d'unicité, ne serait ce que dans Hot; on ne trouve une interprétation géométrique tangible de cette composition que dans $\Hotz$ (où on a une véritable loi de groupoïde), ou dans $H_{T}$ en tant qu'objet de Ind\,Cat\,hot
% foot: Version 22 décembre 2001
\printedpage{71}
% head: A. Grothendieck : Dérivateurs, Ch. VII, Édition des Universités de Montpellier II et Paris VII
\msp{page 112}
ou de Ind\,Hot. (Donc la limite inductive dans $\Hotz$ n'est autre que l'élément de $\Hotz$ défini par $C_{\infty}$.)

Mes efforts pour prouver que
\[
C_{\infty} \longrightarrow C
\]
était un isomorphisme dans Hot, ou ne serait-ce qu'un 1-isomorphisme (i.e. un isomorphisme non seulement pour les $\pi_{0}$, ce qui est toujours vrai, mais aussi pour les $\pi_{1}$ de toutes les composantes), ont échoués. J'ai trouvé qu'il suffisait que le foncteur % [sic] ont échoués -> ont échoué, accord fautif du participe
\[
T \longrightarrow T_{\mathrm{o}}
\]
soit ou bien lisse et coasphérique, ou bien propre et asphérique - donc tout est OK si $T = T_{\mathrm{o}}$, i.e. si $T$ est ordonnée. Mais dans le cas $T = \Ty$, le foncteur $T \longrightarrow T_{\mathrm{o}}$ n'est ni propre, ni lisse (pas même 0-propre ou 0-lisse).

Finalement j'en suis arrivé à la conviction que dans le cas qui nous occupe ($T = \Ty$, $C = \Ch_{\Sigma}(X;a,b)$), la flèche $C_{\infty} \longrightarrow C$ n'est en général \textit{pas} un isomorphisme dans Hot. En effet, dans le cas $a = b$, si cela \texttt{[était vrai]},
\msp{page 113}
on aurait une description du Hot-groupe $\Omega(X,a) = \Ch(X,a)$ (\footnote[23]{On fait ici $\Sigma = \mathrm{Fl}(X)$.}) en termes d'une catégorie $C$ munie d'une loi de composition \textit{associative} (et unitaire) \textit{dans} Cat, redonnant la loi de groupe de $\Omega$ en passant de (Cat) à Hot. Prenant le nerf de $C$, on trouverait un monoïde semi-simplicial qui, dans Hot, donnerait la loi de groupe de $\Omega(X,a)$. Or je me rappelle (\footnote[24]{C'est un erreur - ça doit être pour d'\textit{autres} objets groupe de Hot que les groupes de lacets. J'ai vu dans Quillen que $\Omega(X,x)$ doit se décrire par un groupe semi-simplicial (quand $X$ est un ensemble semi-simplicial, $x$ un sommet). Donc le problème si $C_{\infty} \simeq C$ dans Hot reste entier.}) qu'on m'a expliqué qu'en général, la loi de groupe de $\Omega(X,a)$ ne pourrait pas être obtenue par une loi \textit{associative} sur un ensemble semi-simplicial. Ainsi, l'insuccès de mes efforts semble tenir, non à ma maladresse, mais à la nature des choses. % [sic] C'est un erreur (note 24) -> C'est une erreur, genre

\dgrblock{dgrb341}{The cofinality condition that controls total categories}
\begin{dcomment}{exposition}{p112}{The cofinality condition that controls total categories}
Grothendieck is seeking a hypothesis that allows the base to be replaced without changing the total homotopy type. The appropriate general test is a condition on comma categories, stronger than cofinality of the associated preorders. It gives a precise sufficient criterion without requiring every reindexing functor between fibres to be a weak equivalence.

Let $\alpha:J\to T$ be a functor. If every comma category $\alpha/t$ has contractible nerve, then for every Grothendieck fibration $C\to T$ the projection
\[
 C\times_T J\longrightarrow C
\]
is a weak equivalence of nerves. To prove this, fix $x\in C_t$ and consider the comma category of the displayed projection over $x$. Its objects are triples consisting of $j\in J$, $y\in C_{\alpha j}$, and an arrow $y\to x$ over some $u:\alpha j\to t$. Forgetting $y$ gives a functor to $\alpha/t$. Cartesian lifting supplies a right adjoint: send $(j,u)$ to the cartesian arrow $u^*x\to x$. The adjunction follows from the unique factorization of an arrow through a cartesian lift. Hence these two comma categories have equivalent nerves. The dual of Quillen's Theorem~A now applies; see \citet[Corollary~4.1.3.3]{Lurie2009HTT}, applied to the opposite functor, for a precise modern formulation.

The condition is also necessary if the conclusion is required for \emph{every} fibration over $T$. Use the discrete fibration $T/t\to T$. Its total category has a terminal object, and its pullback is $\alpha/t$. Universal invariance of total homotopy types therefore forces $\alpha/t$ to be contractible. This elementary converse explains why a criterion expressed only in terms of which pairs of objects admit arrows cannot suffice.

For a cofiltered preorder $P$ and a decreasing cofinal sequence $t_n$, the category $\alpha/t$ is the nonempty tail of indices for which $t_n\le t$. It has a terminal object, the least such index. Thus the criterion proves the positive assertion in the text when $T=P$. For $T=\mathbf BG$ and the constant identity sequence, by contrast, $\alpha/t$ is a disjoint union of copies of $\mathbf N^{\mathrm{op}}$ indexed by $G$, so it is not contractible when $G\ne1$. These are the same two cases that the preceding examples separate. This comma criterion characterizes invariance for every fibration. The properness and smoothness conditions in the source instead exploit its hypothesis on parallel pullbacks, as the following proof shows; they need not imply the universal comma criterion.
\end{dcomment}


\dgrblock{dgrb342}{Why the smooth and proper criteria suffice}
\begin{dcomment}{later}{p112}{Why the smooth and proper criteria suffice}
The two sufficient conditions asserted here do follow from the earlier homotopy hypothesis. They are conditions on the base together with this particular class of coefficients; they need not make the chosen sequence cofinal for every fibration. The distinction matters because the hypothesis on parallel pullbacks supplies exactly the local constancy needed in the proof.

Write $q:T\to P=T_{\mathrm o}$, let $\alpha:\mathbf N^{\mathrm{op}}\to T$ be the chosen sequence, and replace the fibration by a strict diagram $F:T^{\mathrm{op}}\to\Cat$ as above. Put $K=NF$, regarded as a diagram of spaces, and $k=q^{\mathrm{op}}$. The fibre $E_p=q^{-1}(p)$ has the single object $p$ and endomorphism monoid $\operatorname{End}_T(p)$. Every endomorphism $u:p\to p$ is parallel to $\id_p$. Hence $F(u)$ is homotopic to the identity, so its nerve is a weak equivalence. Thus $K|_{E_p^{\mathrm{op}}}$ is a locally constant diagram: every one of its structure maps is a weak equivalence.

Use smoothness and properness for the localizer of weak equivalences of nerves. The fibre criterion of \citet[Theorem~1.2(b), printed p.~4]{Maltsiniotis2005Asphericite} says that smoothness makes
\[
 j_p:E_p\longrightarrow p\backslash q
\]
aspherical, meaning that every $j_p/z$ has contractible nerve. By the smooth--proper duality on pp.~4--5, properness makes
\[
 j_p:E_p\longrightarrow q/p
\]
coaspherical, meaning that every $z\backslash j_p$ has contractible nerve. In the first case coasphericity of $q$ says that $p\backslash q$ is contractible; in the second case asphericity of $q$ says that $q/p$ is contractible. Quillen's Theorem~A and its dual therefore show that $E_p$ is contractible in either case.

We shall use the following precise local-constancy fact. If $I$ has contractible nerve and $L:I\to\mathrm{sSet}$ sends every arrow to a weak equivalence, then for every $i\in I$ the canonical maps
\[
 L(i)\longrightarrow\operatorname{hocolim}_I L,
 \qquad
 \operatorname{holim}_I L\longrightarrow L(i)
\]
are weak equivalences. The first is \citet[Corollary~1.18, printed p.~599]{Cisinski2009LocallyConstant}; the second is its dual, applied to the opposite index category and opposite model category. The pointwise homotopy Kan-extension formula is Proposition~1.14 of the same paper, with the right-extension formula obtained by the same duality. These are statements about coherent diagrams, so no choice of mutually compatible homotopies between parallel pullbacks is required.

\emph{Smooth and coaspherical.} Let $G=\mathbf Lk_!K$ be the left homotopy Kan extension. Its value at $p$ is
\[
 G(p)\simeq
 \operatorname{hocolim}_{(p\backslash q)^{\mathrm{op}}}K
 \simeq\operatorname{hocolim}_{E_p^{\mathrm{op}}}K.
\]
Here the second equivalence follows because $j_p^{\mathrm{op}}$ is homotopy final: its under-comma categories are the opposites of the contractible categories $j_p/z$. The local-constancy fact shows that the adjunction unit $K\to k^*G$ is an objectwise weak equivalence. Let $v=(q\alpha)^{\mathrm{op}}:\mathbf N\to P^{\mathrm{op}}$. Cofinality of the chosen sequence in the preorder makes $v$ homotopy final: for $p\in P$, the category $p/v$ is the nonempty tail of indices for which $t_n\le p$, and it has a least element. Consequently
\[
 \operatorname{hocolim}_{\mathbf N}K\alpha^{\mathrm{op}}
 \xrightarrow{\ \sim\ }
 \operatorname{hocolim}_{\mathbf N}v^*G
 \xrightarrow{\ \sim\ }
 \operatorname{hocolim}_{P^{\mathrm{op}}}G
 \simeq\operatorname{hocolim}_{T^{\mathrm{op}}}K.
\]
The last equivalence is composition of left homotopy Kan extensions. By the unit and its adjunction identity, this composite is the canonical comparison of homotopy colimits.

\emph{Proper and aspherical.} Instead take $G=\mathbf Rk_*K$. The pointwise formula and the proper fibre criterion give
\[
 G(p)\simeq
 \operatorname{holim}_{(q/p)^{\mathrm{op}}}K
 \simeq\operatorname{holim}_{E_p^{\mathrm{op}}}K
 \xrightarrow{\ \sim\ }K(p).
\]
Now $j_p^{\mathrm{op}}$ is homotopy initial, which justifies restricting the homotopy limit, and the final arrow is the local-constancy equivalence. These maps are the components of the counit $k^*G\to K$, so this counit is an objectwise weak equivalence. Asphericity of $q$ also makes $k$ homotopy final, since $p/k\cong(q/p)^{\mathrm{op}}$. Both $k$ and $v=k\alpha^{\mathrm{op}}$ therefore give homotopy-colimit equivalences to $\operatorname{hocolim}_{P^{\mathrm{op}}}G$. It follows by two-out-of-three that
\[
 \operatorname{hocolim}_{\mathbf N}\alpha^{\mathrm{op}*}k^*G
 \longrightarrow\operatorname{hocolim}_{T^{\mathrm{op}}}k^*G
\]
is a weak equivalence. The counit transfers this conclusion to $K$.

Finally, the natural fibred form of Thomason's homotopy-colimit theorem identifies the comparison just proved with $C\times_T\mathbf N^{\mathrm{op}}\to C$. The sequential comparison already proved identifies its source with $C_\infty$ in $\Hot$. Thus each of the source's two conditions indeed makes its canonical $C_\infty\to C$ an isomorphism in $\Hot$.

This proof also identifies the exact limit of the statement. It used local constancy on $E_p^{\mathrm{op}}$, supplied by the homotopy hypothesis on $F$. For example, let $T$ be the one-object category of the monoid $\{1,0\}$ with absorbing zero. The natural transformation from the identity functor to the constant functor whose component is $0$ contracts its nerve. The map $q:T\to *$ is both smooth and proper, by the fibre criterion, and both aspherical and coaspherical. But for the constant identity sequence $\alpha$, the category $\alpha/t$ is a disjoint union of two copies of $\mathbf N^{\mathrm{op}}$. Thus these conditions do not imply the universal comma criterion for arbitrary fibrations; the coefficient hypothesis cannot be dropped from the argument.
\end{dcomment}


\dgrblock{dgrb343}{Two disconnected lifting categories}
\begin{dcomment}{examples}{p112}{Two disconnected lifting categories}
The assertion about the projection to the preorder has a small, concrete explanation.  Passing to the preorder remembers that a change of type exists and forgets how many such changes there are.  For smoothness or properness, those forgotten choices must be connected by commuting triangles of lifts.  Natural transformations between interval functors do not by themselves provide such triangles.

Write $q:\Ty\to T_o$ for the projection, with $T_o$ retaining the same objects and having one arrow whenever the corresponding Hom set in $\Ty$ is nonempty.  A necessary condition for $0$-smoothness is the following connected lifting property.  Given $a\in\Ty$, $b\in T_o$, and an arrow $b\to q(a)$, consider the category of arrows $x\to a$ with $q(x)=b$ lying over that arrow; morphisms are commuting triangles over $\mathrm{id}_b$.  This category must be nonempty and connected.  The dual condition, with fixed source, is necessary for $0$-properness.  These are the $W_0$ specializations of the smoothness criterion and its dual, where $W_0$ consists of functors inducing bijections on connected components \citep[Theorem~1.2(b), p.~4]{Maltsiniotis2005Asphericite}.

For smoothness take $b=+-+$ and $a=+$.  The interval of type $b$ has four vertices $0,1,2,3$, and the interval of type $a$ has two.  An increasing endpoint-preserving functor $b\to a$ is forced to collapse the middle, negative edge.  There are exactly two such functors, with vertex lists
\[
 (0,0,0,1)\qquad\hbox{and}\qquad(0,1,1,1).
\]
On the other hand, the only endpoint-preserving increasing endofunctor of $+-+$ is its identity.  To see this, its first internal value must be at most $1$, since a nonconstant image of the first positive edge cannot cross the following negative edge.  Its second internal value must be at least $2$, by the same argument at the last positive edge.  The middle negative edge can then be preserved only when these two values are $1$ and $2$.  As $q$ is the identity on objects, the strict fibre over $b$ has this one object and only its identity arrow.  The lifting category therefore has precisely two objects and only their identities.  It is disconnected, proving that $q$ is not $0$-smooth.

For properness take the fixed source $a=++$ and target $b=+$.  There are exactly two endpoint-preserving increasing functors $a\to b$, with vertex lists
\[
 (0,0,1)\qquad\hbox{and}\qquad(0,1,1).
\]
The only endpoint-preserving increasing endofunctor of $+$ is the identity.  The fixed-source lifting category is consequently another discrete category with two objects.  Thus $q$ is not $0$-proper.

Both examples are compatible with the earlier contractibility theorem for nonempty categories of interval functors.  In that theorem the morphisms are endpoint-fixed natural transformations.  Here they are factorizations through endofunctors of a fixed type.  The distinction explains why connectedness of every mapping category cannot supply the lifting property required on this page.
\end{dcomment}


\dgrblock{dgrb344}{What the preorder forgets about a lift}
\begin{dcomment}{examples}{p112}{What the preorder forgets about a lift}
Grothendieck's proposed sufficient criterion fails for the category of types for a concrete reason: its preorder identifies distinct changes of parameter which cannot be connected inside a lifting category. The definition on manuscript page~99 retains the same objects in $\Ty_o$ and replaces every nonempty set of morphisms by a singleton. It does not identify different objects whose images happen to be isomorphic in that preorder.

Take
\[
 S=(0\longrightarrow1\longleftarrow2\longrightarrow3),
 \qquad T=(0\longrightarrow1).
\]
An endpoint-preserving monotone type map $S\to T$ changes its vertex value from $0$ to $1$ at exactly one of the two forward edges. There are precisely two such maps, with vertex lists
\[
                    (0,1,1,1),\qquad(0,0,0,1).
\]
Both $S$ and $T$ have only the identity endomorphism in $\Ty$. For $S$, a monotone endomorphism can advance by at most one vertex on each edge, since its alternating target has no nonidentity composite of two consecutive edges. Its three increments must sum to three because the endpoints are fixed, so each increment is one. The assertion for $T$ is immediate.

Let $p:\Ty\to\Ty_o$ and consider the unique base arrow $pS\to pT$. An incoming lift with target $T$ has source exactly $S$, and its morphisms are endomorphisms of $S$ commuting with the lifted arrow. Its lifting category is therefore the discrete category on the two displayed maps. The outgoing lifting category with source $S$ is likewise discrete on those maps, now because its morphisms are endomorphisms of $T$. Neither category is connected. The incoming and outgoing lift criteria consequently show that $p$ is neither $0$-smooth nor $0$-proper, and hence neither smooth nor proper. This verifies the source's assertion; failure of a sufficient criterion still does not imply failure of the geometric comparison itself.
\end{dcomment}


Je vais résumer les faits positifs auxquels je suis arrivé dans le diagramme suivant, qui a lieu dans la situation générale $C \longrightarrow T$ envisagée au début du présent paragraphe :

\dgrblock{dgrb345}{The obstruction withdrawn in the footnote}
\begin{dcomment}{context}{p113:Je vais résumer les faits positifs}{The obstruction withdrawn in the footnote}
Grothendieck is trying to explain why an associative composition is easy on all finite paths but elusive on a chosen infinite refinement. The proposed explanation on this page is withdrawn in footnote~24, and the withdrawal is mathematically decisive. A based loop space does admit a strictly associative model. One elementary topological model consists of Moore loops: a loop is a pair $(\ell,\gamma)$, where $\ell\geq0$, $\gamma:[0,\ell]\to X$, and both endpoints are the basepoint. Concatenation adds lengths and joins the maps without rescaling their domains. Its associativity is an equality, and the loop of length zero is a strict unit. Reversal is an inverse up to homotopy, which is exactly the group law required in $\Hot$; a strict monoid model does not require strict inverses.

The simplicial statement is stronger. For a reduced simplicial set $K$, Kan's loop-group construction gives a simplicial group $GK$ and a weak equivalence $K\to\overline W GK$. Consequently $|GK|$ models the based loop space of $|K|$. We checked the precise reducedness hypothesis and the unit equivalence in \citet[Theorem~21]{Stevenson2012}, which gives a modern proof of Kan's theorem. The word ``semi-simplicial'' in this part of the manuscript includes the simplicial formalism with degeneracies; it must not be read as a claim about an arbitrary semisimplicial object without them.

These constructions settle the alleged obstruction, but do not by themselves identify either of Grothendieck's particular categories with the model. That identification is the work still to be done. For $\Sigma=\mathrm{Fl}(X)$, section~13 of the present fragment eventually compares the all-types category with $\Ch_\infty$ and restores the strict concatenation model. The footnote should therefore be read as a correction to an inference, not as an incidental qualification of a theorem. The unsuccessful search described in the surrounding prose is part of the discovery of the later proof.
\end{dcomment}


\msp{page 114}
% foot: Version 22 décembre 2001
\printedpage{72}
% head: A. Grothendieck : Dérivateurs, Ch. VII, Édition des Universités de Montpellier II et Paris VII
\[
\begin{tikzcd}[column sep=small, row sep=huge]
 & & C_{t'} \arrow[d, "{u^{*}\quad\lab{(pour $u : t \longrightarrow t'$\\ dans $T$)}}"] \arrow[ldd, dotted] \arrow[rdd, dotted] & \\
 & & C_{t} \arrow[ld, swap, "{\lab{flèche définie dans\\ Ind\,Cat\,hot\\ (\texttt{[ou dans]}\\ IndHot,\\ Ind\,($\Hotz$))}}"] \arrow[rd, "{\lab[r]{flèche dans Cat\\ (donc dans\\ Cat\,hot,\\ Hot,\\ $\Hotz$)}}"] & \\
\varinjlim C_{t} = & |[draw, circle]| H_{T} \arrow[rd, swap, "{\lab{flèche\\ définie dans\\ Ind\,Cat\,hot\\ ou dans Ind\,Hot\\ ou dans Ind\,$\Hotz$}}"] & \lab{diagramme\\ commutatif\\ dans Ind\,Hot\\ ou Ind\,($\Hotz$)} & C \\
\lab{dans Ind\,Cat\,hot\\ (ou dans Ind\,Hot\\ ou dans Ind\,($\Hotz$))} & & |[draw, circle]| C_{\infty} \arrow[ru, swap, "{\lab[r]{flèche\\ définie\\ dans Hot,\\ donc dans $\Hotz$,\\ en général pas un\\ isomorphisme, même\\ dans $\Hotz$}}"] & \\
 & & = \varinjlim H_{T} \text{ dans } \Hotz &
\end{tikzcd}
\]

\dgrblock{dgrb346}{Strict loop multiplication and the later comparison}
\begin{dcomment}{status}{p113}{Strict loop multiplication and the later comparison}
Grothendieck wants to retain the strictly associative composition already constructed on the total category of paths.  His hesitation here concerns whether that useful algebraic structure can live on a correct model of the loop space.  Footnote~24 retracts the proposed obstruction, and the retraction is mathematically decisive.  A multiplication that is associative only in the homotopy category need not come equipped with the coherent data needed for a strict model.  A loop space, however, already has those data; its origin as a space of loops supplies much more than an isolated group object in the homotopy category.

There is a precise classical replacement.  For a reduced simplicial set $K$, Kan constructs a simplicial group $GK$ and a principal $GK$-bundle over $K$ with weakly contractible total space \citep[Definition~6.3 and Theorem~7.1]{Kan1958Combinatorial}.  After passage to a reduced Kan model, this identifies $GK$ with the based loop homotopy type.  Its multiplication and inverses are strict in every simplicial degree.  Thus the existence of a strictly associative simplicial model cannot disprove the comparison proposed here.  This assertion concerns the existence of a model, rather than a claim that a particular category of paths is that model; the latter still requires a comparison argument.

The same manuscript subsequently supplies that argument for unrestricted paths.  In Chapter~VII, \S13, manuscript pages~157--170, Grothendieck returns to the case $\Sigma=\Fl(X)$ and obtains the comparison through the two projections from the sequential Grothendieck construction.  The endpoint functors and their fibres are essential; a weak equivalence of total categories alone would not justify a statement about their strict fibres.  The elementary lemma that constant paths induce a homotopy equivalence $X\to\operatorname{Fun}(I_\tau,X)$ is left without a written proof at manuscript page~170.  The doubts recorded on the present page therefore belong to the progress of the argument, not to its final conclusion \citep[Ch.~VII, \S13]{GrothendieckDerivateursVII}.

That remaining lemma has a direct proof.  Delete the last vertex of a nontrivial oriented interval, retract it to its neighbour, and let $j$ denote the inclusion of the remaining interval.  The retraction $r$ is a functor, $rj$ is the identity, and either $jr\Rightarrow\mathrm{id}$ or $\mathrm{id}\Rightarrow jr$ is a natural transformation, according to the direction of the deleted edge.  Iterating gives a finite zigzag between the identity of $I_\tau$ and its constant functor at the initial vertex.  Precomposition gives the corresponding zigzag between the identity of $\operatorname{Fun}(I_\tau,X)$ and the composite of evaluation at that vertex followed by the constant-path functor.  The other composite is strictly the identity of $X$.  Thus constant paths give a categorical homotopy equivalence, and hence a weak equivalence of nerves.

For the audit this distinction has two consequences.  The inference from an alleged impossibility of strict loop multiplication is false as written, and is explicitly withdrawn by the author.  The unrestricted path comparison must be recorded with its later positive treatment, whereas the unrestricted \emph{general-fibration} comparison of \S10 still has the counterexamples discussed above.  These statements have different hypotheses.  Neither footnote~24 nor the later path argument removes the base-category obstruction for an arbitrary $C\to T$.
\end{dcomment}


\texttt{[Les annotations dans ce diagramme concernent les flèches en traits pleins.]}

Les objets importants sont $H_{T}$ et $C_{\infty}$, alors que $C$ (sur $T$) et les $C_{t}$ prennent un rôle technique intermédiaire, pour construire $H_{T}$ et $C_{\infty}$. On peut considérer $C_{\infty}$ comme une catégorie (objet de (Cat)), mais dont la construction dépend de nombreux arbitraires, à l'inverse de $H_{T}$ dans Ind\,Cat\,hot. C'est l'objet de $\Hotz$ déduit de $C_{\infty}$ qui a un sens canonique, et peut être décrit comme la limite inductive de $H_{T}$ comme objet de Ind\,$\Hotz$, i.e. comme représentant le même foncteur sur $\mathrm{Hot}_{0}$ que $H_{T}$.

\msp{page 115}
Dans le cas géométrique $T = \Ty$, et $C$ de la forme $\Ch_{\Sigma}(X;a,b)$, donnant lieu à $C_{\infty}(a,b)$, on \textit{n'a pas} de définition naturelle d'un bifoncteur de composition des chemins
\[
C_{\infty}(a,b) \times C_{\infty}(b,c) \ \longrightarrow\ C_{\infty}(a,c).
\]
On peut définir, certes, des accouplements
\[
\Ch_{\Sigma,\tau}(a,b) \times \Ch_{\Sigma,\tau'}(b,c) \ \longrightarrow\ \Ch_{\Sigma,\tau\tau'}(a,c),
\]
ou, avec les notations relatives à la donnée de
\[
\begin{aligned}
&\underline{t} : \mathbf{N}^{o} \longrightarrow \Ty,\\
&\Ch_{n}(a,b) \times \Ch_{m}(b,c) \longrightarrow \Ch_{n+m}(a,c),
\end{aligned}
\]
% foot: Version 22 décembre 2001
\printedpage{73}
% head: A. Grothendieck : Dérivateurs, Ch. VII, Édition des Universités de Montpellier II et Paris VII
mais ces accouplements \textit{ne sont pas} compatibles avec les morphismes de transition des trois systèmes inductifs de catégories $(\Ch_{n})_{n\in\mathbf{N}}$, ils ne le sont \textit{qu'à homotopie près.} Cela suffit cependant pour avoir des accouplements parfaits
\[
H_{T}(a,b) \times H_{T}(b,c) \ \longrightarrow\ H_{T}(a,c),
\]
avec associativité, et (si $\Sigma = \mathrm{Fl}(X)$ tout entier) existence de \textit{vrais} inverses, i.e. de flèches (dans Ind\,Cat\,hot) de `passage au chemin opposé'

\msp{page 116}
\[
H_{T}(a,b) \ \longrightarrow\ H_{T}(b,a), \qquad c \longmapsto c^{o}, \tag*{$(*)$}
\]
donnant lieu au diagramme commutatif exprimant que c'est bien un inverse pour la composition. Une autre façon de le dire, c'est que pour tout objet $Y$ de Cat\,hot (ou de Ind\,Cat\,hot), les ensembles
\[
\mathrm{Hom}_{\mathrm{Cat\,hot}}(Y, H_{T}(a,b)) \overset{\text{déf}}{=} H_{T}^{Y}(a,b)
\]
forment, pour $a,b \in X$ variables, les flèches d'un \textit{groupoïde} (\footnote[25]{Si $\Sigma = \mathrm{Fl}(X)$. Si $\Sigma$ est plus petit, au lieu d'un groupoïde, on trouve une \textit{catégorie} sans plus.}) $\mathrm{Ch}_{\Sigma}^{Y}(X)$ ayant pour ensemble de sommets $\mathrm{Ob}\,X$, et $H_{T}^{Y}(a,b)$ étant formé des flèches de $a$ vers $b$. L'application
\[
H_{T}^{Y}(a,b) \longrightarrow H_{T}^{Y}(b,a), \qquad c \longmapsto c^{o},
\]
déduite du foncteur $(*)$ plus haut n'est autre que l'application de passage à l'inverse. Cette catégorie n'est d'ailleurs autre que celle déduit de % [sic] celle déduit -> celle déduite, accord
\[
X(Y)\Sigma_{Y}^{-1}, \qquad \text{où } \left\{
\begin{aligned}
&X(Y) = \underline{\mathrm{Hom}}_{\mathrm{Cat}}(Y,X)\\
&\Sigma_{Y} \subseteq \mathrm{Fl}(X(Y)),\\
&\text{défini en termes de } \Sigma \text{ `argument par argument'},
\end{aligned}
\right.
\]
à l'aide du foncteur
\[
X \longrightarrow X(Y), \qquad a \longmapsto a_{Y},
\]
\msp{page 117}
associant à tout $a \in X$ le foncteur \textit{constant} sur $Y$ de valeur $a$, en prenant
\[
H_{T}(a,b) \simeq \mathrm{Hom}_{X(Y)\Sigma_{Y}^{-1}}(a_{Y}, b_{Y}) \ .
\]

\dgrblock{dgrb347}{The stage of a concatenation}
\begin{dcomment}{development}{p115}{The stage of a concatenation}
Grothendieck's ind-object accommodates concatenation because a sufficiently refined type can dominate any two given types and their concatenation. The notation $n+m$ is convenient for a multiplicative standard sequence, but does not follow from an arbitrary cofinal sequence alone. The following small calculation locates the missing convention.

\textbf{A slow cofinal sequence.}\ours\quad There is a cofinal decreasing sequence of types $t_n$ for which no type map $t_{n+m}\to t_nt_m$ exists for some $n,m$.

\emph{Proof.} Put $w=\Delta_+\Delta_-$ and
\[
 t_n=w^{k(n)},\qquad k(n)=\lfloor\log_2(n+1)\rfloor.
\]
Use identity maps when $k(n+1)=k(n)$ and collapse the final copies of $w$ otherwise. This gives $t_{n+1}\to t_n$. It is cofinal: every finite signed word is obtained from a sufficiently long alternating word by deleting unwanted edges and expanding the selected edges into its maximal constant-sign runs. For $n=m=3$, however, $t_6=w^2$, whereas $t_3t_3=w^4$. A type map from a word to a reduced alternating target must use at least one source edge for every target edge, since the image of a single source edge is a constant-sign run. Four source edges cannot cover the eight edges of $w^4$. Thus there is no such map.\hfill$\square$

The obstruction is visible on path categories themselves. Take $X=I_{w^4}$, with endpoints $a,b$, and let $c$ be its middle vertex. The inclusions of its two four-edge halves give objects of $\uCh_{t_3}(X;a,c)$ and $\uCh_{t_3}(X;c,b)$. The category $\uCh_{t_6}(X;a,b)$ is empty: every nonidentity arrow of this alternating interval joins adjacent vertices, so four transitions cannot cover the graph distance eight between its endpoints. Hence a pairing with target stage $6$ cannot exist in this example, even with $\Sigma=\mathrm{Fl}(X)$.

For a general chain the correct target is a chosen $k=k(n,m)$ admitting $t_k\to t_nt_m$. Cofinality guarantees such a $k$, and passage to a larger stage or a different chosen map makes no difference in $\Cathot$, by the homotopy theorem. Equivalently, one may choose the standard sequence $t_n=w^n$, where $t_nt_m=t_{n+m}$. We checked the type-map description in the preceding part of this chapter and the usual filtered-colimit formalism; we did not find this particular indexing correction recorded in the literature. Its purpose is to make the intended ind-construction exact, not to alter that construction.
\end{dcomment}


\dgrblock{dgrb348}{Parameters turn the ind-object back into localization}
\begin{dcomment}{exposition}{p116}{Parameters turn the ind-object back into localization}
Grothendieck ends this part by explaining why the ind-object is geometrically useful: its values on parameter categories recover ordinary localization of diagram categories. The assertion has a direct proof, and it supplies associativity and inverses at the same time.

For an actual small category $Y$, write $X(Y)=\uHom(Y,X)$ and let $\Sigma_Y$ consist of natural transformations whose components belong to $\Sigma$. It contains identities and is closed under composition. A functor
\[
 Y\longrightarrow\uCh_{\Sigma,\tau}(X;a,b)
\]
is the same thing as a functor $Y\times I_\tau\to X$ whose endpoints are the constant functors $a_Y,b_Y$, with identity endpoint components on morphisms in $Y$, and whose backward arrows belong pointwise to $\Sigma$. By exchanging the two variables, it is exactly a $\Sigma_Y$-path of type $\tau$ in $X(Y)$ from $a_Y$ to $b_Y$. This identification respects natural transformations and every type-reindexing functor. The two total categories over $\Ty$ are therefore isomorphic.

Taking connected components and applying the localization theorem of \S7 gives
\[
 \begin{split}
 H_T^Y(a,b)
 &:=\Hom_{\operatorname{Ind}(\Cathot)}(Y,H_T(a,b))\\
 &\cong\Hom_{X(Y)[\Sigma_Y^{-1}]}(a_Y,b_Y).
 \end{split}
\]
This corrects two notational slips in the printed formulas: the ambient category on manuscript page~116 is $\operatorname{Ind}(\Cathot)$, and the left side on page~117 needs the superscript $Y$. No finiteness of $Y$ is required for this formula. The uniform type of a $Y$-family is retained by the ind-object; replacing it by $C_\infty$ is the separate step for which finite presentation was needed.

Concatenation on either side is the same operation on functors $Y\times I_\tau\to X$, so this bijection respects composition and units. It follows that the sets $H_T^Y(a,b)$ are the hom-sets of the full subcategory of $X(Y)[\Sigma_Y^{-1}]$ on the constant diagrams. They form a groupoid when $\Sigma$ is all arrows: the reversed path represents the inverse, and the backtracking contractions of \S7 are natural in $Y$. Yoneda in the ind-category turns these pointwise identities into the associativity, unit, and inverse identities of the ind-valued path system. Finite products of ind-objects are represented by the corresponding product diagrams, so the pairings have exactly the domain stated in the manuscript.

If $Y$ itself is an ind-object, the abstract hom-set on the first line still makes sense. For a presentation $Y=(Y_i)$ it is $\lim_i H_T^{Y_i}(a,b)$. The expression $\uHom(Y,X)$ on the second line is only the stated functor category when $Y$ is an actual small category; it should not be silently applied to a formal ind-object. The natural extension is the displayed inverse limit on source indices.
\end{dcomment}


\dgrblock{dgrb349}{What the ind-object represents}
\begin{dcomment}{exposition}{p117:associant à tout $a \in X$}{What the ind-object represents}
The ind-object retains a uniform bound on the complexity of a family of paths. This is why it can encode localization before a single category has acquired a satisfactory geometric interpretation. With $[Y,D]_{\mathrm{cat}}=\pi_0\underline{\mathrm{Hom}}(Y,D)$, the formula is
\[
 H_T^Y(a,b)=\mathop{\rm colim}_n[Y,C_n(a,b)]_{\mathrm{cat}}
 \cong\mathrm{Hom}_{X(Y)[\Sigma_Y^{-1}]}(a_Y,b_Y).
\]
Here $\Sigma_Y$ consists of natural transformations whose components belong to $\Sigma$, and the right side uses the constant diagrams. A morphism on the right is represented by a finite zigzag of natural transformations. Its finite length is a bound shared by every object of $Y$; it therefore gives a map into one stage on the left. Refinement and the elementary relations of localization give the identifications in the colimit. This is the mechanism already discussed in the commentary on Chapter~V, manuscript pages~30--33, concerning the parameter category and ind-representation. The formula printed on page~117 omits the superscript $Y$ from $H_T^Y$; without it an ind-object is being equated to a set.

For finite $Y$, the natural map
\[
 \mathop{\rm colim}_n[Y,C_n]_{\mathrm{cat}}
       \longrightarrow [Y,\mathop{\rm colim}_n C_n]_{\mathrm{cat}}
\]
is bijective. Indeed a functor uses finitely many objects, arrows, and relations, all of which lift to one stage. A finite zigzag of natural transformations uses finitely many more data, so equality also occurs at one stage. Finite presentation, rather than any topological fibrancy, is doing the work. For infinite $Y$, a functor to the colimit can have an unbounded range of stages, and this argument stops.

Composition of localized zigzags is associative because juxtaposition of finite strings is associative. When all arrows of $X$ are inverted, reversal supplies the inverse string. Hence the family of sets $H_T^Y(a,b)$ is a groupoid with object set $\mathrm{Ob}(X)$. For a general $\Sigma$ it is only a category. Passing from this set-valued representation to mapping \emph{spaces} is an additional step: the notation does not identify the ordinary nerve of every functor category with a derived mapping object.
\end{dcomment}


Dans le même cas géométrique, je me suis convaincu d'ailleurs que $C_{\infty}$ était, en tant qu'objet de Hot (pas seulement de $\Hotz$) \textit{unique à} Hot-\textit{isomorphisme près,} mais (chose étrange) peut-être pas à Hot-isomorphisme \textit{unique} près. De façon générale, revenant aux
% foot: Version 22 décembre 2001


\dgrblock{dgrb350}{From path components to mapping spaces}
\begin{dcomment}{later}{p117}{From path components to mapping spaces}
The last formula brings the discussion back to its original purpose: families of paths should describe morphisms after localization, and should continue to do so when an indexing category is introduced.  At the level of sets of components, the preceding arguments give a complete description.  To retain higher homotopies, a modern treatment keeps the diagram of zigzag categories together with the coherences in its changes of shape.

One precise formulation is the hammock formula.  For a relative infinity-category $(\mathcal C,\mathcal W)$, let $\mathrm{Zig}_{Z}(\mathcal C,\mathcal W)_{x,y}$ denote the infinity-category of zigzags of shape $Z$ with fixed endpoints and pointwise weak-equivalence transformations.  Write $|{-}|$ for groupoid completion, which for an ordinary category is represented by the realization of its nerve.  Then, with colimit taken in spaces,
\[
 \operatorname{Map}_{\mathcal C[\mathcal W^{-1}]}(x,y)
 \simeq
 \mathop{\operatorname{colim}}_{Z\in\mathrm{Zig}_{\partial}^{\mathrm{op}}}
 \bigl|\mathrm{Zig}_{Z}(\mathcal C,\mathcal W)_{x,y}\bigr|.
\]
Here $\mathrm{Zig}_{\partial}$ has endpoint-preserving changes of shape.  Arakawa and Cnossen prove this formulation in their 2026 preprint, Definition~1.7, Theorem~1.9 and Corollary~1.11; Remark~1.10 identifies composition with concatenation \citep{ArakawaCnossen2026}.  Their \S4 compares the construction with Dwyer--Kan's original hammock localization.

The notation must not obscure a difference in definitions.  Grothendieck's $\Ch_{\Sigma,\tau}$ on these pages is a \emph{full} subcategory of the functor category: membership constrains the backward arrows of an object, while morphisms are arbitrary natural transformations with the specified endpoint condition.  The hammock construction just quoted also constrains the vertical components of those transformations.  A comparison of the two models for general $\Sigma$ is an additional theorem, not a consequence of recognizing the same words of signs.  When all arrows are marked this difference disappears, and the later comparison within Chapter~VII is directly relevant.  This qualification is particularly useful when passing from $X$ to $X(Y)$: it prevents an equality of localized morphism \emph{sets} from silently being promoted to an identification of mapping \emph{spaces}.
\end{dcomment}


\dgrblock{dgrb351}{Independence of the sequential homotopy type}
\begin{dcomment}{development}{p117}{Independence of the sequential homotopy type}
Grothendieck's final distinction between an isomorphism and a uniquely determined isomorphism can be made precise without solving the comparison with the total category $C$.  The two questions separate: a bad base can obstruct $C_\infty\to C$ while all the sequential models still have the same weak homotopy type.  The following argument gives the latter conclusion under the abstract hypotheses of \S10.  Its ingredients are standard telescope arguments; we include their application here because we did not locate this exact reachability-preorder statement in the sources checked.  For the usual sequential homotopy-colimit construction see \citet[\S14.9, Proposition~14.10 and Corollary~14.11]{DuggerHocolim}.

\textbf{Proposition (independence of choices).}\ours\quad
Let $T$ be a nonempty small category and let $F:T^{\mathrm{op}}\to\Cat$ be a strict functor.  Suppose that $F(u)$ and $F(v)$ are joined by a finite zigzag of natural transformations whenever $u,v$ are parallel arrows of $T$.  Let
\[
 t_0\longleftarrow t_1\longleftarrow\cdots,
 \qquad s_0\longleftarrow s_1\longleftarrow\cdots
\]
be two chains whose objects are cofinal to the left in the reachability preorder of $T$.  Then
\[
 \varinjlim_n F(t_n)\quad\hbox{and}\quad\varinjlim_m F(s_m)
\]
have isomorphic images in $\Hot$.  The assertion includes arbitrary choices of the arrows in either chain.  It does not assert a preferred isomorphism.

\textit{Proof.}
We first record precisely which properties of sequential homotopy types will be used.  For a sequence of small categories $A_n$, the nerve commutes with its filtered colimit.  Indeed, a simplex is a finite string of objects and arrows subject to finitely many incidence equations; these data, and equalities between two presentations of them, occur at some common stage.  Hence
\[
 N\bigl(\varinjlim_n A_n\bigr)=\varinjlim_n N A_n.
\]
For a sequence of simplicial sets $K_n$, the canonical map
\[
 \operatorname{Tel}(|K_n|)\longrightarrow
       \bigl|\varinjlim_n K_n\bigr|
\]
is a weak equivalence.  Here the telescope glues the upper end of
$|K_n|\times[0,1]$ to the next space by its transition map.  To check this assertion on homotopy groups, replace simplicial sets by $\operatorname{Ex}^{\infty}K_n$.  This replacement commutes with filtered colimits: $\operatorname{Sd}\Delta^r$ is finite for each $r$, so $\operatorname{Ex}$ commutes with them, and the sequential iteration does also.  A filtered colimit of Kan complexes is Kan, because every horn is finite.  Maps and homotopies defining a homotopy group use finite simplicial data, so homotopy groups, with a basepoint represented at some stage, commute with this filtered colimit.  On the telescope side a map or a homotopy from a compact sphere or cylinder has bounded height.  Each bounded initial telescope retracts to its final space.  Its homotopy groups therefore give the same direct limits.  The canonical map induces these identifications, including the statement about components, and is a weak equivalence.

We also need invariance when each transition map is replaced by a homotopic map.  Suppose $f_n,g_n:A_n\to A_{n+1}$ are continuous maps with homotopies $H_n$ from $g_n$ to $f_n$.  A map from $\operatorname{Tel}(f)$ to $\operatorname{Tel}(g)$ is the identity on every level $A_n$.  On the cylinder for $f_n$, it first traverses the cylinder for $g_n$, then follows $H_n$ in level $A_{n+1}$.  More explicitly it sends $(x,r)$, for $0\leq r\leq\tfrac12$, to $(x,2r)$ in the $g_n$ cylinder, and for $\tfrac12\leq r\leq1$ to $H_n(x,2r-1)$ in level $A_{n+1}$.  The endpoint at $r=1$ is $f_n(x)$, so these maps glue.  Reversing the homotopies gives a map back.  Their composites are homotopic to the identities, fixing the levels, by cancelling a path immediately followed by its reverse and reparametrizing the cylinders.  These cancellations are defined separately on each cylinder and agree at its ends; they consequently glue on the entire telescope.  No compatibility between different $H_n$ is required.  Finally, passing to a cofinal subsequence preserves the weak homotopy type of a telescope, by the same bounded-height calculation of homotopy groups.

We now compare the chains.  Choose indices, increasing without bound, and arrows in $T$ to obtain an interleaving
\[
 t_{n_0}\longleftarrow s_{m_0}\longleftarrow
 t_{n_1}\longleftarrow s_{m_1}\longleftarrow\cdots.
\]
At each step this is possible by cofinality of the opposite chain; after finding an index, moving farther down that chain allows the indices to increase.  Apply $F$ and take nerves and realizations.  The even subsequence has objects $|NF(t_{n_i})|$.  Its bonding map is obtained from the composite of two of the newly chosen arrows.  That composite and the original transition $t_{n_{i+1}}\to t_{n_i}$ are parallel arrows of $T$.  By hypothesis their pullbacks are connected by natural transformations; taking nerves turns these into a homotopy of continuous maps.  The preceding telescope construction therefore identifies the homotopy type of this even subsequence with that of the original $t$-sequence.  The odd subsequence has the corresponding property for the $s$-sequence.  Both are cofinal subsequences of the same interleaved telescope.  Combining these comparisons with the telescope-to-colimit weak equivalences proves the proposition.\hfill$\square$

\end{dcomment}
\begin{otherlanguage}{english}
\dgroverview{VII-b}{Chapter VII \textperiodcentered\ \textup{VII-b} \textperiodcentered\ Overview}
\selectlanguage{english}
\noindent\textbf{Paths, localization, and the information in an indexing category}



\medskip
These pages occupy the middle of Chapter~VII.  They begin with the end of the combinatorial analysis of path types, establish the ordinary category of fractions by means of path categories, and then examine whether a category containing all shapes of paths has the same homotopy type as a sequentially enlarged model.  Their central concern is a compatibility problem.  Concatenation is completely satisfactory before taking limits: it is functorial, associative and unital.  The sequential models make other calculations accessible, but their transition maps need not respect that composition strictly.  Grothendieck is trying to preserve both advantages in one construction.

The category $\underline{\mathrm{Ty}}$ consists of finite oriented intervals, equipped with their spatial order and their two endpoints.  Its arrows preserve that order, the categorical direction of each traversed edge, and the endpoints.  Concatenation makes it a strict monoidal category.  The opening calculation describes its nerve degree by degree as a free monoid: a composable string of type maps splits uniquely over the edges of its initial interval.  The data attached to one initial edge are successive finite weightings recording how many edges replace it at each stage.  The commentary on the first part of the chapter explains this calculation through manuscript page~57.  The commentary here adds the distinction needed for its later use: a face can send an indecomposable generator to a product.  Degreewise freeness is therefore not freeness on a simplicial set.

For an ordinary category $X$, a path of type $\tau$ is a functor $I_\tau\to X$.  Between paths of the same type one keeps natural transformations, and between different types one also keeps pullback along interval maps.  These categories assemble into a category fibred over $\underline{\mathrm{Ty}}$.  Fixing endpoints $a,b$ means taking the strict fibre of the endpoint functor: the components of a natural transformation at both ends must be identities.  This requirement is responsible for several of the correct diagrams in the text.  Dropping it would change the hom-sets obtained after passing to components.

If $\Sigma$ is a wide subcategory of $X$, we allow backward edges only when they lie in $\Sigma$, and write $\Ch_\Sigma(X;a,b)$ for the resulting total fixed-endpoint path category.  Closure of $\Sigma$ under identities and composition is exactly what makes this construction stable under every type map.  A source edge can collapse, or cover a consecutive block of target edges; the two possibilities require precisely these two closure conditions.  No calculus-of-fractions hypothesis is needed at this stage.  The fibres are full subcategories of functor categories, so their internal vertical transformations remain unrestricted.

The first decisive result is
\[
 \operatorname{Hom}_{X[\Sigma^{-1}]}(a,b)
      \cong \pi_0\Ch_\Sigma(X;a,b).
\]
Concatenation induces the composition on the right.  A one-edge path followed by its reversal is connected to a constant path through a ladder with identity endpoint components, and the constant path is connected to the empty path by a type collapse.  Thus every arrow of $\Sigma$ becomes invertible.  For a functor $F:X\to Y$ which inverts $\Sigma$, evaluation of a zigzag in $Y$ is unchanged under either a ladder or a type change.  In a ladder, the equations for its successive squares telescope to
\[
 \eta_n E_F(c)=E_F(c')\eta_0;
\]
the endpoint identities give equality of effects.  The argument uses inverses of the backward horizontal edges but never inverses of the internal vertical components.  It supplies the omitted well-definedness check for the map $\beta$ on manuscript page~68.  The resulting evaluation has the required unique factorization property.  This is the classical localization problem of \citet[Chapter~I, \S1]{GabrielZisman1967}, with the witnesses for its relations retained as categorical data.

Two examples make the theorem tangible.  Inverting one arm of a span creates the fraction across the span and produces a category with a walking-arrow skeleton.  Inverting both arrows of the category $a\rightrightarrows b$ produces a connected groupoid with automorphism group $\mathbf Z$; the nerve of the original category is a circle.  The complete calculations distinguish cancellation of backtracking from the survival of a genuine loop.  They also show why a relative localization need not be a groupoid.

The next task is to establish that concatenation really acts on all morphisms of the total path categories.  A ladder can be glued to a second ladder because their common endpoint components are identities.  A change of type can be glued to a second change of type because the common endpoint has a specified image.  Writing a total morphism as a type map together with a vertical transformation makes the functoriality and strict associativity transparent.  The commentary gives this calculation explicitly.  It also repairs an edition diagram which glues over the two-element set $\{0,1\}$ where the construction requires a single identified endpoint.

Sections~8 and~9 then explain why the particular interval map inducing a pullback ceases to matter after categorical homotopy.  The relevant homotopy relation joins functors by finite zigzags of natural transformations.  Grothendieck proves that every nonempty category of endpoint-preserving interval maps is connected, and asks whether it is contractible.  We give a complete categorical contraction.  After a forced initial collapse when the first signs disagree, the first target run can be filled by the first source edge.  Clamping the remaining vertex images to the end of that run defines an idempotent functor and a natural transformation in the direction determined by the run's sign.  The retract is a smaller interval-mapping category.  Induction contracts it to a point.  This is marked as a development of the present commentary, without a claim of priority.

Precomposition transports these natural transformations to the path categories.  Because the endpoint components remain identities, the homotopy takes place over $X\times X$ and restricts to every fixed-endpoint fibre.  In modern notation, the diagram of fibres becomes constant on parallel arrows after passage to $\Cathot$, the category whose morphisms are categorical-homotopy classes of functors.  This is the bridge to the abstract discussion of a fibration $C\to T$ in \S10.

Here three objects must be distinguished.  There is the total category $C$, the ind-object determined by its fibres in $\Cathot$, and an ordinary sequential colimit $C_\infty$ formed after choosing a chain cofinal in the reachability preorder $T_{\mathrm o}$.  The source assumes that this preorder is cofiltered.  It does not assume that parallel arrows of $T$ can be equalized.  Forgetting that difference would conceal the principal obstruction: for the constant point-valued diagram over a one-object group $BG$, the ind-object and every sequential model are points, while the total category is $BG$.

There is a further failure on manuscript page~103.  The reachability preorder of a slice $T/t_0$ need not be the lower ideal of $t_0$ in $T_{\mathrm o}$.  The category with two parallel arrows already disproves that assertion.  The lower ideal may still be used by restricting to the full subcategory of objects which admit an arrow to $t_0$, but this is only cofinality on reachability preorders.  Comparison of total categories requires the actual comma categories.  The commentary formulates and proves the corresponding sufficient-and-necessary condition for base change to preserve the homotopy type of every Grothendieck fibration, using Quillen's Theorem~A \citep[\S1]{Quillen1973}.  The source's separate properness and smoothness criteria are also proved under its hypothesis on parallel pullbacks.  Contractibility of the strict fibres of $T\to T_{\mathrm o}$, together with local constancy of the coefficients there, allows homotopy Kan extension to descend the diagram to the preorder.  This argument explains why those sufficient criteria can hold without the universal comma condition \citep{Maltsiniotis2005Asphericite,Cisinski2009LocallyConstant}.

The sequential colimit itself is less problematic.  The nerve preserves filtered colimits, and for a sequence its homotopy colimit maps by a weak equivalence to its ordinary colimit.  We prove that two cofinal chains give isomorphic objects of $\Hot$ under the source's homotopy hypothesis: interleave cofinal subsequences, compare their bonding maps by homotopies, and compare the resulting telescopes.  This answers the existence part of the independence question.  The proof neither chooses a preferred isomorphism nor promotes it to a categorical homotopy equivalence in $\Cathot$.

Finite probes explain part of the remaining caution.  A finite category is finitely presented for ordinary categorical filtered colimits, which gives the asserted formula for $\Cathot$-morphisms into $C_\infty$.  Its nerve need not have the homotopy type of a finite CW complex.  The one-object category of the group of order two is a concrete counterexample to unrestricted compactness after passage to $\Hot$: its mod-two cohomology persists in every degree, whereas each finite-dimensional skeleton of its nerve has bounded cohomological dimension.  Replacing $\Cathot$ by $\Hot$ therefore needs a separate argument.  The mapping-set formulas at the end of manuscript page~117 are audited with their parameter $Y$ and their ambient ind-category restored in the commentary.

The proposed obstruction from strict loop multiplication on manuscript page~113 is already withdrawn in footnote~24.  Kan's loop-group theorem supplies strict simplicial group models \citep[Theorem~7.1]{Kan1958Combinatorial}.  For unrestricted paths, the specific comparison is taken up positively later in this same chapter, \S13; it is not a current open problem.  The constant-path lemma left unwritten at manuscript page~170 is supplied here by successively retracting the end vertices of an oriented interval.  Modern hammock localization gives a separate, precise mapping-space theory.  Its requirement that vertical arrows also be weak equivalences must be checked when comparing it with Grothendieck's full path fibres \citep{DwyerKan1980Calculating,ArakawaCnossen2026}.
\end{otherlanguage}
\dgrchunkend
